\subsection{Cycle classes with the manuscript's orientation}

We use the analytic comparison $c_X:H^r_{\dR}(X/\C)\to H^r(X^{\mathrm{an}},\C)$ and its inverse $\alpha_X$. Thus all scalar formulas below have the same Betti-to-de-Rham direction as the triangle. The comparison for arbitrary smooth complex algebraic varieties is Grothendieck's Theorem~$1'$ \cite{Grothendieck1966}; for smooth proper varieties and the trace convention we also use \cite[Theorem 1.4 and the cycle/trace paragraphs of \S1, pp.9--11 in the revised version]{Deligne1982Cycles}. We give the normalization argument explicitly.

\begin{result}{F.1.1}{The global cycle comparison, including singular cycles}{Proved}
Let $X/\C$ be smooth and proper, and let $z$ be an integral algebraic cycle of codimension $p$. Choose $i$ and orient each complex coordinate line by $(1,i)$. In ordinary, untwisted cohomology,
\[
 c_X\bigl(P_{\mathrm{alg}}(z)\bigr)=(2\pi i)^pP_{\mathrm{top}}(z),
 \qquad
 \alpha_X\bigl(P_{\mathrm{top}}(z)\bigr)=(2\pi i)^{-p}P_{\mathrm{alg}}(z).
 \tag{F.1.1}
\]
The equality holds with supports before passing to ordinary cohomology. The integral topological group maps to complex cohomology through its quotient by torsion.
\end{result}
\begin{proof}
For a smooth divisor with local equation $z_1=0$, its algebraic fundamental class is the residue class of $dz_1/z_1$. The positively oriented small normal circle has integral $2\pi i$. The topological Thom class has value one on the oriented normal disc. Applying the connecting maps of the punctured-disc pair therefore identifies the former class with $2\pi i$ times the latter. For a smooth codimension-$p$ inclusion, take transverse holomorphic coordinates $z_1,\ldots,z_p$. The product residue and the product Thom class give the factor $(2\pi i)^p$. The product ordering is the complex orientation ordering, so no permutation of real one-dimensional factors is inserted.

For an irreducible, possibly singular support $Z\subset X$ of codimension $p$, put $S=Z\setminus Z_{\mathrm{reg}}$. Each component of $S$ has complex codimension at least $p+1$ in $X$. Cohomology with supports in an oriented real $2d$-manifold identifies with Borel--Moore homology by cap product with its orientation class. A complex algebraic set of dimension at most $d-p-1$ has a triangulation of real dimension at most $2d-2p-2$; its Borel--Moore homology vanishes in larger degrees. Hence
\[
 H_S^{2p}(X,\C)=H_S^{2p+1}(X,\C)=0.
\]
The support localization sequence gives an isomorphism
\[
 H_Z^{2p}(X,\C)\longrightarrow
 H_{Z_{\mathrm{reg}}}^{2p}(X\setminus S,\C).
\]
The de Rham comparison with supports follows from the ordinary comparison for $X$ and $X\setminus Z$ by their localization triangles. Consequently the same restriction is an isomorphism in de Rham cohomology with supports. The locally defined algebraic fundamental class and the oriented topological fundamental class are therefore the unique extensions of their restrictions on the smooth locus. The local equality already proved extends uniquely across $S$. Additivity handles arbitrary cycles and their multiplicities. Forgetting supports proves (F.1.1). This also verifies the formerly uncertain support-to-ordinary arrow on 18:21: it is the natural map of localization triangles and commutes with comparison.

For $X=\mathbf P^1$, choose $u\in H^2(X,\Z)$ with integral one and let $v=P_{\mathrm{alg}}(\mathrm{point})$ in $H^1(X,\Omega_X^1)$. The case $p=1$ gives $\alpha_X(u)=(2\pi i)^{-1}v$, exactly 18:29. Changing $i$ to $-i$ reverses $u$ and the sign of $2\pi i$ together; the embedded rank-one lattice is unchanged. Deligne's definitions (2.1.13)--(2.1.14) \cite{Deligne1971HodgeII} use this same joint orientation convention.
\end{proof}

\begin{result}{F.1.2}{Proper direct image and every cohomological degree}{Proved}
Let $f:X\to Y$ be a proper morphism of smooth complex algebraic varieties, of pure dimensions $d,e$, and put $p=e-d$. The cohomological Gysin maps have types
\[
 f_*^{\B}:H^r(X,\C)\to H^{r+2p}(Y,\C),\qquad
 f_*^{\dR}:H^r_{\dR}(X)\to H^{r+2p}_{\dR}(Y).
\]
With the residue-normalized de Rham trace,
\[
 \alpha_Y f_*^{\B}=(2\pi i)^{-p}f_*^{\dR}\alpha_X.
 \tag{F.1.2}
\]
This includes negative $p$, and is compatible with composition. It proves the degree shift and scalar of 18:23 for $r=2i$ as well as for odd $r$.
\end{result}
\begin{proof}
First suppose $X,Y$ proper. Normalize algebraic trace by $\operatorname{Tr}_{\dR}P_{\mathrm{alg}}(\mathrm{point})=1$. Result F.1.1 gives
\[
 \int_X c_X(a)=(2\pi i)^d\operatorname{Tr}_{X,\dR}(a)
 \quad(a\in H^{2d}_{\dR}(X)).
\]
For $x\in H^r_{\dR}(X)$ and $y\in H^{2d-r}_{\dR}(Y)$, the definition of Gysin by duality and the projection formula give
\begin{align*}
 \int_Y c_Y(f_*^{\dR}x)\,c_Y(y)
  &=(2\pi i)^e\operatorname{Tr}_{Y,\dR}(f_*^{\dR}x\,y)\\
  &=(2\pi i)^e\operatorname{Tr}_{X,\dR}(x\,f^*y)\\
  &=(2\pi i)^p\int_X c_X(x)\,f^*c_Y(y).
\end{align*}
The degrees match because $2e-(r+2p)=2d-r$. Nondegeneracy of Poincare duality proves $c_Yf_*^{\dR}=(2\pi i)^pf_*^{\B}c_X$, which is (F.1.2). For smooth varieties that are not proper, the same proof pairs ordinary cohomology with compactly supported complementary cohomology. Properness of $f$ makes the required pullback on compact supports well-defined; the trace normalization is local and is the same residue normalization. The analytic-holomorphic version uses the holomorphic de Rham resolution and that trace, with the identical computation. Relative codimensions add for a composite, so the factors multiply to the factor for the composite; functoriality of the dual trace maps proves the full assertion.
\end{proof}

\begin{result}{F.1.3}{Correspondences, twists and the cancellation on 18:25}{Proved}
Let $X,Y/\C$ be smooth proper, $d=\dim X$, and $z\in\operatorname{CH}^r(X\times Y)$. Put $t=r-d$. The two untwisted actions are
\[
 z_*^{\B}(x)=p_{2*}^{\B}(p_1^*x\,P_{\mathrm{top}}(z)),\qquad
 z_*^{\dR}(x)=p_{2*}^{\dR}(p_1^*x\,P_{\mathrm{alg}}(z)),
\]
from degree $m$ to degree $m+2t$. They satisfy
\[
 \alpha_Yz_*^{\B}=(2\pi i)^{d-r}z_*^{\dR}\alpha_X.
 \tag{F.1.3}
\]
Consequently comparison is compatible with the degree-zero correspondence action ($r=d$), and with the general action after the target is written $H^{m+2t}(Y)(t)$.
\end{result}
\begin{proof}
The projection $p_2:X\times Y\to Y$ has relative codimension $-d$. Apply F.1.2 to it and F.1.1 to $z$, using compatibility of comparison with pullback and cup product. The two factors are $(2\pi i)^d$ and $(2\pi i)^{-r}$, whose product gives (F.1.3). The Tate line $\Q(t)_{\B}=(2\pi i)^t\Q$ contributes $(2\pi i)^t$ to the Betti-to-de-Rham comparison, cancelling $(2\pi i)^{-t}$. Thus this is an equality of typed morphisms of realizations of pure weight $m$.

For correspondences $z:X\rightsquigarrow Y$ and $w:Y\rightsquigarrow Z$, their composite class is $p_{13*}(p_{12}^*z\,p_{23}^*w)$. Projection formula and the base-change square on the triple product identify its action with $w_*z_*$. These are the ordinary Gysin/cup operations just checked; the cycle degrees are even, and the twists add. The diagonal has codimension $d$ and acts as the identity. This proves functoriality on the correspondence category, and on summands and twists by applying the corresponding projectors. The doubtful exponent on 18:25 must be the dimension of the projected-out factor for an arbitrary correspondence; with the page's degree-zero notation $z\in A^n(X^n\times Y)$ it is $n$ and cancels the cycle factor. We retain both transcription alternatives and do not use the final displayed $\Omega_X$ on a cohomology group of $Y$: the typed target is $\Omega_Y$.
\end{proof}

\begin{result}{F.1.4}{An unchosen complex isomorphism and its exact lattice dependence}{Proved; unrestricted minus-sign reading refuted}
Let $k$ be an abstract field isomorphic to $\C$, let $X/k$ be smooth, and choose an isomorphism $\sigma:\C\to k$, in the direction written on 18:19. Set $X_\sigma=X\otimes_{k,\sigma^{-1}}\C$. The transported Betti cycle class in $H^{2p}_{\dR}(X/k)$ is
\[
 P^{(p)}_\sigma=\sigma(2\pi i)^{-p}P_{\mathrm{alg}}.
 \tag{F.1.4}
\]
Changing $\sigma$ to $\sigma c$, with $c$ complex conjugation, multiplies this class by $(-1)^p$ and preserves the whole embedded integral Betti lattice (modulo torsion). Different complete archimedean topologies can give different lattices, already for $\mathbf P^1$.
\end{result}
\begin{proof}
Extend the ordinary comparison by scalar transport along $\sigma$. It is $\sigma$-semilinear on the complex Betti coefficients and $k$-linear after $H_{\B}\otimes_{\Z}k$ is formed. Applying this map to F.1.1 gives (F.1.4). Replacing $\sigma$ by $\sigma c$ changes $\sigma(2\pi i)$ to its negative. The antiholomorphic map $X_\sigma\to X_{\sigma c}$ gives an integral cohomology isomorphism. Combined with coefficient conjugation, its comparison maps coincide in $H_{\dR}(X/k)$; hence the image lattices coincide in every degree. The cycle generator acquires the stated parity sign because its normal bundle has complex rank $p$.

The unrestricted single-minus reading is false: on $\mathbf P^2$, with $h$ the hyperplane generator, conjugation sends $h$ to $-h$ and fixes the nonzero class $h^2$ of a point. The final minus sign and the missing exponent in the lattice formula of 18:19 are therefore valid in the immediately preceding $p=1$ example, not as formulas for arbitrary $p$.

For the topology dependence, identify $k=\C$ as an abstract field. Choose a transcendence basis containing $\pi$, keep $i$ and the other basis elements fixed, send $\pi$ to $2\pi$, and extend the resulting automorphism of $\Qb$ with that basis to its algebraic closure $\C$. Such an extension exists by the extension theorem for embeddings of algebraic closures. For $X=\mathbf P^1$, the two embedded lattices in $kv$ are
\[
 (2\pi i)^{-1}\Z v\quad\hbox{and}\quad(4\pi i)^{-1}\Z v.
\]
They have index two and are unequal. This proves dependence of the lattice, with its exact map, independently of any unknown period relation.
\end{proof}

\begin{result}{F.1.5}{The two readings of ``primitifs'' on 18:29}{Proved exact image; an unrestricted reading refuted}
In the one-dimensional line $H^1(\mathbf P^1_K,\Omega^1)=Kv$, as complete archimedean topologies and the accompanying orientation choices vary, the comparison images of an oriented integral generator are exactly
\[
 \{a v:a\in K\text{ is transcendental over }\Q\}.
 \tag{F.1.5}
\]
For a fixed one of these lattices the primitive elements are its two generators. Without a lattice, ``primitive'' has no specified meaning. The reading ``all possible'' is correct for the displayed set of comparison generators; the reading ``every nonzero vector of $Kv$'' is false.
\end{result}
\begin{proof}
The coefficient is $a=\sigma(2\pi i)^{-1}$ and is transcendental because $2\pi i$ is transcendental. Conversely, for a given transcendental $a\in K$, extend $2\pi i$ and $a^{-1}$ to transcendence bases. Their equal cardinalities follow from $K\simeq\C$. A bijection of the bases sending the former element to the latter, followed by an extension to algebraic closures, gives an isomorphism $\sigma:\C\to K$ with $\sigma(2\pi i)^{-1}=a$. Pull back the usual complex topology and its orientation along $\sigma$; it is a complete archimedean field topology realizing $av$. An embedded cyclic lattice $\Z av$ has primitive elements $\pm av$. In particular $v$ itself, with algebraic coefficient one, never occurs. This decides the mathematical quantifiers of both frozen readings C07 without deciding the handwritten word.
\end{proof}

\begin{result}{F.1.6}{Real descent of the diagonal integral lattice}{Proved}
For a smooth proper variety $X/\R$, set $L=H^r(X(\C),\Z)/\mathrm{torsion}$. Integral Betti--etale comparison and Betti--de-Rham comparison embed $L$ diagonally in
\[
 \left(\prod_\ell H^r_{\mathrm{et}}(X_\C,\Z_\ell)/\mathrm{torsion}\right)
 \times \bigl(H^r_{\dR}(X/\R)\otimes_\R\C\bigr).
\]
The image is stable under $\Gal(\C/\R)$, acting on the first factor by geometric Galois action and on the second by coefficient conjugation. This is the integral assertion C08 on 18:29.
\end{result}
\begin{proof}
Let $f_\infty:L\to L$ be the integral involution induced by conjugating complex points. Each etale comparison $c_\ell$ satisfies $c_\ell f_\infty=c\,c_\ell$, by functoriality of finite-coefficient comparison and inverse limit; on the torsion-free quotients the resulting map is $L\otimes\Z_\ell\simeq H^r_{\mathrm{et}}/\mathrm{torsion}$. On de Rham cohomology, the exact semilinear identity is
\[
 (1\otimes c)\alpha_X=\alpha_X(f_\infty\otimes c).
\]
For $x\in L$, all coordinates of the conjugated diagonal image are therefore the coordinates of $f_\infty x\in L$. This proves stability, including the sign on odd Tate powers, without assuming the involution is trivial.
\end{proof}

\subsection{The mixed-Hodge index on 18:9}

\begin{result}{F.2.1}{The corrected weight-graded twist}{Proved; unshifted formula refuted}
For an integral or rational mixed Hodge structure $E$, using increasing $W$ and decreasing $F$,
\[
 W_i(E(n))=W_{i+2n}(E)\otimes\Q(n),\qquad
 F^a(E(n))=F^{a+n}(E)\otimes\C,
\]
with the evident integral tensor line in the integral case. Consequently
\[
 \operatorname{Gr}_i^W(E(n))=
 \bigl(\operatorname{Gr}_{i+2n}^W E\bigr)(n).
 \tag{F.2.1}
\]
In particular the index in $\operatorname{Gr}_i^W(E(-1))$ is $i-2$ on the right.
\end{result}
\begin{proof}
The Tate object has weight $-2n$ and bidegree $(-n,-n)$ by Deligne's Definition (2.1.13). Tensoring shifts weights and Hodge degrees by precisely these amounts. Taking the quotient $W_i/W_{i-1}$ gives (F.2.1) because tensoring with a rank-one free module is exact. As a complete counterexample to the unshifted reading, take $E=\Q(0)$ and $i=0$: $\operatorname{Gr}_0^W(E(-1))=0$, whereas $(\operatorname{Gr}_0^W E)(-1)=\Q(-1)\ne0$. We leave the unshifted line in the transcription as evidence and correct its mathematics here; R5's filtration shifts themselves are retained.
\end{proof}

\subsection{F.3. Additional proofs required by the note inventory}

The assertion inventory retains every requested note entry, including MB03b and the eight cote-18 rows of Standing Question~I. Its rows record the original claim identifier, frozen page and line, original cancellation or uncertainty, present mathematical status, proof dependency, and exact residual obligation. The additional arguments below supply the dependencies that were not already proved in Parts~A--F.2. A mathematical correction does not decide a doubtful handwritten symbol.

\begin{result}{F.3.1}{Oda's comparison and the characteristic-$p$ Hodge extension}{Known; scope and extension identified}
Let $k$ be perfect of characteristic $p$, including $p=2$, and let $A/k$ be an abelian variety. In Oda's contravariant, left-Dieudonne-module convention there is a canonical isomorphism
\[
 M(A[p])\xrightarrow{\sim}H^1_{\dR}(A/k).
\]
It identifies $VM(A[p])$ with $H^0(A,\Omega_A^1)$ and the quotient with $H^1(A,\mathcal O_A)$, and hence identifies their exact sequences. This resolves the characteristic-$p$ extension question in I.P04 and II.RH18 in precisely this abelian scope.
\end{result}
\begin{proof}
These are the isomorphism and submodule identification of Oda's Corollary~5.11, printed p.131 \cite{AuditOda1969}. Its proof composes Proposition~4.1 and Theorem~5.10. The same page identifies the $V$-submodule using the Frobenius-twisted Dieudonne module of the Frobenius kernel. The operators $F,V$ on the domain are semilinear and satisfy $FV=VF=p=0$; transporting them through the isomorphism makes the statement a comparison of the additional structures, not only of vector-space dimensions. The resulting exact sequence is
\[
0\longrightarrow H^0(A,\Omega_A^1)\longrightarrow H^1_{\dR}(A/k)
\longrightarrow H^1(A,\mathcal O_A)\longrightarrow0.
\]
Thus the extension with its distinguished maps is recovered from $A[p]$. No canonical splitting of this sequence is asserted. Mixed-characteristic lifts and their Galois realizations require the additional maps and admissibility conditions proved in Part~B. We inspected the printed source page; this is not a claim to have inspected an archive manuscript scan.
\end{proof}

\begin{result}{F.3.2}{The rational Lefschetz operator images and an integral counterexample}{Proved}
On the primitive string $V_m=\bigoplus_{j=0}^m\Q v_j$, put
\[
 Ev_j=v_{j+1},\quad \Lambda v_j=v_{j-1},\quad
 Fv_j=j(m-j+1)v_{j-1},\quad Hv_j=(2j-m)v_j,
\]
with zero at the endpoints. Then
\[
 [E,F]=H,\qquad[H,E]=2E,\qquad[H,F]=-2F.
\]
On $W=\bigoplus_{m=0}^dV_m$ the actual operator subalgebras satisfy
\[
 \Q[E,\Lambda]=\Q[E,F,H]
 =\prod_{m=0}^d\operatorname{Mat}_{m+1}(\Q).
\]
On the displayed integral string of length two, however,
\[
 \Z[E,F,H]\ne\Z[E,\Lambda]=\operatorname{Mat}_3(\Z).
\]
These prove the exact rational assertions of II.MB03--MB03b and refute their unqualified integral extension for this lattice.
\end{result}
\begin{proof}
Apply each commutator to $v_j$ and subtract; the endpoint formulas give the same identities. If $\pi_j$ projects to $\Q v_j$, the conversion in $\End_\Q(V_m)$ is
\[
 F=\sum_{j=1}^m j(m-j+1)\Lambda\pi_j,
 \qquad
 \Lambda=\sum_{j=1}^m\frac{F\pi_j}{j(m-j+1)}.
\]
Put $B=1-E\Lambda$ and $T=1-\Lambda E$. These project to the bottom and top of each string. Therefore
\[
 B_m=B\Lambda^mTE^mB
\]
projects onto the bottom of the length-$m$ string alone: $E^m$ kills shorter strings, while $T$ kills the displayed image in longer ones. For $0\leq a,b\leq m$, the operator $E^bB_m\Lambda^a$ sends $v_a$ in that string to $v_b$ and kills every other basis vector. These are its matrix units, proving the assertion for $\Q[E,\Lambda]$.

The Casimir $C=H^2+2H+4FE$ acts on $V_m$ as $m(m+2)$. These eigenvalues are distinct for $m\geq0$. Polynomial interpolation in $C$ projects onto each string, and interpolation in $H$ there gives every $\pi_j$. The conversion formula consequently recovers $\Lambda$ from $E,F,H$. Conversely the matrix-unit result places $F,H$ in $\Q[E,\Lambda]$. This proves equality as subalgebras of $\End(W)$. For multiplicity spaces $P_m$, the same formulas act on $\bigoplus P_m\otimes V_m$ by identity on each $P_m$.

For $m=2$, both $F$ and $H$ vanish modulo~2. Thus reduction of $\Z[E,F,H]$ is contained in $\mathbf F_2[E]$. Every such operator sends $v_2$ into $\mathbf F_2v_2$, whereas $\Lambda v_2=v_1$. Hence $\Lambda\notin\Z[E,F,H]$. The integral matrix-unit formula proves $\Z[E,\Lambda]=\operatorname{Mat}_3(\Z)$ on this string. No division-free identification has been inferred from a rational formula.

The attached exact certificate checks commutators, conversions and Casimir for $m=0,\ldots,10$, all matrix units for $d=0,\ldots,5$, and the modulo-2 obstruction. The basis-vector proof above covers all lengths; the finite checks are not substituted for it.
\end{proof}

\begin{result}{F.3.3}{Weight, Cartan and filtration gradings with their comparison maps}{Proved}
For a smooth projective $d$-fold, on $H^i(X)$ put $w_X(t)=t^i$ and $\lambda_X(t)=t^{i-d}$. Then $\lambda_X(t)=t^{-d}w_X(t)$, and both conjugate a degree-$j$ operator by $t^j$. For the character $\varepsilon$ on $\Q(-1)$, the tensor weight gives $\varepsilon w=t^2$. A filtration splitting with $\varepsilon j=t$ is a separately specified map.

The displayed Lefschetz representation on $H^*(\mathbf P^1)$ does not factor through the Tate representation $\rho(z)=\operatorname{diag}(1,z)$, but satisfies
\[
 \rho(z)E\rho(z)^{-1}=zE,\qquad
 \rho(z)F\rho(z)^{-1}=z^{-1}F.
\]
The map $E:h^0\to h^2(1)$ is equivariant. On elliptic $H^1$ the Lefschetz Lie-algebra action is zero, so it is not the nonzero standard action of the normalized period $\SL_2$ there.
\end{result}
\begin{proof}
The weight formulas follow on each graded summand; tensor compatibility gives $w_{\Q(1)}(t)=t^{-2}$. The central tensor grading is the one in \cite[Definition 5.1]{DM1982}. For an explicit filtration splitting, take $G=\GL_2$, upper triangular $B$, diagonal Levi $L$, $\varepsilon=\det$ and $j(t)=\operatorname{diag}(t,1)$. It is central in $L$, has centralizer $L$, and $\varepsilon j=t$. But $j(t)E_{12}j(t)^{-1}=tE_{12}$, so it is not central in $G$. The central cohomological weight is $w(t)=tI$, with determinant $t^2$. This calculates exactly the normalization issue in I.P14, I.M03 and II.MB05.

For $H^*(\mathbf P^1)=\Q e_0\oplus\Q e_2$,
\[
 E=\begin{pmatrix}0&0\\1&0\end{pmatrix},\quad
 F=\begin{pmatrix}0&1\\0&0\end{pmatrix},\quad
 H=\operatorname{diag}(-1,1).
\]
Every $\rho(z)$ fixes $e_0$, while $\exp(sE)e_0=e_0+se_2$, proving the failed factorization. Direct multiplication gives the two conjugation formulas. Twisting $h^2$ by $(1)$ cancels its Tate character, proving equivariance of the point-class morphism. On elliptic $H^1$, raising and lowering have targets $H^3$ and $H^{-1}$, both zero, and $H=i-d=0$. This proves the exact representation comparison used in I.M02 and II.MB06. It makes no assertion that the surrounding mathematical theories lack other comparison maps.
\end{proof}

\begin{result}{F.3.4}{Reconstruction and the specified modern motivic comparisons}{Known; maps and hypotheses checked}
For a neutral rigid abelian $\Q$-linear tensor category $\mathcal C$ with exact faithful fibre functor $\omega$, the tensor-automorphism construction gives an equivalence
\[
 \mathcal C\xrightarrow{\sim}\Rep_\Q\bigl(\Aut^\otimes(\omega)\bigr),
\]
and tensor-isomorphisms between fibre functors form the corresponding torsor. With compatible Betti fibre functors and a fixed $\Qb\subset\C$, rational Nori motives have the exact affine-group-scheme sequence
\[
 1\longrightarrow G_{\mathrm N}(\Qb)
 \longrightarrow G_{\mathrm N}(\Q)
 \longrightarrow\underline{\Gal(\Qb/\Q)}\longrightarrow1.
\]
For a field embedded in $\C$ and a principal ideal coefficient domain, the specified Ayoub and Nori Betti motivic Hopf algebras are isomorphic by the constructed inverse maps. The pure rational Nori category is equivalent to Andre's motivated category. General Nori Hodge fullness remains the explicitly stated conjectural obligation; fullness on the abelian-generated pure category over $\C$ is established.
\end{result}
\begin{proof}
For a morphism $a:M\to N$, the tensor automorphism acts by the exact naturality equation
\[
 g_N\omega(a)=\omega(a)g_M.
\]
The equivalence and the Isom torsor are \cite[Theorems 2.11 and 3.2]{DM1982}. On cycle correspondences the morphism is $p_{2*}(p_1^*(-)\cup\operatorname{cl}(z))$, with the degree and twist proved in F.1.3. The triple-product push--pull formula proves composition. Adding projectors and Tate twists gives the named pure categories; this is not a claim that an unspecified historical category was constructed on the supplied pages.

Base change of Nori motives induces the first arithmetic-group arrow, while inclusion of Artin motives induces the last arrow to the proconstant group. Exactness follows from \cite[Theorem 5.10, \S5.8, Theorem 10.2 and Remark 10.3]{AuditKahn2025}. Theorem 5.10 no longer assumes semisimplicity. Kahn's p.589, footnote~1, identifies the normality gap in the older proof of Deligne--Milne Proposition 6.23(a); his argument supplies the needed exactness. Exactness does not specify a splitting or identify its kernel with the identity component.

The Hopf maps $\varphi_a,\varphi_n$ of \cite[Theorem 9.1, p.41]{AuditChoudhuryGallauer2017} are inverse. Applying $\Spec$ reverses them and gives the group-scheme isomorphism with its specified Betti realization. An integral representation and its chosen lattice are transported; a rational representation alone has no lattice selected by this operation.

The pure semisimplicity and equivalence with motivated motives are \cite[Theorems 6.3.3 and 6.4.1, pp.57 and 60--61]{AuditArapura2013}. The functor sends algebraic-cycle correspondences to motivated ones and commutes with Betti realization. The general Hodge-fullness assertions are Conjectures 7.1.1--7.1.2, p.62, while Corollary 7.1.4, p.63, treats the abelian-generated category. Over $\R$ the additional descent compatibility is retained as in C.1.1. Neither the Hopf isomorphism nor pure semisimplicity proves the unspecified six global essential-image claims or the general algebraic-cycle surjectivity required in C.1.2.
\end{proof}

\begin{result}{F.3.5}{Numerical quotient, centres and central Tate extension data}{Known and proved as specified}
The identity on algebraic cycles induces the surjection
\[
 \operatorname{Hom}_{\mathrm{hom}}(M,N)
 \twoheadrightarrow\operatorname{Hom}_{\mathrm{num}}(M,N).
\]
Rational numerical motives are semisimple. A cohomology realization factors through this quotient exactly when its kernel acts trivially; on the relevant smooth projective products this is numerical equals homological equivalence for the correspondence classes.

Semisimplicity does not imply a trivial centre. Nor does a normalized group $\SL_2$ with a surjection to the Tate torus determine the full group: $(\GL_2,\det)$ and $(\Gm\times\SL_2,\operatorname{pr}_1)$ have that kernel and are not isomorphic. Their explicit comparison is the central isogeny $(q,A)\mapsto qA$, which squares the specified Tate character.
\end{result}
\begin{proof}
A homologically zero cycle has zero pairing with every complementary cycle, so is numerically zero. This proves the quotient map. Semisimplicity is \cite[Theorem 1, p.448]{AuditJannsen1992}. Factoring a functor through a quotient category is exactly annihilation of its kernel. Kunneth and Poincare duality identify a correspondence's cohomology class with its action on all cohomological degrees; hence the annihilation criterion is the stated equality of equivalence relations. The general equality is retained as an unresolved standard-conjecture obligation, rather than inferred from semisimplicity.

Every rational representation of $\Gm$ decomposes into integer-character spaces, proving semisimplicity of its representation category. Its entire group is the kernel of conjugation, giving the centre counterexample. For the second assertion, the centres of $\GL_2$ and $\Gm\times\SL_2$ are respectively connected $\Gm$ and disconnected $\Gm\times\mu_2$. Both displayed character kernels are $\SL_2$. The map $(q,A)\mapsto qA$ has kernel $\{(1,I),(-1,-I)\}$ and determinant $q^2$. These computations exhibit both the nonisomorphism and the surviving relation, so the missing central data are explicit.

For the historical A--D dictionary we checked \cite[pp.195--197]{AuditGrothendieck1969}: printed A is Lefschetz on algebraic classes, B is algebraicity of the lowering correspondence, and C is algebraicity of Kunneth projectors. The operational manuscript A and C therefore match printed A and B. Printed D is numerical equals homological equivalence in \cite[p.17, Proposition 5-1]{AuditKleiman1994}. The proposed identification of the remaining two historical letters lacks standalone definitions in the supplied witnesses and stays a documentary obligation.

The newly acquired originals narrow that obligation without
erasing it. The letter to Serre of 27 August 1965, cote~14,
pp.43--44, explicitly defines its cycle groups modulo
algebraic equivalence and its letter B by positive
definiteness of the primitive middle-dimensional intersection
form \cite{AuditGrothendieckSerre1965}. The comparison of
those cycle groups with the later cohomological convention
is the actual surjection
\[
 \bigl(\operatorname{CH}^r(X)/{\sim_{\rm alg}}\bigr)\otimes\Q
 \longrightarrow \operatorname{im}(\operatorname{cl}_H),
 \qquad [z]\longmapsto\operatorname{cl}_H(z).
\]
It is well defined because an algebraic family has constant
cycle class on a connected parameter curve; it is surjective
by the definition of its target and commutes with products
and multiplication by the hyperplane class. This records
the comparison without identifying its kernel with zero.
The Coates letters of 4 January 1967 and 6 January 1966
give, respectively, the stated relations among A--D and
the inverse-Lefschetz characterization of C, but do not
write standalone definitions of their B and D
\cite[pp.1--3]{AuditGrothendieckCoates1967}
\cite[p.1]{AuditGrothendieckCoates1966}.
The 1965 definition therefore supplies a dated, explicit
meaning of B; the particular B and D of the 1967 letter
remain to be located. The finite-source check does not
claim that no such definition exists elsewhere.
\end{proof}

\begin{result}{F.3.6}{The actual maps relating the anabelian and linear statements}{Known and proved as specified}
Mochizuki's sub-$p$-adic Hom theorem concerns dominant maps and open continuous arithmetic fundamental-group homomorphisms, modulo geometric inner automorphisms. For proper curves, the induced linear comparison is geometric pro-$\ell$ abelianization followed by the Albanese identification with the Tate module of the Jacobian. Tate's finite-field and Faltings's number-field Hom theorems govern the resulting abelian-variety maps.

For a centre-free geometric group, the arithmetic extension is recovered as the pullback of $\Aut\to\operatorname{Out}$ along its outer Galois action. For proper real hyperbolic curves, section classes recover connected real components, but their conjugation actions on geometric abelianization coincide. Stix's cusp-inertia criterion requires the non-Tate-weight condition, an open arithmetic normalizer image, and the specified cyclotomic character action.
\end{result}
\begin{proof}
The precise Hom theorem is \cite[Theorem A, PDF p.3, equivalently Theorem 16.5 and its following remark]{AuditMochizuki1999}. It includes the profinite and geometrically pro-$p$ versions with their specified hypotheses. For proper curves with chosen base points, the Albanese map induces
\[
 \pi_1(X_{\bar k})^{\mathrm{ab},\ell}\simeq T_\ell J_X.
\]
A curve morphism gives its covariant Jacobian map on these modules; the geometric inner ambiguity disappears after abelianization. Over a finite field the integral Hom identification is \cite[Main Theorem, p.134]{AuditTate1966}; over a number field we use the rational Faltings Hom theorem as already checked in D.3.2. The resulting linear Hom space is not universally filled by curve maps: on a genus-$g\geq2$ curve, multiplication by two on the Jacobian gives an invertible rational Tate-module map. Riemann--Hurwitz forces a nonconstant self-map of the curve to have degree one, so it preserves the principal polarization, whereas multiplication by two pulls that polarization back to four times itself.

For $1\to\Delta\to\Pi\to\Gamma\to1$ with $Z(\Delta)=1$, define
\[
 \Pi\longrightarrow
 \Gamma\mathop{\times}_{\operatorname{Out}(\Delta)}\Aut(\Delta),
 \qquad p\longmapsto(\bar p,\operatorname{Ad}(p)|_\Delta).
\]
Its kernel is $Z(\Delta)$. For a pair $(\gamma,a)$ on the right, lift $\gamma$ to $p$; equality of outer classes gives $a=\operatorname{Ad}(\delta)\operatorname{Ad}(p)$ for some $\delta\in\Delta$, so $\delta p$ maps to the pair. This proves the isomorphism, with the topological version obtained by the continuous maps between the profinite groups in question. The curve-group hypotheses are \cite[Lemma 1.3.1 and Corollary 1.3.2]{AuditMochizuki2004}; the author's correction notes for both papers were checked and do not alter these locators.

For a smooth proper real curve of genus at least two, \cite[Theorem 1.3 and Example 5.8]{AuditPal2011} applies with all primes and gives the component-to-section-class bijection. Lifts of conjugation for two sections have the form $l$ and $\delta l$. Their actions on $\Delta^{\mathrm{ab}}$ coincide because inner conjugation by $\delta$ is trivial there. The full abelianized extension or the section cocycle can contain additional data; this calculation concerns its conjugation action.

The non-Tate condition in \cite[Definition 26]{AuditStix2012} is $\operatorname{Hom}_{G_k}(\Z_\ell(1),T_\ell A)=0$ for all abelian varieties and primes. For number fields, choose a good prime away from $\ell$; arithmetic Frobenius has eigenvalue $q$ on the twist and eigenvalues of absolute value $q^{1/2}$ on $T_\ell A$, proving the vanishing. Definition~29 includes both the open arithmetic normalizer image and its cyclotomic conjugation action. Theorem~30 then identifies the nontrivial closed pro-$\ell$ subgroups of the geometric fundamental group lying in cusp inertia. The full LM6 source has now been obtained and the exact A.50 enlargement is checked in F.3.10. Its residual weight-one character still does not supply the cyclotomic-action premise needed to close A.22. F.3.15 independently re-proves the pro-prime recognition result; F.3.16--F.3.17 give actual additional cover maps and their precise remaining visibility limit.
\end{proof}

\begin{result}{F.3.7}{Real fixed points and the corrected compact-centralizer assertion}{Refuted and proved as specified}
The equality of $H^*(X(\R),\Q)$ with the conjugation invariants in $H^*(X(\C),\Q)$ is false, including the unrestricted degree-zero reading. The raw Weil-element centralizer in $G(\R)$ need not be compact. In the polarizable real category of C.1.1, after quotienting a finite tensor-generator group by the central weight image, the centralizer of the Weil element is maximal compact. The two-element normalizer written in the cancelled paragraph then equals that centralizer.
\end{result}
\begin{proof}
For $X=\mathbf P^1_\R$, the complex $H^1$ is zero, while $X(\R)=\mathbf{RP}^1$ is a circle with $H^1=\Q$. Complex conjugation on the sphere is reflection across the real equator. Choosing the representative in one closed hemisphere identifies the quotient with a closed disk. Its rational cohomology equals the invariant sphere cohomology: reflection acts by $-1$ on $H^2$. This computes the quotient example in II.RH08 as well as the failed fixed-point equality. For $X=\Spec\C$ over $\R$, the two complex points are interchanged, their $H^0$ invariants equal $\Q$, and the real locus is empty with $H^0=0$.

The pure Tate group $\Gm$ has Weil element $C=1$ and centralizer $\R^\times$, providing the raw compactness counterexample. Now put $\dot G=G/\im(w)$ and let $C$ be the image of $h(i)$. Then $C^2=w(-1)=1$ in $\dot G$, so the normalizer of $\{C,C^{-1}\}$ equals $Z_{\dot G}(C)$. In particular $f_\infty C f_\infty^{-1}=C^{-1}$ places $f_\infty$ in it.

Choose a faithful weight-zero representation of $\dot G$ and a polarization $q$. Such a representation exists because representations of the quotient are exactly those on which the central weight image acts trivially; in the stated category they are polarizable. A finite direct sum of their polarizations gives an invariant real symmetric $q$ such that $B(x,y)=q(x,Cy)$ is positive definite. If $*$ denotes its adjoint, invariance of $q$ and $C^2=1$ give
\[
 \operatorname{Ad}(C)(g)=(g^*)^{-1}.
\]
Hence $K=Z_{\dot G(\R)}(C)=\dot G(\R)\cap O(B)$ is compact. If a compact subgroup $H$ contains $K$, then $C\in H$, so $H$ is stable under $\operatorname{Ad}(C)$. For $g\in H$ we obtain $g^*\in H$ and $g^*g\in H$. This positive self-adjoint matrix has all positive and negative powers in a compact set. Each eigenvalue and its reciprocal therefore have bounded powers, forcing all eigenvalues to be one. Thus $g^*g=1$ and $g\in K$, proving maximality. The proof permits nonconnected finite-dimensional groups. For a pro-group we make this assertion only on each finite tensor-generator quotient.

The compact-band theorem \cite[Theorem 4.17(a)]{DM1982} is related by this specified positive form and real structure. The entire paragraph on 18:38 remains cancelled as documentary evidence. Its live real-splitting assertion is proved separately in C.4.2: $\beta$ is defined over $\R$ and $\R(\gamma)=\R(\alpha)$. We do not repeat that result or infer algebraicity over $\Qb$ from it.
\end{proof}

\begin{result}{F.3.8}{Integral Chow specialization and its integral cycle-class square}{Known; construction checked}
For $X$ smooth over a discrete valuation ring $R$, there is an integral specialization ring homomorphism
\[
 \operatorname{sp}:\operatorname{CH}^*(X_\eta)
 \longrightarrow\operatorname{CH}^*(X_s)
\]
which sends a prime cycle to the special fibre of its flat closure, with its scheme-theoretic multiplicities. For $R$ complete, $X$ smooth proper, and $\ell$ invertible in $R$, the map on geometric fibres commutes with the integral $\Z_\ell$ cycle-class map under the proper smooth specialization isomorphism. This retains torsion and answers the stated DVR version of I.P07--P08.
\end{result}
\begin{proof}
Let $i:X_s\to X$ and $j:X_\eta\to X$. Closure of cycles gives the surjection $j^*:\operatorname{CH}^r(X)\to\operatorname{CH}^r(X_\eta)$. Its kernel is the image of $i_*$: a rational equivalence on the generic fibre extends by closing each supporting variety and taking the same rational function; its discrepancy is supported on $X_s$. Cartier Gysin factors through rational equivalence, and
\[
 i^!i_*a=c_1(N_{X_s/X})a=0
\]
because the normal line is trivial, cut out by a uniformizer. Thus $\operatorname{sp}(z)=i^!\widetilde z$ is independent of the lift $\widetilde z$ and uses integral coefficients. The precise Gysin facts are \cite[Lemmas 42.30.2--42.30.3]{AuditStacksChow}. The closure is flat over the DVR, and the Gysin class is its special fibre with multiplicities.

Compatibility with proper pushforward, flat pullback and refined Gysin is \cite[Proposition 20.3, p.399]{AuditFulton1998}. Applying this to the diagonal regular immersion of the smooth scheme proves preservation of intersection products, as in Corollary~20.3 on the same page. For geometric fibres over a complete DVR, pass to finite extensions of its fraction field and their integral closures, and then to the direct limit. Example~20.3.5, p.400, gives exactly the commutative square with integral $\Z_\ell(r)$ coefficients. Its construction restricts the relative class of a flat cycle closure to the generic and special fibres. Proper smooth base change identifies the latter cohomology groups; rational-equivalence descent then gives the square on Chow groups. No tensoring with $\Q$ is made.

The source status is precise. Fulton's printed \S20.3, Proposition and Corollary~20.3, and Examples 20.3.1 and 20.3.5 were opened in a text reproduction of the primary book. The university-hosted PDF located separately contains only front matter and was not used for these statements. The original SGA~6 scan mirror \cite{AuditSGA6} was downloaded; PDF pages 586--587 show printed pp.579 and 581, including the beginning of Remark~7.13.11 and its continuation, but the mirror omits printed p.580 containing (7.13.12). We do not claim that missing original page was inspected. The supplied annotated 18:7 transcription remains the archive witness, while the modern result has the checked Fulton and Stacks proof above.
\end{proof}

\begin{result}{F.3.9}{The transversal dimension on 18:15}{Proved correction; reading retained}
At an isolated transverse intersection with a codimension-$p$ subvariety of a smooth complex $n$-fold, a smooth transversal has dimension $p$. The frozen dimension label $n-p$ in I.P28 has that value only when $n=2p$.
\end{result}
\begin{proof}
At a smooth point of the two subvarieties, transversality and isolation give
\[
 T_yX'+T_yY=T_yX,\qquad T_yX'\cap T_yY=0.
\]
Taking dimensions yields $\dim X'+(n-p)=n$. This proves the precise correction under ordinary complex-dimensional grading. The written $n-p$ and the possibility of an unstated convention remain in the reading record; the dimension computation is not paleographic evidence.
\end{proof}

\begin{result}{F.3.10}{The complete LM6 character statement and its repair}{Proved repair; residual problem retained}
The complete LM6 A.50 proves descent of a \emph{power} of a
normalizer character. Its noncuspidal alternative produces a
weight-one Jacobian character; it does not supply the
cyclotomic-action premise needed to close LM6 A.22. The
previously unavailable full statement and proof have now been
read. Their first Sylow-subgroup step is valid after the
explicit enlargement $N\mapsto NA$ below.
\end{result}
\begin{proof}
Use the original notation: $U/K$ is an affine hyperbolic curve
over a number field, $1\to F\to E_U\xrightarrow{p}G_K\to1$
its arithmetic fundamental group, $A\ne1$ a closed procyclic
subgroup of $F$, and $N\subset E_U$ a closed subgroup
normalizing $A$ with $p(N)=G_L$ open. For $A_\ell\ne1$,
the action on $A_\ell\simeq\Z_\ell$ is
$\theta_\ell:N\to\Z_\ell^\times$. The conclusion is
\[
 \lambda_\ell:G_L\longrightarrow\Z_\ell^\times,
 \qquad \lambda_\ell p=\theta_\ell^d,
 \qquad
 d=\begin{cases}\ell-1,&\ell>2,\\1,&\ell=2.\end{cases}
 \tag{F.3.10a}
\]
The original proof immediately chooses a Sylow subgroup of
$N\cap F$ containing $A_\ell$, although the statement only
assumes that $N$ normalizes $A$. Define
$\widetilde N=NA$ and $\widetilde D=(N\cap F)A$.
Multiplication $N\times A\to E_U$ has compact image, so
$\widetilde N$ is closed; normalization makes it a subgroup.
The equality $p(na)=p(n)$ proves
$p(\widetilde N)=G_L$ and
$\widetilde N\cap F=\widetilde D$. Conjugation gives the
well-defined extension
$\widetilde\theta_\ell(na)=\theta_\ell(n)$: the subgroup
$A$ centralizes $A_\ell$, so the formula is independent of
the factorization. Restriction to $N$ recovers the original
character and every original hypothesis.

Now $A_\ell\subset\widetilde D$ is normal. A pro-$\ell$
Sylow subgroup $S$ containing it exists. Closed pro-$\ell$
subgroups of the free profinite group $F$ are free pro-$\ell$:
cohomological dimension is at most one on closed subgroups,
and the pro-$\ell$ freeness criterion applies; the criterion
is also recalled in \cite[\S1.5, Lemma 1.7 and the proof of
Proposition 1.6]{AuditGille2000}. For $g\in S$, the closed
subgroup generated by $A_\ell$ and $g$ is metabelian, since
it has abelian normal subgroup $A_\ell$ and procyclic quotient.
A free pro-$\ell$ group on at least two generators is not
metabelian. One explicit finite quotient witnessing this is
the group generated by $1+X,1+Y$ in the units of
$\mathbf F_\ell\langle X,Y\rangle/(X,Y)^6$:
the degree-five part of
$[[1+X,1+Y],[[1+X,1+Y],1+X]]$, with $[a,b]=aba^{-1}b^{-1}$,
has coefficient $-1$
at $XYYXX$. This survives every prime. Thus the generated
subgroup is procyclic and $g$ centralizes $A_\ell$.
The image of a Sylow subgroup is a Sylow subgroup of the
image, so $\widetilde\theta_\ell(\widetilde D)$ has no
nontrivial pro-$\ell$ part. In $\Z_\ell^\times$ its exponent
divides $d$ (it is trivial for $\ell=2$). The character
$\widetilde\theta_\ell^d$ therefore descends through the
compact quotient $\widetilde N/\widetilde D\simeq G_L$;
restriction gives exactly (F.3.10a), continuously.

For the remaining conclusion choose an open characteristic
$V\subset F$ avoiding a generator of $A_\ell$ and put
$H=VA_\ell$. The nontrivial cyclic quotient $H/V$ makes
$A_\ell\to H^{\rm ab}$ nonzero; torsion-freeness of the
abelianization of the open free profinite group $H$ makes
its pro-$\ell$ part injective. The corresponding finite cover
$Y\to U$ descends over a finite extension. Intersecting
$\widetilde N$ with its open arithmetic group yields an
open arithmetic image $G_{L_1}$. Its geometric kernel lies
in $H$, hence acts trivially on $H^{\rm ab}$. The detected
line has an actual descended character
$\psi:G_{L_1}\to\Z_\ell^\times$ satisfying
\[
              \lambda_\ell|_{G_{L_1}}=\psi^d.
 \tag{F.3.10b}
\]
The homology map to $T_\ell J_{\overline Y}$ has boundary
kernel, a quotient of the cusp permutation module tensored
with $\Z_\ell(1)$. After splitting the cusp fields, a line
in that kernel has character $\chi_\ell$. Stix's criterion
then places $A_\ell$ in a cusp inertia group, and his
Proposition 16(3) places the abelian group $A$ in the same
inertia group; see \cite[Proposition 16, Definition 29 and
Theorem 30]{AuditStix2012}. Otherwise the line injects into
the rational Tate module of the Jacobian. It has Frobenius
weight one and Hodge--Tate exponents zero or one, by the
abelian-variety comparison used in B.1.4. A good-prime
Frobenius then has absolute value $q^{1/2}$ on $\psi$ and
$q$ on $\chi_\ell$, so no positive power of the former
equals the corresponding power of the latter.

This is precisely the surviving alternative in the full
source \cite[A.50, pp.33--34; A.22, pp.16--17]{AuditLM6Full2026}.
For $U=(y^2=x^3-x)\setminus\{O\}$ over $\Q(i)$, A.22
retains the existence or exclusion of a nontrivial closed
procyclic subgroup in the common kernel
$R=\bigcap_\ell\ker(F\to F^{(\ell)})$ with open arithmetic
normalizer image. Its further claimed restriction in A.36
is $R_{\rm nil}(F')\subset F''$. This audit verifies the
power-character calculation, its finite-cover alternative,
and its need for the stated enlargement; it does not assert
a solution of that residual subgroup problem or a new proof
of all the later LM6 restrictions. The archive's ``Open''
status for A.22 is retained.
\end{proof}

\begin{result}{F.3.11}{Abstract semilinear actions can fail effective descent}{Refuted; explicit counterexample proved}
An arbitrary automorphism-indexed abstract semilinear action
does not always arise from algebraic descent, even on an
affine finite \emph{\'etale} scheme with two geometric points.
This closes the unqualified reading in II.P16.02; it does
not affect the actual fpqc tensor-Isom construction in
II.P16.01.
\end{result}
\begin{proof}
Let $q_1,q_2,\ldots$ be distinct rational primes and let
$K=\bigcup_n\Q(\sqrt{q_1},\ldots,\sqrt{q_n})$.
Independence of their square classes gives
$\Gamma=\Gal(K/\Q)=\prod_{n\ge1}\mathbf F_2$.
The finite-support subgroup $D=\bigoplus_{n\ge1}\mathbf F_2$
is dense. The vector $(1,1,\ldots)$ has nonzero image in
the vector-space quotient $\Gamma/D$. Extend that vector
to an algebraic basis and take its coordinate functional.
This constructs a nonzero group homomorphism
$\eta:\Gamma\to\mathbf F_2$ with $D\subset\ker\eta$.
Its proper kernel is dense, so $\eta$ is not continuous.

On the algebra $K\times K$ let $\sigma\in\Gamma$ act by
applying $\sigma$ to both coefficients and then swapping the
two factors when $\eta(\sigma)=1$. The coefficient action
commutes with the swap, and the homomorphism identity for
$\eta$ proves the exact semilinear action law. Suppose this
action were obtained from a $\Q$-scheme by scalar extension.
Faithfully flat descent of finite \emph{\'etale} morphisms
would make that scheme finite \emph{\'etale} of degree two.
Its two geometric points are defined over a finite extension,
so their permutation action has open kernel. The specified
action on the two points has kernel $\ker\eta$, which is
not open, a contradiction. The same counterexample extends
to $\overline\Q/\Q$ by restriction of automorphisms to $K$.

The exact relation to the effective descent theorem is the
map from a continuous finite Galois action to its algebraic
descent datum. Our action lies outside its image because its
permutation map is discontinuous. Thus no assertion about
the effectiveness of the particular motivic cocycle is made
without its actual tensor-Isom or fpqc realization.
\end{proof}

\begin{result}{F.3.12}{Integral tensor forms and the finite-type finiteness theorem}{Known theorem; exact tensor and lattice maps proved}
Fix an integral tensor fibre functor $\omega_0$ and put
$\mathcal G=\Aut^\otimes(\omega_0)$. Its fppf-local tensor
forms are classified by $\mathcal G$-torsors through
\[
 [\omega]\longmapsto
       [\Isom^\otimes(\omega_0,\omega)].
 \tag{F.3.12a}
\]
For affine finite-type $\mathcal G$ over $\Z$, this is a finite
set. More generally the same conclusion holds over the ring
of $S$-integers of a number field for a finite set $S$.
No connectedness,
reductivity or smoothness of $\mathcal G$ is needed for this
statement. Over these Dedekind rings all fpqc-local forms
are already fppf-local, without a flatness assumption on
$\mathcal G$, and the same finiteness applies. An
unrestricted proalgebraic version is false.
\end{result}
\begin{proof}
The Isom sheaf in (F.3.12a) has the right action
$p\cdot g=p\circ g$ and is locally nonempty by the definition
of a form. It is a torsor. In the reverse direction a right
$\mathcal G$-torsor $P$ gives, object by object,
\[
                  \omega_P(X)=P\times^{\mathcal G}\omega_0(X).
 \tag{F.3.12b}
\]
Finite locally free module descent constructs these modules;
the given linear maps, unit, symmetry and tensor products
descend with them. Over a trivialization of $P$ the two
constructions are inverse, and faithful flatness descends
those inverse identifications. Thus the fppf-local forms
are exactly $H^1_{\rm fppf}(\Z,\mathcal G)$.
Here is the precise passage from fpqc to fppf over the
specified Dedekind ring $R$. A fpqc $\mathcal G$-torsor $P$
is affine and finitely presented, by fpqc descent of those
properties from $\mathcal G$. Choose a faithfully flat
$R$-algebra $B$ trivializing $P$. A section in $P(B)$
factors through a finitely generated $R$-subalgebra
$B_0\subset B$, because $P$ is finitely presented.
The $R$-module $B_0$ is torsion-free, since $B$ is flat;
over a Dedekind ring this implies flatness. For every
maximal ideal $\mathfrak m$ of $R$, the equality
$\mathfrak m B_0=B_0$ would imply
$\mathfrak m B=B$, contradicting faithful flatness of $B$.
Thus $B_0$ is faithfully flat. Noetherianity of $R$ makes
$B_0$ finitely presented, so $\Spec B_0\to\Spec R$ is an
fppf cover. Its section $s$ gives the actual trivialization
$\mathcal G_{B_0}\to P_{B_0}$, $g\mapsto s\cdot g$;
this is an isomorphism by the torsor identity, checked after
the fpqc cover. No flatness of $\mathcal G$ was used.
The primary theorem is
\cite[Proposition 5.1, author preprint pp.8--10]{AuditGilleMoretBailly2013}.
It proves finiteness for every affine finite-type group
scheme over the specified ring. The groups attached to
different forms are the inner forms
$\Aut^\otimes(\omega_P)=\Aut_{\mathcal G}(P)$, so only
finitely many isomorphism classes of their integral-point
groups arise in this setting.

Here is the equally explicit lattice formulation underlying
the Borel reference on 15:213. Let $G\subset\GL(V)$ be a
linear algebraic $\Q$-group, let $\Lambda_0\subset V$ be a
full lattice, and put
$K_f=\{g\in G(\mathbf A_f):g\widehat\Lambda_0=\widehat\Lambda_0\}$.
Then
\[
 G(\Q)\backslash G(\mathbf A_f)/K_f
 \longrightarrow
 \{\text{lattices in the $G$-genus of $\Lambda_0$}\}/G(\Q),
 \qquad [a]\longmapsto[V\cap a\widehat\Lambda_0]
 \tag{F.3.12c}
\]
is a bijection. Local lattice intersection recovers the
prescribed completions. Conversely two lattices in $V$
agree at all but finitely many primes; their local
$G(\Q_p)$-isomorphisms can therefore be chosen as the identity
there, giving an adele. Right multiplication by $K_f$ fixes
the local lattices; left multiplication by $G(\Q)$ is exactly
the stated global equivalence. This proves all identifications
in (F.3.12c). Its source is finite by
\cite[\S2.1 and Theorem 5.1, p.19]{AuditBorel1963}, including
nonconnected $G$; compact open subgroups are commensurable,
so the theorem's integral stabilizer can be replaced by $K_f$.
The corresponding discrete automorphism group is exactly
$G(\Q)\cap aK_fa^{-1}$.

For the unrestricted assertion take the tensor category of
finite free $\Z$-modules graded by
$D=\bigoplus_{i\ge1}\Z/2$. The usual forgetful functor has
group $\mathcal G=\prod_{i\ge1}\mu_2$. For every sign
sequence $\epsilon=(\epsilon_i)$ define a functor with the
same underlying modules and tensor constraint on homogeneous
degrees $a,b\in D$
\[
                 J^\epsilon_{a,b}
                     =\prod_i\epsilon_i^{a_i b_i}.
 \tag{F.3.12d}
\]
Only finitely many factors occur. The carry identity
$ab+(a\mathbin\oplus b)c=bc+a(b\mathbin\oplus c)$ for
$a,b,c\in\{0,1\}$ proves associativity of this constraint;
symmetry and the unit are immediate. These functors are
already fppf-locally tensor isomorphic: on the finite free
faithfully flat cover $\Z[i]/\Z$, choose $t_i=1$ when
$\epsilon_i=1$ and $t_i=i$ when $\epsilon_i=-1$, and use
the scalar $\prod_i t_i^{a_i}$ on degree $a$.
The identities $t_i^2=\epsilon_i$ give exactly the tensor
compatibility with (F.3.12d). A tensor isomorphism over
$\Z$ between two such functors has scalar
$u_i\in\Z^\times$ on the line of degree $e_i$. Its tensor
square forces $\epsilon_i=\epsilon_i' u_i^2=\epsilon_i'$.
Distinct sign sequences therefore give distinct integral
tensor forms. This supplies infinitely many actual forms,
with their isomorphism torsors, for one proalgebraic group.

The historical finite-$\alpha$ claim now has a precise
remaining gate: identify $\mathbb T_\alpha$ with the integral
form groupoid of an actual affine finite-type automorphism
model. Its stated fpqc locality then suffices by the
Dedekind-ring argument above. Once that identification is provided, (F.3.12a--c)
give the claimed finiteness. The arithmetic finiteness theorem
itself is no longer an unresolved part of I.M09.01. The
supplied pages 15:207--213 describe the gerbe and its
(pro)group but do not finish that integral-model construction.
\end{proof}

\begin{result}{F.3.13}{The integral automorphism model of an admissibly generated finite stage}{Proved for the specified exact category}
Let $R$ be noetherian, let $\mathcal C$ be a small rigid symmetric $R$-linear exact tensor category, and let
$\omega:\mathcal C\to\operatorname{Proj}_f(R)$ be an exact faithful tensor functor. Fix an object $M$ from which every object of $\mathcal C$ is obtained, up to isomorphism, by finitely many direct sums, tensor products, duals, and admissible subobjects and quotients. For an $R$-algebra $A$, put $\omega_A(X)=\omega(X)\otimes_R A$. Restriction to $M$ represents the functor
\[
 A\longmapsto\operatorname{Aut}^{\otimes}(\omega_A)
\]
by a closed subgroup scheme of $\operatorname{GL}_R(\omega M)$, hence by an affine group scheme of finite presentation over $R$.

For $R=\mathcal O_{F,S}$, with $F$ a number field and $S$ a finite set of finite primes, the isomorphism classes of fpqc-local tensor forms of this particular $\omega$ are consequently finite, by the explicit torsor equivalence and Dedekind-ring approximation in F.3.12. No smoothness, flatness, connectedness or reductivity of the resulting automorphism scheme is needed.
\end{result}
\begin{proof}
Write $V=\omega M$ and $C=\mathcal O(\operatorname{GL}_R(V))$. The universal automorphism $u$ and its inverse have coefficients in $C$. For each object choose one finite expression by the indicated operations, together with the chosen admissible sequences and final isomorphism. Tensor, dual and sum operations define matrices polynomial in the coefficients of $u,u^{-1}$.

At an admissible subobject $i:N\hookrightarrow P$, exactness gives
\[
 0\longrightarrow\omega N\xrightarrow{\omega i}\omega P\longrightarrow Q\longrightarrow0
\]
with $Q$ finite projective. This sequence splits as $R$-modules. Choose a retraction $r:\omega P\to\omega N$, and impose the polynomial equations
\[
 (1-(\omega i)r)\,u_P(\omega i)=0,
 \qquad
 (1-(\omega i)r)\,u_P^{-1}(\omega i)=0.
\]
Their meaning over every $A$ is preservation of $\omega_A N$ by $u_P$ and $u_P^{-1}$. The induced mutually inverse matrices on $\omega_A N$ are $r u_P(\omega i)$ and $r u_P^{-1}(\omega i)$. The split injection remains injective after every base change, so these are the unique matrices making naturality for $i$ hold. At an admissible quotient $q:P\twoheadrightarrow N$, its projective kernel and projective quotient similarly give a split sequence. Preservation of the kernel by $u_P,u_P^{-1}$ gives a unique induced automorphism; for a chosen section $s$ its matrix is $(\omega q)u_Ps$. These formulas handle any finite expression inductively, including successive subquotients. Finite projective modules may be written as images of idempotent matrices in finite free modules, so all displayed preservation conditions are ordinary coefficient equations even without a global basis.

Now impose, for every morphism $a:X\to Y$ of $\mathcal C$, the equations
\[
 u_Y\omega(a)=\omega(a)u_X,
\]
and impose the tensor, unit and dual compatibility equations. In particular these impose compatibility between the independently chosen expressions for $X\otimes Y$ and the tensor product of the chosen expressions for $X,Y$. Let $I\subset C$ be the ideal generated by all these equations and the subquotient preservation equations. The construction gives, for every $R$-algebra $A$, mutually inverse maps between $\operatorname{Hom}_{R\text{-alg}}(C/I,A)$ and tensor automorphisms of $\omega_A$: a point supplies the constructed components; an automorphism supplies its component at $M$; naturality and split exactness force all its other components. Thus the assertion is scheme-theoretic and includes algebras with nilpotents and torsion. It is not merely a calculation on field-valued points.

The algebra $C$ is finitely presented over noetherian $R$, hence noetherian. Its ideal $I$ is finitely generated. Therefore $\operatorname{Spec}(C/I)$ is of finite presentation over $R$. Composition and inversion of the natural automorphisms restrict those of $\operatorname{GL}(V)$, giving the claimed closed subgroup scheme. Finally apply F.3.12 to this actual automorphism scheme and the actual map $\omega'\mapsto\operatorname{Isom}^{\otimes}(\omega,\omega')$.
\end{proof}

\paragraph{Historical boundary.}
This construction removes the extra representability hypothesis for the exact finite stages just defined. The supplied 15:213 uses $\mathbb M_\alpha$ and $T_\alpha$, and 15:223 calls $G_M$ an algebraic group over $\mathbb Z$. Identifying those historical stages with this admissibly generated integral category still requires their definition; finite generation of the rational category alone does not supply the split integral presentations used in this proof. The unrestricted proalgebraic tensor category in F.3.12 has explicitly infinitely many forms.

\begin{result}{F.3.14}{A precise hypersurface normalization construction}{Proved; historical identification retained}
Every smooth integral projective $d$-fold $X$ over a characteristic-zero field $k$ admits a finite birational morphism
\[
 \nu:X\longrightarrow Y\subset\mathbf P_k^{d+1}
\]
onto an integral hypersurface, and $\nu$ identifies $X$ with the normalization of $Y$.
This is an actual geometric construction relevant to the proposed use of singular hypersurfaces. The pullback on ordinary cohomology need not be an isomorphism: the explicit nodal cubic below has $H^1(Y(\C),\Q)=\Q$ and $H^1(X(\C),\Q)=0$, with the kernel specified by its normalization sequence.
\end{result}
\begin{proof}
The birational projection statement is explicitly present in the original EGA~V prenotes, Corollary~14.6, with its preceding Proposition~14.5 and the following warning that the projection need not be an immersion \cite[original scan p.89]{AuditGrothendieckEGAVOriginal}. We give the construction and normalization proof in the stated characteristic-zero scope.

Choose a very ample bundle $L$ and a projective embedding defined by a basis of $H^0(X,L)$. Because $k$ is infinite, choose $d+1$ sections $s_0,\ldots,s_d$ having no common zero. One can construct them successively: at each stage choose a section avoiding all the finitely many generic points of the current intersection; after $d$ cuts the intersection is finite, and the last section avoids those points. This gives
\[
 g=[s_0:\cdots:s_d]:X\longrightarrow\mathbf P_k^d,
 \qquad g^*\mathcal O(1)=L.
\]
A positive-dimensional geometric fibre would contain a projective integral curve, on which $L$ would both have positive degree and be trivial. Thus $g$ is quasi-finite; it is proper and hence finite \cite[Tag 02LS, Lemma 37.44.1]{AuditStacksFiniteNormalization}. Its image is closed of dimension $d$, so it is all of $\mathbf P^d$. Consequently
\[
 F=k(s_1/s_0,\ldots,s_d/s_0)\subset E=k(X)
\]
is a finite separable extension.

Complete the sections to a basis. Their ratios $r_j=s_j/s_0$ generate $E$ over $F$. There is a $k$-linear combination $r=\sum_j a_jr_j$ which is a primitive element. Indeed, for every pair of distinct $F$-embeddings $\tau,\tau'$ of $E$ into a separable closure, the equation
\[
 \sum_j a_j\bigl(\tau(r_j)-\tau'(r_j)\bigr)=0
\]
defines a proper $k$-linear subspace of $k^N$. Finitely many such proper subspaces do not cover $k^N$. Choose the coefficients outside their union and set $t=\sum_j a_js_j$. The morphism
\[
 \nu=[s_0:\cdots:s_d:t]:X\longrightarrow\mathbf P_k^{d+1}
\]
is everywhere defined, has pullback $\mathcal O(1)$ equal to $L$, and is finite by the same fibre argument. Its scheme-theoretic image $Y$ is integral of dimension $d$. A homogeneous height-one prime in the polynomial coordinate ring of $\mathbf P^{d+1}$ is principal, so $Y$ is a hypersurface. The primitive-element choice gives $k(Y)=F(r)=E$, hence $\nu$ is birational.

For an affine open $\operatorname{Spec}A\subset Y$, write its inverse image as $\operatorname{Spec}B$. The ring $B$ is finite over $A$, has fraction field $E$, and is integrally closed because $X$ is smooth. Thus $B$ is the integral closure of $A$ in $E$: every element integral over $A$ is also integral over $B$, hence belongs to $B$, and the reverse inclusion follows from finiteness. This proves the normalization assertion with its exact map; compare the integral-closure argument of \cite[Tag 0AB1, Lemma 29.55.8]{AuditStacksFiniteNormalization}.

For the cohomological boundary, take
\[
 Y=\{y^2z=x^2(x+z)\}\subset\mathbf P^2_\C,
 \qquad
 \nu([u:v])=[v(u^2-v^2):u(u^2-v^2):v^3].
\]
The formula has no common zero and parametrizes $Y$ birationally. It is its normalization, with the two points $[1:1]$ and $[-1:1]$ over the node $P=[0:0:1]$. On the complex analytic spaces there is an exact sequence
\[
 0\longrightarrow\Q_Y\longrightarrow\nu_*\Q_{\mathbf P^1}
 \xrightarrow{\delta}\Q_P\longrightarrow0,
 \qquad \delta(a,b)=a-b
\]
at the node; away from the node the first map is an isomorphism. Since $\nu$ is finite, its higher direct images on these constant sheaves vanish. The map from global constants to $H^0(Y,\Q_P)=\Q$ is zero. The long exact sequence therefore gives
\[
 H^0(Y,\Q_P)=\Q\xrightarrow{\sim}H^1(Y(\C),\Q),
 \qquad
 \nu^*:H^1(Y(\C),\Q)\longrightarrow H^1(\mathbf P^1(\C),\Q)=0.
\]
Thus the geometric reduction comes with specified normalization data; it does not alone provide an ordinary-cohomology identification for the singular image.

The original cote14 scans111--113 (archive pp110--112) contain the global cycle-Lefschetz restrictions, Gysin maps, and punctured-cone formulations; scan131 (archive p130) contains a generic-member/base-locus diagram and its cycle-image formula. These inspected pages do not yet identify a proof carrying out the particular singular-hypersurface proposal in the 1967 Coates letter. I.M05.01 and II.MB08.01 retain that exact documentary obligation. The present construction is independently proved here, not attributed to an unseen manuscript passage.
\end{proof}

\begin{result}{F.3.15}{CM pro-prime normalizers: an independent primary-source proof}{Proved; the separate integral lifting claim is not used}
Let $K=\Q(i)$, $E_0:y^2=x^3-x$, $U=E_0\setminus\{O\}$, and
\[
 F=\pi_1(U_{\overline K}),\qquad P_\ell=F^{(\ell)},\qquad
 E_\ell=\pi_1(U)/\ker(F\to P_\ell).
\]
For every prime $\ell$, a nontrivial closed procyclic subgroup $A\subset P_\ell$ whose normalizer in $E_\ell$ has open image in $G_K$ is contained in an actual cusp inertia subgroup of $P_\ell$.
\end{result}
\begin{proof}
This independently re-proves the recognition assertion of LM6 A.20. The arithmetic nonvanishing input is Ishii's author version arXiv:2312.04175v2, with its exact standing hypotheses checked below; the final journal full text was not obtained. The distinct integral class-four lifting claim at $\ell=5$ in the prior return is not an input or a conclusion here.
Write $N=N_{E_\ell}(A)$ and let $p:E_\ell\to G_K$ be the natural projection.

\paragraph{1. Actual finite-cover homology detects a generator.}
The free pro-$\ell$ group $P_\ell$ is torsion-free, so $A\simeq\mathbf Z_\ell$; fix a generator $u$. Choose a finite $\ell$-group quotient $P_\ell\to Q$ in which $u$ has nonidentity image $\bar u$. For
\[
 H=\operatorname{preimage}(\langle\bar u\rangle)\subset P_\ell,
\]
the composite $H\to\langle\bar u\rangle\to\mathbf Z/\ell$ sends $u$ to a nonzero element. Thus $u$ has nonzero image in the torsion-free abelianization $H^{\rm ab}$.

The inverse image of $H$ in $F$ defines an actual finite etale geometric cover $Y\to U_{\overline K}$. It descends to a finite extension of $K$. Moreover $\pi_1(Y_{\overline K})^{(\ell)}\simeq H$. To check this exact identification, any finite $\ell$-group quotient of $\pi_1(Y_{\overline K})$ has normal closure over $U_{\overline K}$ with an $\ell$-group as Galois group: take the Galois closure of $Y$, which is already an $\ell$-cover, and the compositum of its finitely many conjugate $\ell$-covers. The resulting group is a subgroup of an iterated wreath product of finite $\ell$-groups. Hence all these quotients factor through $P_\ell$, proving the asserted maximal-pro-$\ell$ identification.

Take an arithmetic covering subgroup for a descent of $Y$ and intersect it with $N$. Its arithmetic image is still open, since that covering subgroup is open. Its geometric kernel acts inner on $H$ and hence trivially on $H^{\rm ab}$. Therefore the nonzero line spanned by $u$ in
\[
 H^{\rm ab}\otimes\mathbf Q_\ell
 \simeq H_1(Y_{\overline K},\mathbf Q_\ell)
\]
is stable for the ordinary Galois action of a finite number-field extension. This construction uses neither a section of $N$ nor a splitting of an arithmetic extension. The homology exact sequence for the smooth compactification $\overline Y$ has boundary term of weight two and proper term of weight one. Thus the normalizer character on this line has Frobenius weight one or two.

\paragraph{2. The first nonzero term has degree one or two.}
Use the completed free Lie algebra over $\mathbf Q_\ell$ on $V=H_1(E_{0,\overline K},\mathbf Q_\ell)$. The pro-$\ell$ Magnus embedding is injective, and its logarithm gives a nonzero series $U_u=\log u$ with first nonzero homogeneous term $u_d\in L_d(V)$. For clarity, the Magnus embedding here comes from the canonical completed-group-algebra identification
\[
 \mathbf Z_\ell[[P_\ell]]\simeq\mathbf Z_\ell\langle\!\langle X,Y\rangle\!\rangle,
 \qquad x_0\mapsto1+X,\quad y_0\mapsto1+Y.
\]
The group is separated by its finite quotients, so its inclusion in the units of its completed group algebra is faithful. After the formal substitutions $x=\log(1+X)$, $y=\log(1+Y)$, the logarithm of a group-like element is primitive, hence lies in the completed free Lie algebra. No integral division by two or by $\ell$ is used to infer a group identity.

Conjugation by a geometric element acts trivially on the first homogeneous term. If $n\in N$ acts by $u\mapsto u^{\theta(n)}$, then
\[
 p(n)(u_d)=\theta(n)u_d.
\]
The character therefore descends to $p(N)$, and at a good arithmetic Frobenius of residue cardinality $q$, its complex absolute value is $q^{d/2}$, since $V$ has weight one. Comparison with the actual detector in Step~1 gives $d=1$ or $d=2$.

\paragraph{3. Degree one at inert primes and at two.}
The curve has trace $-2$ over $\mathbf F_5$, so a good Frobenius at a prime over5 has eigenvalues
\[
 \alpha=-1+2i,\qquad\bar\alpha=-1-2i.
\]
Their ratio is not a root of unity: it lies in $\mathbf Q(i)$ and is none of its four roots of unity $1,-1,i,-i$. A positive power of this Frobenius lies in any prescribed open subgroup of $G_K$. When $\ell=2$ or $\ell$ is inert in $K$, $i\notin\mathbf Q_\ell$, and no positive power of either eigenvalue belongs to $\mathbf Q_\ell$, since equality with its conjugate would make $(\alpha/\bar\alpha)^n=1$. Thus the characteristic polynomial of every such power is irreducible over $\mathbf Q_\ell$, and $V$ has no one-dimensional subspace stable under an open arithmetic subgroup. This excludes $d=1$ at those primes.

\paragraph{4. The split-prime source input and its exact degree.}
Now let $\ell$ split in $K$, hence $\ell\ge5$. The CM eigenspaces give a decomposition $V=\mathbf Q_\ell x\oplus\mathbf Q_\ell y$, with $[i]$ acting by $i$ and $-i$. Every line stable under an open arithmetic subgroup is one of these two lines. Indeed the two CM characters remain distinct after finite extension; the preceding Frobenius~5 argument proves this for $\ell\ne5$, and at $\ell=5$ use Frobenius~13, whose eigenvalues are $3\pm2i$ and whose ratio is again not a root of unity.

Put $K_\infty=K(E_0[\ell^\infty])$. Nakamura's (4.3)--(4.4) gives\cite[p.~214]{AuditNakamura1995} a canonical representative of every outer Galois action from $G_{K_\infty}$ which fixes the boundary and is the identity through derivation degree three:
\[
 G_{K_\infty}\longrightarrow\Gamma^*_{1,1}(4).
\]
This particular representative matters: reduction to a metabelian group without retaining the boundary-fixing representative would lose the following information. Nakamura's (4.7)--(4.9) defines the comparison
\[
 \operatorname{gr}_4\Gamma^*_{1,1}\longrightarrow
 \operatorname{gr}_4\Psi^*\simeq\operatorname{Sym}^2V\otimes\det(V).
\]
It is surjective by Lemma~4.10 for odd $\ell$\cite[p.~217]{AuditNakamura1995}. The exact sequence (4.9.2) gives
\[
 \dim\operatorname{gr}_4\Gamma^*_{1,1}
 =2\dim L_5(V)-\dim L_6(V)=2\cdot6-9=3.
\]
The target also has dimension~3, so the displayed map is an isomorphism after tensoring with $\mathbf Q_\ell$.

For $[i]$, the three weight spaces in its target have characters $(3,1),(2,2),(1,3)$ and eigenvalues $-1,1,-1$. The arithmetic outer action commutes with the $K$-defined endomorphism $[i]$, so its degree-four image lies in the middle line. Ishii's Corollary~4.10 gives a nonzero character\cite[p.~21]{AuditIshii2024v2} $\kappa_{(2,2)}$ on $G_{K_\infty}$. All standing assumptions hold: $K=\mathbf Q(i)$ has class number one, the full CM is defined over $K$, and $\ell\ge5$ splits. The curve has discriminant~64, so it has good reduction at $\ell$. Its Hasse invariant is the coefficient of $x^{\ell-1}$ in $(x^3-x)^{(\ell-1)/2}$, namely $(-1)^{(\ell-1)/4}\binom{(\ell-1)/2}{(\ell-1)/4}$ for $\ell\equiv1\pmod4$; this is nonzero modulo $\ell$, verifying ordinary reduction directly. His Theorem~1.3, equivalently Nakamura's Corollary~4.12 with the shifted index convention, places this character in the coefficient of $U_1U_2$ in the quadratic part of the comparison power series, multiplied by the nonzero factor $(1-\ell^2)^{-1}$\cite[Theorem~1.3,p.~3]{AuditIshii2024v2}. Thus the middle degree-four outer action is nonzero.

A nonzero continuous homomorphism to the torsion-free additive group $\mathbf Z_\ell$ remains nonzero on every open subgroup: otherwise its image would be finite. In particular there exists
\[
 \sigma\in G_{K_\infty}\cap p(N)
\]
with nonzero middle degree-four action. This supplies the required arithmetic element without choosing a section of $N\to p(N)$.

\paragraph{5. Independent calculation of the middle derivation.}
Let $\omega=[x,y]$, $F_0(x)=y$, $F_0(y)=0$, and define
\[
 \epsilon_4(x)=\operatorname{ad}_x^4y,\quad
 \epsilon_4(y)=[y,\operatorname{ad}_x^3y]-[[x,y],\operatorname{ad}_x^2y],
 \qquad T_4=[F_0,\epsilon_4].
\]
Direct expansion gives $T_4(\omega)=0$ and, in standard Lyndon brackets,
\[
 T_4(x)=-2b_{xxxyy}-b_{xxyxy},\qquad
 T_4(y)=2b_{xxyyy}+3b_{xyxyy}.
\]
The three special derivations $\epsilon_4,T_4,[F_0,T_4]$ form a basis: the kernel of $(D(x),D(y))\mapsto[D(x),y]+[x,D(y)]$ has dimension~3, and these three pairs are linearly independent. In the metabelian Lie quotient, with $L'/L''=\mathbf Q_\ell[X,Y]\omega$, their pairs have the form
\[
 (D(x),D(y))=(X\beta\omega,Y\beta\omega),\qquad
 \beta=X^2,\ 2XY,\ 2Y^2
\]
respectively. Hence their comparison determinant is4 and the middle line is exactly $\mathbf Q_\ell T_4$. The sign of $\omega$ relative to Nakamura's inverse-commutator boundary changes only a nonzero scalar and is immaterial to nonvanishing.

Most importantly,
\[
 T_4(x)\notin[x,L_4],\qquad T_4(y)\notin[y,L_4].
\]
In the relevant bidegree, each bracket image has dimension~1, and adjoining the indicated vector raises the rank to2. The independently written exact script \texttt{check\_CM\_degree\_four.py} verifies the kernel dimension, all three special derivations, the three metabelian polynomials, and both rank obstructions. It also independently counts the reductions at~5 and~13. These computations are finite algebra certificates; the preceding primary-source result supplies the arithmetic nonvanishing.

\paragraph{6. The full conjugacy equation, including lower inner terms.}
Suppose $d=1$ and, after interchanging the two CM lines if needed, write
\[
 U_u=\log u=a x+\text{terms of degree at least2},\qquad a\in\mathbf Q_\ell^\times.
\]
No primitivity of the original abelian image is assumed. Choose the $\sigma$ in Step~4 and a lift $n\in N$. Since $\sigma$ acts trivially on $V$ and $a\ne0$, the normalizer exponent is $\theta(n)=1$. Let $\widetilde\sigma$ denote Nakamura's boundary-fixing representative. Its first nonidentity derivation is $cT_4$ with $c\ne0$, whereas the representative furnished by $n$ differs from it by an inner automorphism. Thus, in the rational completed free Lie algebra,
\[
 \widetilde\sigma(U_u)=e^{\operatorname{ad} W}U_u
\]
for some $W$.

Comparing degree~2 first gives $[W_1,x]=0$, hence $W_1=b x$. Replace the conjugating element $\exp W$ on the right by
\[
 \exp W\,\exp\bigl(-(b/a)U_u\bigr).
\]
This preserves its conjugation action on $U_u$ and removes its degree-one term. Comparing degrees3 and4 now forces the degree-two and degree-three terms to vanish. The fact used here is $C_L(x)=\mathbf Q_\ell x$: in the free associative algebra a homogeneous element commuting with $x$ is a scalar multiple of a power of $x$, and among these only the degree-one term is a nonzero primitive Lie element in characteristic zero.

Consequently the degree-five equation is
\[
 acT_4(x)=[W_4,a x],
\]
contradicting Step~5. The same argument with $y$ uses the second rank obstruction. This excludes degree one for every split prime, and completes the part of the prior argument that had left its inner correction implicit.

\paragraph{7. Degree two gives the actual cyclotomic action and actual inertia.}
Thus $d=2$. Since $L_2(V)=\mathbf Q_\ell\omega$, write $u_2=a\omega$ with $a\ne0$. For every $n\in N$, comparison on this term gives
\[
 a\theta(n)\omega=a\chi_\ell(p(n))\omega.
\]
Canceling the nonzero scalar yields $\theta=\chi_\ell\circ p$ on the actual normalizer. Nakamura's Theorem~2.1.1\cite[p.~94]{AuditNakamura1994}, applied to the full class of finite $\ell$-groups, now puts $u$ and $A$ in an actual cusp inertia subgroup of $P_\ell$. The conclusion concerns the pro-$\ell$ group itself; the rational Lie algebra was used only to rule out the degree-one case.

\end{proof}

\paragraph{Verification boundary.}
The displayed proof also supplies the two steps previously left implicit: the ordinary-reduction condition in Ishii's input, and the removal of a lower inner conjugator when the original abelian image of $u$ is not primitive. It uses no chosen splitting of the arithmetic normalizer. The exact numerical normalization $\kappa/48$ and the integral class-four lift at~5 remain outside this re-proof's scope. The full profinite common-kernel subgroup in LM6 A.22 remains to be constructed or excluded.

\begin{result}{F.3.16}{A finite-cover CM Prym and its actual punctured quotient maps}{Proved; a precise special case of the residual homology test}
For $U=(y^2=x^3-x)\setminus\{O\}$ over $K=\Q(i)$, there is a degree-$32$ finite etale cover $Y\to U$ whose smooth compactification has genus~3 and whose proper Prym contains an elliptic curve isogenous to $E_0:y^2=x^3-x$. Thus a general assertion that these Pryms have no CM weight-one characters is false.

For this same cover, every closed procyclic subgroup of $\pi_1(Y_{\overline K})$ with open arithmetic normalizer image has zero image in the proper homology $H_1(\overline Y,\Z_\ell)$ for every $\ell$. Whenever it has nonzero image in $H_1(Y_{\overline K},\Z_\ell)$ for some $\ell$, it is contained in one cusp inertia group. These assertions use the explicit maps on the same punctured curve constructed below.
\end{result}
\begin{proof}
Take the smooth plane quartic
\[
 C:\ x^4+y^4=z^4
\]
and the smooth proper genus-one curve $D$ with affine equation
\[
 s^2=1-t^4,\qquad O_D=(0,1).
\]
The rational formulas
\[
 q_x:C\longrightarrow D,\quad (t,s)=(y/z,x^2/z^2),\qquad
 q_y:C\longrightarrow D,\quad (t,s)=(x/z,y^2/z^2)
\]
extend uniquely to morphisms of smooth proper curves. Both have degree two: their function-field extensions adjoin respectively a square root of $s$, and the displayed sign-changing involutions are nontrivial deck involutions.

Here $D$ is the original CM elliptic curve over $K$, with a completely explicit isomorphism. Put
\[
 X=\frac{2(1+s)}{t^2},\qquad Y=\frac{4(1+s)}{t^3}.
\]
Then $Y^2=X^3+4X$. The inverse on the dense open set is
\[
 t=\frac{2X}{Y},\qquad s=\frac{2X^3}{Y^2}-1.
\]
Thus this is an isomorphism of smooth proper genus-one curves, sending $O_D$ to the point at infinity. With $a=1+i$, one has $a^4=-4$, so $X=a^2x_E$, $Y=a^3y_E$ identifies the cubic with $E_0$ over $K$.

The four branch values of $q_x$ are precisely $s=0$, $t^4=1$. Indeed the divisor of $s$ has simple zeroes at those four points and poles of order two at each of the two points at infinity of $D$; adjoining $\sqrt{s}$ therefore ramifies exactly at its zeroes. These branch values belong to $D[4]$. To see this without a torsion-coordinate assertion, let
\[
 j(t,s)=(-t,s),\qquad h(t,s)=(t,-s),\qquad P=(0,-1).
\]
The involution $j$ fixes the origin and is nontrivial, hence is $[-1]$ on the elliptic curve. It fixes $P$, so $2P=0$. The involution $h$ sends $O_D$ to $P$ and acts by $-1$ on the regular differential $dt/s$. An automorphism of an elliptic curve is the translation by its value at the origin followed by an origin-preserving automorphism; in characteristic zero the derivative determines the latter. Hence $h(Q)=P-Q$. Its fixed points, exactly the four branch values above, satisfy $2Q=P$ and thus $4Q=0$.

Consequently
\[
 f=[4]\circ q_x:C\longrightarrow D
\]
is a degree-$32$ morphism branched only over $O_D$, since $[4]$ is etale. It induces an actual finite etale map
\[
 C\setminus f^{-1}(O_D)\longrightarrow D\setminus\{O_D\}\simeq U.
\]
There are $28$ removed geometric points: four branch values of $q_x$ in $D[4]$ have one point above each, and the other twelve points of $D[4]$ have two. Thus $2-2g(C)-28=-32$, agreeing with the degree and $\chi(U)=-1$.

Write $q_{x,*}:J_C\to D$ for the norm map and $q_y^*:D\to J_C$ for pullback. The involution $\iota_x(x,y,z)=(-x,y,z)$ fixes $q_x$ and satisfies $q_y\iota_x=[-1]q_y$. Consequently
\[
 q_{x,*}\iota_x^*q_y^*=q_{x,*}q_y^*,\qquad
 \iota_x^*q_y^*=-q_y^*.
\]
It follows that $2q_{x,*}q_y^*=0$, hence
\[
 q_{x,*}q_y^*=0.
\]
The final implication is exact: the image of a morphism of abelian varieties is connected, whereas the image of a morphism annihilated by two lies in the finite etale group $D[2]$. Also $q_{y,*}q_y^*=[2]$, the norm--pullback identity for a degree-two map. Therefore, for every prime $\ell$,
\[
 q_y^*:T_\ell D\hookrightarrow T_\ell J_C,
 \qquad \operatorname{im}(q_y^*)\subset
 \ker\bigl(f_*:T_\ell J_C\to T_\ell D\bigr),
 \quad f_*=[4]q_{x,*}.
\]
Injection follows by applying $q_{y,*}$ and using torsion-freeness of $T_\ell D$, including $\ell=2$. For odd $\ell$, $\tfrac12q_{y,*}$ is an integral retraction. The connected image $q_y^*(D)$ lies in the connected Prym kernel of $q_{x,*}$ and is isogenous to $D\simeq E_0$.

All these maps are defined over $K$. For every rational prime $\ell$ split in $K$, the two idempotents of $K\otimes\mathbf Q_\ell\simeq\mathbf Q_\ell\times\mathbf Q_\ell$ give $G_K$-stable one-dimensional summands of $T_\ell D\otimes\mathbf Q_\ell$. Their images are actual weight-one CM lines in this proper Prym. This is a statement about a finite cover and its homology. It does not lift either line to a procyclic subgroup of $\pi_1(U_{\overline K})$, does not place such a lift in $R_{\rm nil}(F')$, and does not close A.22. The construction refutes precisely the suggested exclusion of all CM characters from all of the relevant finite-cover Pryms.

\paragraph{The actual punctured maps eliminate this cover's proper homology alternative.}
We now apply the independently proved recognition theorem F.3.15 to $E_0$. The preceding CM Prym is not an example where one must artificially add punctures to obtain the elliptic quotient map: the required maps exist on the same open curve.

Let $Y=C\setminus f^{-1}(O_D)$. The map $q_y$ restricts to $Y\to D\setminus\{O_D\}$, since $q_y^{-1}(O_D)$ consists of $(0:1:1)$ and $(0:-1:1)$, whose $q_x$ images are the four-torsion points $(t,s)=(\pm1,0)$. The map $q_x$ itself also restricts to that punctured target because $q_x^{-1}(O_D)\subset f^{-1}(O_D)$.

The third quotient is
\[
 D_z:s_z^2=t_z^4+1,\quad O_z=(0,1),\qquad
 q_z:C\to D_z,\quad(t_z,s_z)=(x/y,z^2/y^2).
\]
Its origin has inverse image $(0:1:1),(0:1:-1)$, both in $f^{-1}(O_D)$. Thus $q_z$ restricts to $Y\to D_z\setminus\{O_z\}$ as well. The same formulas $X=2(1+s_z)/t_z^2$, $Y_E=4(1+s_z)/t_z^3$ give $Y_E^2=X^3-4X$. Scaling by $a=\sqrt2$, with $X=a^2x_E$, $Y_E=a^3y_E$, identifies $D_z$ with $E_0$ over $K'=K(\sqrt2)$. Passage to this finite field extension preserves the open arithmetic normalizer condition by taking the inverse image of $G_{K'}$.

The product of the three norm maps
\[
 J_C\longrightarrow D\times D\times D_z
\]
is an isogeny. Here is an explicit verification. On the affine chart $z=1$, write
\[
 \omega_1=\frac{dx}{y^3},\quad
 \omega_2=\frac{x\,dx}{y^3},\quad
 \omega_3=\frac{y\,dx}{y^3}.
\]
These are the usual three regular differentials on the smooth plane quartic. Differentiating $x^4+y^4=1$ gives $dy=-x^3dx/y^3$, and hence
\[
 q_x^*(dt/s)=-\omega_2,\qquad
 q_y^*(dt/s)=\omega_3,\qquad
 q_z^*(dt_z/s_z)=\omega_1.
\]
The pullback of differentials is an isomorphism. The product norm is consequently an isogeny in characteristic zero, and its map on the product of rational Tate modules is an isomorphism.

Let $A$ be a closed procyclic subgroup of the geometric fundamental group of $Y$, with open arithmetic normalizer image. For every prime $\ell$, its image in $H_1(C_{\overline K},\mathbf Z_\ell)$ is zero. Indeed, nonzero proper homology has nonzero image under one of the three quotient maps, by the isogeny just proved. The corresponding homomorphism of actual arithmetic fundamental groups carries $A$ to a procyclic subgroup of the maximal pro-$\ell$ geometric group of a copy of $E_0\setminus\{O\}$, and the normalizer still has open arithmetic image. The image is nonzero because its proper abelianization is nonzero. F.3.15 makes it cuspidal, whereas inertia at the unique puncture maps to zero in proper $H_1(E_0,\mathbf Z_\ell)$. This is a contradiction.

It follows in particular that whenever such an $A$ has nonzero image in $H_1(Y_{\overline K},\mathbf Z_\ell)$, that image lies in the boundary submodule. After a finite extension splitting the boundary points it is acted on through the cyclotomic character. A nonzero homomorphism from the torsion-free procyclic $A_\ell\simeq\mathbf Z_\ell$ to this torsion-free module is injective, so its action is the actual cyclotomic action on $A_\ell$. Nakamura's full profinite recognition theorem\cite[Theorem~2.1.1,p.~94]{AuditNakamura1994} then places $A_\ell$ in a cusp inertia group of $Y$, and the cusp centralizer statement places all of $A$ there\cite[Lemma~2.1.2,p.~95]{AuditNakamura1994}. Here $Y$ has~28 punctures, so each peripheral generator belongs to a free generating system and the cited primitive-generator lemma applies directly. This is an exact finite-cover homology recognition statement, not the assertion that this one cover detects every nontrivial subgroup in LM6 A.22.

\end{proof}

\begin{result}{F.3.17}{The precise reach and limitation of this cover tower}{Proved; no arithmetic normalizer is asserted for the invisible elements}
Let $F=\pi_1(U_{\overline K})$, $P_p=F^{(p)}$,
\[
 R=\bigcap_p\ker(F\to P_p),\qquad S=R_{\rm nil}(F').
\]
The cover of F.3.16 and the Fermat power tower below have geometric covering subgroups containing $\ker(F\to P_2)$ and hence $R$. There are actual nontrivial closed procyclic subgroups of $S$ whose images in the first homology of every member of this tower vanish. More generally this holds for every family of open subgroups of $F$ containing $R$. The exhibited subgroups are not asserted to have open arithmetic normalizer image, so this conclusion does not decide LM6 A.22.
\end{result}
\begin{proof}
The subgroup $H=\pi_1(Y_{\overline K})\subset F$ contains $\ker(F\to P_2)$, and therefore contains every candidate $C\subset R$ in A.22. Indeed, on geometric function fields the degree-$16$ multiplication cover $[4]:D\to D$ is Galois with group $D[4]\simeq(\mathbf Z/4)^2$. Over its function field, $q_x$ is quadratic. The normal closure of the composition is obtained by adjoining the at most sixteen conjugate square roots. Its Galois group fits in
\[
 1\longrightarrow A\longrightarrow G\longrightarrow D[4]\longrightarrow1,
 \qquad A\subset(\mathbf Z/2)^{16},
\]
so $G$ is a finite $2$-group. The subgroup corresponding to the given connected cover contains the kernel of this quotient of $F$, and hence the kernel of the maximal pro-$2$ quotient.

The Fermat power maps give an explicit infinite family of further covers. For $N=4\cdot2^r$, put
\[
 C_N:\ x^N+y^N=z^N,\qquad
 a_N:C_N\to C,\quad(x:y:z)\mapsto(x^{N/4}:y^{N/4}:z^{N/4}),
 \qquad f_N=f\circ a_N.
\]
The map $a_N$ is Galois over $\overline K$ with group $(\mu_{N/4})^2$ and is ramified only above the coordinate-zero points of $C$. All those points map under $q_x$ to $D[4]$: $x=0$ was checked above; $y=0$ gives $O_D$ or $P\in D[2]$; and the two points of $D$ at infinity are fixed by $j=[-1]$, hence are in $D[2]$. Thus
\[
 Y_N=C_N\setminus f_N^{-1}(O_D)\longrightarrow U
\]
is finite etale. Iterating the normal-closure argument shows that its geometric Galois closure is a $2$-group. Every corresponding subgroup $H_N$ therefore contains $\ker(F\to P_2)$ and $R$.

This family cannot, as a purely geometric assertion, detect every nontrivial subgroup of the remaining kernel by ordinary first homology. In fact, for \emph{any} family of open subgroups $H_i\subset F$ containing $R$, one has the actual inclusions
\[
 R'\subset H_i',\qquad R'\longrightarrow H_i^{\rm ab}=0.
\]
Here primes mean closed commutator subgroups. The subgroup $S'$ is a nontrivial subgroup of $R'$ and already lies in the refined A.22 kernel $S$. To prove nontriviality, choose a surjection $\varphi:F\twoheadrightarrow A_5$, using a $5$-cycle and a $3$-cycle as images of the two free generators. Since $A_5$ is perfect, $\varphi(F')=A_5$. The quotient $F'/S$ is pronilpotent, so
\[
 A_5/\varphi(S)
\]
is a finite nilpotent quotient of the perfect group $A_5$, and is trivial. Hence $\varphi(S)=A_5$ and $\varphi(S')=A_5' = A_5$, proving $S'\ne1$. Moreover $S\subset R$: every map $F\to P_p$, restricted to $F'$, has pro-$p$ target, and therefore kills $S$.

Choose $g\in S'$ with $\varphi(g)\ne1$. Then $\overline{\langle g\rangle}$ is an actual nontrivial closed procyclic subgroup of $S'\subset S$, and its image in every $H_i^{\rm ab}$ is zero. This construction does \emph{not} assert that its arithmetic normalizer has open image. The missing arithmetic-normalizer requirement remains exactly the A.22 requirement. It proves the precise limitation of an argument using only the displayed covers, not a counterexample to A.22 itself. A successful finite-cover detector for all candidates must use the normalizer condition in an additional way, or range beyond this family.
\end{proof}

\begin{result}{F.3.18}{A rank-one Jacobian character gives an actual CM quotient}{Proved; exact quotient and Prym maps supplied}
Let $J$ be an abelian variety over a number field $L$, let $\ell$ be a prime,
and let $W\subset V_\ell(J)$ be a nonzero one-dimensional
$\mathbf Q_\ell$-subspace stable under $G_L$.
There is an $L$-simple abelian variety $B$ and an actual surjective
$L$-homomorphism $q:J\to B$ such that $V_\ell(q)|_W\ne0$ and
\[
 \operatorname{End}^0_L(B)=E,\qquad [E:\mathbf Q]=2\dim B.
\]
Thus $B$ is of CM type, with its CM endomorphisms defined over $L$.
There is a place $\lambda\mid\ell$ of $E$ with $E_\lambda=\mathbf Q_\ell$;
the image of $W$ is the corresponding one-dimensional summand of
$V_\ell(B)$.
\end{result}
\begin{proof}
Choose an $L$-isogeny decomposition of $J$ into powers of $L$-simple
abelian varieties. Its rational inverse followed by one factor projection,
and then multiplication by a positive integer clearing denominators,
is an actual homomorphism $q:J\to B$ which is nonzero on $W$ for at least
one factor $B$. This map is surjective, since its image is a nonzero
abelian subvariety of the $L$-simple $B$. Denote its one-dimensional
image on $W$ by $W_B$.

Put $D=\operatorname{End}^0_L(B)$, a division algebra with centre $E$,
and write $[D:E]=d^2$. For a fixed embedding $L\hookrightarrow\mathbf C$,
the vector space $H_1(B(\mathbf C),\mathbf Q)$ is a left $D$-vector space
of an integer dimension $m\ge1$. Consequently
\[
 2\dim B=m d^2[E:\mathbf Q].
\]
The Betti--etale comparison preserves every endomorphism of $B$, so
$V_\ell(B)$ is, as a left $D\otimes\mathbf Q_\ell$-module, the direct
sum of $m$ copies of its regular module.

Faltings' semisimplicity and endomorphism theorem supply a $G_L$-equivariant
projector $e$ with image $W_B$ and identify it with an idempotent in
$D\otimes\mathbf Q_\ell$.
The exact original locators are \cite[\S5, S\"atze3--4,p.~360]{GenFaltings1983};
the corresponding Hom theorem is Korollar1,p.~361. In particular the
projector is allowed to have $\ell$-adic coefficients: it is not silently
asserted to descend to a rational idempotent of $D$.

For each $\lambda\mid\ell$, write
\[
 D\otimes_E E_\lambda\simeq M_{r_\lambda}(D_\lambda),\qquad
 [D_\lambda:E_\lambda]=\delta_\lambda^2,\qquad
 r_\lambda\delta_\lambda=d.
\]
An idempotent in this factor of rank $s_\lambda$ over the division
algebra $D_\lambda$ is conjugate to a diagonal matrix with
$s_\lambda$ entries equal to one. Its image on the left regular module
has $\mathbf Q_\ell$-dimension
\[
 s_\lambda r_\lambda\delta_\lambda^2[E_\lambda:\mathbf Q_\ell]
 =s_\lambda d\delta_\lambda[E_\lambda:\mathbf Q_\ell].
\]
It follows that the actual rank of $e$ on $V_\ell(B)$ is
\[
 1=m\sum_{\lambda\mid\ell}
 s_\lambda d\delta_\lambda[E_\lambda:\mathbf Q_\ell].
\]
Every summand is a nonnegative integer. There is exactly one nonzero
summand, and all its factors are one. In particular $m=d=1$,
$D=E$, $[E:\mathbf Q]=2\dim B$, and the active completion
$E_\lambda$ equals $\mathbf Q_\ell$. This proves the stated CM quotient
and the exact summand identification. No assertion that $\dim B=1$
has been used or obtained.

For completeness, $E$ is a CM field, not only a commutative algebra
of the required degree. Choose a polarization over $L$. Its Rosati
involution $\dagger$ preserves $E=\operatorname{End}^0_L(B)$.
The associated positive real inner product on $H_1(B(\mathbf C),\mathbf R)$
makes $a^\dagger$ the adjoint of $a$, so
\[
 \operatorname{Tr}_{E/\mathbf Q}(aa^\dagger)>0\qquad(a\ne0).
\]
Here the equality with the ordinary representation trace uses the
already proved rank-one $E$-module structure of rational homology.
The complex structure commutes with $E$ and hence, on the free
rank-one $E\otimes\mathbf R$-module, belongs to $E\otimes\mathbf R$.
Its square is $-1$, which is impossible in a real factor of that
algebra. Therefore $E$ has no real embeddings.

The Rosati involution is nontrivial: for an identity involution the
positive trace form $a\mapsto\operatorname{Tr}_{E/\mathbf Q}(a^2)$
would force every archimedean factor to be real. Indeed a complex
factor has a vector with negative square trace, and rational
coordinates are dense in the real scalar extension, contradicting
positivity. For the fixed field $E^+$ of the nontrivial involution,
the same argument applied to the restricted trace form shows that
$E^+$ is totally real. Since $[E:E^+]=2$ and $E$ is totally imaginary,
$E$ is a CM field, as claimed.

\paragraph{Actual curve and Prym maps.}
For a smooth proper curve $\overline Y/L$, choose a point after a finite
extension and use its Abel--Jacobi morphism
\[
 a:\overline Y\longrightarrow J_{\overline Y},\qquad
 y\longmapsto[y-y_0].
\]
The composite $b=q\circ a:\overline Y\to B$ is nonconstant whenever
$q$ is nonzero: the image of $a$ generates the Jacobian, so a constant
$q\circ a$ (necessarily zero at $y_0$) would imply $q=0$.
The induced homology map is exactly $V_\ell(q)$ under the Abel--Jacobi
isomorphism, and therefore detects the specified line.

In the A.22 finite-cover setting, let
$f:\overline Y\to E_0$ have degree $n$ and let
$P=\ker(f_*:J_{\overline Y}\to E_0)^0$.
The explicit homomorphism
\[
 v=n\,\mathrm{id}_{J_{\overline Y}}-f^*f_*
\]
has image in $P$, since $f_*f^*=[n]$ and its image is connected.
On a line $W\subset\ker V_\ell(f_*)$, it acts by multiplication by
$n$, so it is nonzero on $W$. Apply the preceding argument to the
line $V_\ell(v)(W)\subset V_\ell(P)$. It gives an actual quotient
$q_P:P\to B$ of CM type and an actual nonconstant morphism
\[
 \overline Y\xrightarrow{a}J_{\overline Y}
 \xrightarrow{v}P\xrightarrow{q_P}B
\]
which detects $W$. This construction proves the precise CM-factor
consequence for the residual weight-one branch; it does not merely
attach a representation-theoretic name to the line.

\paragraph{The remaining morphism obligations.}
For $Y=\overline Y\setminus T$, the displayed map $b$ restricts to
$Y\to B\setminus\{0\}$ exactly when $b^{-1}(0)\subset T$, by the
definition of inverse image of the open complement. That containment
has not been established for an arbitrary finite cover in A.22.
Moreover the rank calculation allows CM abelian varieties of dimension
greater than one; it does not produce one of the punctured elliptic
targets covered by F.3.15. F.3.16 supplies both the elliptic targets and
the necessary fibre containment for its particular Fermat cover by
explicit equations. An application to all A.22 covers still requires
the corresponding actual maps and a theorem applying to those targets.
\end{proof}

\begin{result}{F.3.19}{An absolutely simple CM surface in a finite-cover Prym}{Proved; the actual geometric maps are explicit}
For $U=E_0\setminus\{O\}$ with $E_0:y^2=x^3-x$, there is a degree-$800$ finite etale cover over $L=\mathbf Q(i,\zeta_5)$ whose smooth compactification is the genus-$171$ Fermat curve $C_{20}$. Its proper Prym contains an abelian surface isogenous to the absolutely simple CM Jacobian of $Y^2=1-X^5$. The inclusion is injective on integral Tate modules at every prime, including $2$ and $5$. The same open cover has the explicit map to the once-punctured genus-two curve described below. No procyclic lift with open arithmetic normalizer is asserted.
\end{result}
\begin{proof}
Work over $L=\mathbf Q(i,\zeta_5)$. Let
\[
 C_{20}:x^{20}+y^{20}=z^{20},\qquad
 a:C_{20}\longrightarrow C_4,\quad(x:y:z)\longmapsto(x^5:y^5:z^5).
\]
The target $C_4$ and the map $f_4=[4]q_x:C_4\to E_0$ are those of
F.3.16. The degree of $a$ is25; its deck group over $L$ is
$(\mu_5)^2$, and it ramifies only above the coordinate-zero points of
$C_4$. F.3.17 proves that all those points map under $q_x$ into $D[4]$.
Consequently
\[
 f_{20}=f_4\circ a:C_{20}\to E_0
\]
has degree~800 and is branched only over $O$. Thus
\[
 Y_{20}=C_{20}\setminus f_{20}^{-1}(O)\longrightarrow
 U=E_0\setminus\{O\}
\]
is an actual finite etale cover. The smooth plane curve $C_{20}$ has
genus~171.

Let $X_5$ be the smooth proper genus-two curve with affine equation
\[
 X_5:\quad Y^2=1-X^5,
\]
and let $\infty$ be its unique point at infinity. The formulas
\[
 q:C_{20}\longrightarrow X_5,\qquad
 X=(x/z)^4,\quad Y=(y/z)^{10}
\]
give a morphism of smooth proper curves. It has degree~40. Indeed on
function fields $L(C_{20})=L(u,v)$ with $u^{20}+v^{20}=1$, while
$L(X_5)=L(u^4,v^{10})$. The extension has degree at most40 by the
displayed fourth and tenth roots. The40 distinct automorphisms
$(u,v)\mapsto(\alpha u,\beta v)$ for
$\alpha\in\mu_4$, $\beta\in\mu_{10}$ fix this subfield; hence the
degree is exactly40.

\paragraph{CM and absolute simplicity, with the algebra embedding obstruction proved.}
Put $B=J_{X_5}$. The automorphism
\[
 \tau:(X,Y)\longmapsto(\zeta_5X,Y)
\]
fixes $\infty$ and induces an endomorphism of $B$. Its eigenvalues on
the basis $dX/Y,X\,dX/Y$ of regular differentials are
$\zeta_5,\zeta_5^2$; on the complex-conjugate part they are
$\zeta_5^{-1},\zeta_5^{-2}$. Thus its minimal polynomial on rational
first homology is $\Phi_5$. The faithful homology representation of
endomorphisms gives an actual embedding
\[
 E=\mathbf Q(\zeta_5)\hookrightarrow\operatorname{End}^0_L(B).
\]
This has degree~4, equal to $2\dim B$, so $B$ is of CM type over $L$.

We prove absolute simplicity without assuming an unproved description
of the entire endomorphism algebra. A nonsimple abelian surface over
$\overline L$ is isogenous to a product of two elliptic curves.
If the factors are nonisogenous, its endomorphism algebra is
$D_1\times D_2$, where each $D_i$ is $\mathbf Q$ or an imaginary
quadratic field. A unital embedding of the degree-four field $E$ into
this product would give unital embeddings into each factor, which is
impossible by their dimensions.

For two isogenous elliptic factors, the algebra is $M_2(D)$ with
$D=\mathbf Q$ or an imaginary quadratic field. When $D=\mathbf Q$,
every matrix satisfies its degree-two characteristic polynomial, so
the image of a primitive element of $E$ cannot have degree four.
When $D$ is imaginary quadratic, its central copy in $M_2(D)$ commutes
with $E$. The field $E$ is cyclic quartic and its only quadratic
subfield is the real field $\mathbf Q(\sqrt5)$: the Galois group
$(\mathbf Z/5)^\times$ is cyclic of order4 and complex conjugation
is its unique subgroup of order2. Hence $D$ is not contained in $E$,
and $E\cap D=\mathbf Q$. Since $D/\mathbf Q$ is quadratic,
$E\otimes_{\mathbf Q}D$ is a field of degree~8 over $\mathbf Q$.
Multiplication of the commuting copies defines a unital homomorphism
\[
 E\otimes_{\mathbf Q}D\longrightarrow M_2(D).
\]
Its source is a field, so this homomorphism is injective. Both sides
have dimension~8 over $\mathbf Q$, making it an isomorphism of
algebras. That is impossible: the source is commutative and the
target contains the explicitly noncommuting matrix units
$e_{12},e_{21}$, whose two products are $e_{11},e_{22}$.
This rules out every elliptic isogeny decomposition and proves that
$B$ is absolutely simple.

\paragraph{The actual Prym inclusion.}
The norm and pullback maps satisfy
\[
 q_*q^*=[40]:B\to B.
\]
Every homomorphism $B\to E_0$ is zero: a nonzero such homomorphism
would be surjective, and its connected kernel would be a
one-dimensional abelian subvariety of the absolutely simple surface
$B$. In particular the actual composite
\[
 B\xrightarrow{q^*}J_{C_{20}}\xrightarrow{f_{20,*}}E_0
\]
is zero. The image of $q^*$ is connected and therefore belongs to
$P=\ker(f_{20,*})^0$. The identity $q_*q^*=[40]$ shows that its kernel
is finite, so it is an abelian subvariety of $P$ isogenous to $B$.
For every prime $\ell$ it gives the exact injection
\[
 T_\ell B\xrightarrow{q^*}T_\ell P\subset T_\ell J_{C_{20}},
\]
including $\ell=2,5$: applying $q_*$ multiplies a kernel vector by40,
and Tate modules are torsion-free. Rationally, $\tfrac1{40}q_*$ is a
retraction onto this factor.

For any $\ell\equiv1\pmod5$, the algebra
$E\otimes\mathbf Q_\ell\simeq\mathbf Q_\ell^4$ acts on
$V_\ell B$ with four one-dimensional eigenspaces. Every endomorphism
of this $E$-action is defined over $L$, so all four lines are
$G_L$-stable. Their images in $V_\ell P$ are genuine weight-one
CM lines in this finite-cover Prym.

The map to the punctured genus-two curve is also defined on the same
open cover: $q^{-1}(\infty)$ consists precisely of the coordinate-zero
points $z=0$, and these have already been removed in $Y_{20}$.
Thus $q$ restricts to
\[
 Y_{20}\longrightarrow X_5\setminus\{\infty\}.
\]
The explicit target has genus two, and its Jacobian has just been
proved to have no elliptic isogeny factor. This identifies the precise
geometric target beyond the elliptic maps in F.3.16. The construction
provides no procyclic lift of these Tate-module lines with open
arithmetic normalizer. LM6 A.22 remains unresolved.
\end{proof}

\begin{result}{F.3.20}{Degree-six recognition for every once-punctured elliptic curve}{Proved; all number fields and all primes}
Let $E$ be an elliptic curve over a number field $K$, put
$U=E\setminus\{O\}$, and let $\ell$ be any prime. Set
\[
 P=\pi_1(U_{\overline K})^{(\ell)},\qquad
 \Pi=\pi_1(U)/\ker\bigl(\pi_1(U_{\overline K})\to P\bigr),
 \qquad p:\Pi\to G_K.
\]
Every nontrivial closed procyclic $A\subset P$ whose arithmetic
normalizer $N=N_\Pi(A)$ has open image under $p$ is contained in an
actual cusp inertia subgroup of $P$.
\end{result}
\begin{proof}
We work in the maximal pro-$\ell$ group and keep the arithmetic normalizer $N$ throughout.

\paragraph{1. The exact primary input.}

Nakamura~1996, Lemma~4.5 and Corollary~4.6,p.~116, give the surjection
from the graded arithmetic Galois image of a fixed affine hyperbolic
curve to its universal arithmetic quotient. Equations(4.7.1)--(4.7.2),
p.~117, identify this quotient and its action\cite[pp.~115--117]{AuditNakamura1996}: in degree $m$ it is acted
on by the similitude character to the power $m/2$. For $(g,n)=(1,1)$,
degree~6 is nonzero by Corollary~4.2 and the degree~6 entry of(4.3),p.~115.
These are unconditional statements in the original article; the
standing prime in Section~2,p.~101 is any rational prime\cite[p.~101]{AuditNakamura1996}. The
$\operatorname{End}(J)=\mathbf Z$ condition in Theorem B,p.~100,
concerns its additional centralizer-finiteness conclusion and is not
a premise of Lemma~4.5 or Corollary~4.6.

In explicit terms, put
\[
 G_K(m)=\{\sigma\in G_K:\rho_U(\sigma)\in\Gamma_{1,1}(m)\}.
\]
The degree-six image $\mathfrak a_6\subset
\operatorname{gr}_6\Gamma_{1,1}\otimes\mathbf Q_\ell$ has a nonzero
$G_K$-equivariant quotient
\[
 q_6:\mathfrak a_6\twoheadrightarrow Q_6,
 \qquad G_K\text{ acts on }Q_6\text{ by }\chi_\ell^3.
\]
Here $\mathfrak a_6$ is the rational span of the image of the actual
group $G_K(6)$ in $\Gamma_{1,1}(6)/\Gamma_{1,1}(7)$.
This is the quotient used below; no general normalizer assertion is
attributed to Theorem B.

\paragraph{2. The first Magnus degree.}
The finite-quotient detector of F.3.15, Steps~1--2 applies to every
once-punctured elliptic curve. For clarity its precise mechanism is
as follows. A generator $u$ of $A\simeq\mathbf Z_\ell$ survives a
finite $\ell$-group quotient. The inverse image of its cyclic image
is an open subgroup whose torsion-free abelianization detects $u$.
It is the maximal pro-$\ell$ group of an actual finite etale cover;
normal closures of the relevant iterated $\ell$-covers prove that
identification. Descending the cover and intersecting its arithmetic
subgroup with $N$ gives the line of $u$ in actual curve homology,
without choosing a section of $N$. Its Frobenius weight is one or
two, by the proper-homology and boundary exact sequence.

The injective Magnus embedding and its rational logarithm give a
first nonzero term $u_d$ in the degree$d$ free Lie algebra on
$V=V_\ell(E)$. Geometric conjugation is trivial on that term, so the
normalizer exponent $\theta$ descends to $p(N)$ and satisfies
$\sigma(u_d)=\theta(\sigma)u_d$. The degree$d$ term has Frobenius
weight$d$. Comparing the same character with the detector shows
$d=1$ or2. These weight statements use arithmetic Frobenius on
homology, so $V_\ell(E)$ has eigenvalue absolute value$q^{1/2}$.

\paragraph{3. The only degree-one case needing an obstruction.}
Extend $K$ finitely so that all geometric endomorphisms of $E$ are
defined. Openness of the arithmetic normalizer image is preserved.
Write $H=p(N)$ after this replacement. Faltings' semisimplicity and
endomorphism theorem, original \S5,S\"atze3--4,p.~360, give\cite{GenFaltings1983}
\[
 \operatorname{End}_{G_L}(V_\ell E)
 =\operatorname{End}^0_L(E)\otimes\mathbf Q_\ell
\]
for every further finite extension $L/K$.
For a non-CM curve this algebra is $\mathbf Q_\ell$, which has no
rank-one projector on the two-dimensional $V$. For a CM curve at a
prime not split in its CM field, the algebra is a quadratic field,
which also has no nontrivial idempotent. Semisimplicity therefore
excludes a one-dimensional subrepresentation of any open arithmetic
subgroup in those cases.

The remaining case is split CM, where
\[
 V=\mathbf Q_\ell x\oplus\mathbf Q_\ell y,
 \qquad \rho_V=\chi_1\oplus\chi_2,
 \qquad \chi_1\chi_2=\chi_\ell.
\]
These are the only stable lines. The ratio $\chi_1/\chi_2$ has
infinite image on $H$. Otherwise a finite extension would make the
two characters equal, and the displayed endomorphism algebra would
become $M_2(\mathbf Q_\ell)$ instead of the split quadratic algebra,
contradicting Faltings' theorem and the full geometric CM
endomorphism algebra. Since $\mathbf Q_\ell$ has only finitely many
roots of unity, choose an actual $\tau\in H$ for which
$\chi_1(\tau)/\chi_2(\tau)$ is not a root of unity.

\paragraph{4. The character projection produces an actual arithmetic element.}
Use Nakamura~1995,(4.3)--(4.4),p.~214, to lift\cite{AuditNakamura1995} the degree-at-least-three
outer action canonically to a boundary-fixing representative.
The isomorphism of graded pieces there shows that an element of
$G_K(6)$ has representative in $\Gamma^*_{1,1}(6)$.
Its leading derivation $D$ has degree~6 and is special:
\[
 D([x,y])=0.
\]
Let $\mathcal D_6$ denote this finite-dimensional space over
$\mathbf Q_\ell$; the integral action gives it a fixed
$\mathbf Z_\ell$-lattice containing the actual image. Set
\[
 M=\operatorname{im}\bigl(H\cap G_K(6)\longrightarrow
       \operatorname{gr}_6\Gamma^*_{1,1}\bigr).
\]
This is a compact, hence closed, additive subgroup of a finite free
$\mathbf Z_\ell$-module. It is a $\mathbf Z_\ell$-module: multiplication
by an $\ell$-adic integer is the limit of integer multiples, and
closedness keeps that limit in $M$. It is invariant under
$T=\operatorname{Ad}(\tau)$, since $G_K(6)$ is normal and
$\tau\in H$. Moreover $M\otimes\mathbf Q_\ell=\mathfrak a_6$;
restriction to the finite-index subgroup $H\cap G_K(6)$ does not
change the rational image.

The eigencharacters of $T$ on $\mathcal D_6\subset
\operatorname{Hom}(V,L_7(V))$ have the form
$\chi_1^a\chi_2^b$ with $a+b=6$. Because the chosen ratio is not
a root of unity, the eigenvalue
\[
 c_0=\chi_\ell(\tau)^3
\]
occurs precisely for the middle weight $(a,b)=(3,3)$.
For the finite set of distinct eigenvalues $c$ of $T$, the polynomial
\[
 e(T)=\prod_{c\ne c_0}\frac{T-c}{c_0-c}
\]
is the projection to this middle weight. Its coefficients belong
to $\mathbf Q_\ell$. Choose an integer $r\ge0$ such that all
coefficients of $\ell^r e(T)$ are in $\mathbf Z_\ell$.

By Step~1 choose $m\in M$ with $q_6(m)\ne0$. Equivariance gives
$q_6(e(T)m)=q_6(m)$, since $T$ acts on the quotient by $c_0$.
Then
\[
 m_0=\ell^r e(T)m\in M,\qquad q_6(m_0)\ne0.
\]
Membership in $M$ follows from its $\mathbf Z_\ell$-module property
and invariance under all powers of $T$. By the definition of $M$,
there is an actual
\[
 \sigma\in H\cap G_K(6)
\]
whose degree-six derivation is $m_0$. Thus its first nonzero action
is the nonzero special derivation $D=m_0$ in middle weight $(3,3)$.
The use of closed image and cleared denominators is essential here:
a rational projected vector alone would not be a specified Galois
element.

\paragraph{5. The complete middle-weight calculation.}
Write $b_w$ for the standard Lyndon bracket of the word $w$ in the
ordered alphabet $x<y$. The degree-seven bidegree$(4,3)$ words are
\[
 (xxxxyyy,\ xxxyxyy,\ xxxyyxy,\ xxyxxyy,\ xxyxyxy),
\]
and the bidegree$(3,4)$ words are
\[
 (xxxyyyy,\ xxyxyyy,\ xxyyxyy,\ xxyyyxy,\ xyxyxyy).
\]
In these two ordered bases respectively, define
\[
\begin{array}{c|rr}
 &D_1&D_2\\\hline
 D(x)&(1,-5,-6,1,-2)&(-3,-9,-6,-3,-2)\\
 D(y)&(-1,4,6,8,0)&(3,12,6,0,8).
\end{array}
\]
Both derivations are special by direct expansion. They span the
middle-weight special space: the two source spaces have dimension~5
each, the target $L_8(4,4)$ has dimension~8, and
\[
 (a,b)\longmapsto[a,y]+[x,b]
\]
is surjective. The latter follows from
$L_8=[x,L_7]+[y,L_7]$, itself the Jacobi-identity consequence of
generation by $x,y$, after taking bidegree$(4,4)$.

The degree-six middle Lyndon words are
$(xxxyyy,xxyxyy,xxyyxy)$. Appending their three brackets with $x$
to the columns $D_1(x),D_2(x)$ gives the matrix
\[
 \begin{pmatrix}
 1&-3&1&0&0\\-5&-9&0&1&0\\-6&-6&0&0&1\\
 1&-3&0&0&0\\-2&-2&0&0&0
 \end{pmatrix},\qquad\det=-8.
\]
For $y$ the corresponding matrix is
\[
 \begin{pmatrix}
 -1&3&-1&0&0\\4&12&-1&-1&0\\6&6&0&0&-1\\
 8&0&1&0&-1\\0&8&0&-1&1
 \end{pmatrix},\qquad\det=-8.
\]
Therefore every nonzero middle-weight special $D$ satisfies
\[
 D(x)\notin[x,L_6],\qquad D(y)\notin[y,L_6].
\]
The determinants are nonzero over $\mathbf Q_\ell$ also at
$\ell=2$. No assertion of invertibility over $\mathbf Z_2$ is used.

\paragraph{6. The actual normalizer conjugacy equation through degree seven.}
Suppose $d=1$ and $\log u=ax+$ higher terms, $a\ne0$; the other
line is identical using the second matrix. Lift the $\sigma$ from
Step~4 to $n\in N$. Since $\sigma$ is the identity on $V$, its
normalizer exponent is1. Its actual action centralizes $u$.
The canonical boundary-fixing representative differs by an inner
automorphism, so in the rational completed free Lie algebra
\[
 \widetilde\sigma(\log u)
 =\exp(\operatorname{ad}W)(\log u).
\]
Comparison in degree~2 gives $W_1=bx$. Replacing the conjugating
element $\exp W$ by
$\exp W\exp(-(b/a)\log u)$ preserves its conjugation action and
removes its degree-one term. Comparison in degrees3,4,5,6 then
successively gives $W_2=W_3=W_4=W_5=0$, using the exact centralizer
$C_L(x)=\mathbf Q_\ell x$. This holds without primitivity of the
original abelian image.

In degree~7 the resulting equality is
\[
 aD(x)=[W_6,ax],
\]
contradicting Step~5. The second matrix treats $\log u=ay+$ higher
terms. This excludes degree one in every remaining split-CM case.

\paragraph{7. Degree two concludes in the actual pro-$\ell$ group.}
Only $d=2$ remains. The one-dimensional space $L_2(V)$ is
$\det V=\mathbf Q_\ell(1)$. Canceling the nonzero first term of
$\log u$ identifies the actual normalizer character with
$\chi_\ell\circ p$. Nakamura~1994,Theorem~2.1.1,p.~94, now applies\cite{AuditNakamura1994}
to the class of finite $\ell$-groups and places $A$ in actual cusp
inertia. The finite extension used to define endomorphisms leaves
this geometric cusp subgroup unchanged.

\paragraph{Finite certificate.}
The independent exact program \texttt{check\_elliptic\_degree\_six.py} constructs the Lyndon brackets from associative words, checks that both displayed derivations preserve $[x,y]$, computes the full middle-space dimension, and obtains both displayed matrices with determinant $-8$. Its arithmetic input is the cited original theorem; the closed-image argument above supplies the actual Galois element. Root independently ran the certificate and checked the source and proof.

\paragraph{Scope.}
The theorem extends the family of punctured elliptic targets available
in F.3.16 and uses an opened original degree-six source. It does not
construct a detecting elliptic morphism for an arbitrary finite
cover in A.22, and does not treat the higher-dimensional CM quotient
constructed in F.3.19. No conclusion about the full profinite
common-kernel subgroup is drawn.
\end{proof}

\begin{result}{F.3.21}{Actual genus-two normalizer recognition at nonsplit primes}{Proved; exact finite-cover and cyclotomic maps}
Let $X_5:Y^2=1-X^5$ be its smooth proper genus-two model over a
number field containing $\mathbf Q(\zeta_5)$, let $O=\infty$ be its
unique point at infinity, and put $U_5=X_5\setminus\{O\}$.
For every prime $\ell=5$ or $\ell\not\equiv1\pmod5$, a nontrivial
closed procyclic subgroup of
\[
 P=\pi_1(U_{5,\overline K})^{(\ell)}
\]
with open arithmetic normalizer image is contained in an actual
cusp inertia subgroup. The assertion remains valid over every
finite number-field extension. F.3.22 below supplies the remaining split primes
$\ell\equiv1\pmod5$ for this same curve. The full profinite
A.22 requirement remains unresolved.
\end{result}
\begin{proof}
We use the actual detecting cover to identify the normalizer character before applying cusp recognition.

\paragraph{1. The exact absence of a rational line at these primes.}

By F.3.19 the Jacobian $B=J_{X_5}$ is absolutely simple and has
an embedding $E=\mathbf Q(\zeta_5)\subset\operatorname{End}^0_K(B)$.
The full geometric endomorphism algebra is a division algebra
$D$, and $H_1(B(\mathbf C),\mathbf Q)$ is a faithful finite-dimensional
left $D$-vector space. Thus $[D:\mathbf Q]$ divides~4. Since it
already contains the degree-four field $E$, one has
\[
 \operatorname{End}^0_{\overline K}(B)
 =\operatorname{End}^0_L(B)=E
\]
for every finite extension $L/K$.
Faltings' semisimplicity and endomorphism theorem, original
\S5, S\"atze~3--4, p.~360, identify\cite{GenFaltings1983}
\[
 \operatorname{End}_{G_L}(V_\ell B)=E\otimes\mathbf Q_\ell.
\]
The module $V_\ell B$ is free of rank one over this algebra, by
the endomorphism-compatible Betti--\'etale comparison and the
degree-four action on four-dimensional rational homology.

For $\ell\ne5$, each factor of $E\otimes\mathbf Q_\ell$ has degree
$f=\operatorname{ord}_5(\ell)$: the primitive fifth roots have
Frobenius orbits of precisely that size over $\mathbf F_\ell$,
and the separable factors lift by Hensel's lemma. For $\ell=5$,
\[
 \Phi_5(T+1)=T^4+5T^3+10T^2+10T+5
\]
is Eisenstein, so the local algebra is a field of degree~4.
At the primes in the theorem every field factor therefore has
degree greater than one. A nonzero idempotent on this rank-one
module has rank equal to a nonempty sum of those degrees, and
cannot have rank one. Semisimplicity would supply exactly such
an idempotent for any one-dimensional $G_L$-stable subspace.
There is consequently no rational Galois line even after an
arbitrary finite extension.

\paragraph{2. An actual finite cover detects the cyclic generator.}
Let
\[
 \Pi=\pi_1(U_5)/\ker(\pi_1(U_{5,\overline K})\to P),
 \quad p:\Pi\to G_K,\quad N=N_\Pi(A),\quad H=p(N).
\]
The open subgroup $H$ is also closed because $N$ is compact.
Since $P$ is a free pro-$\ell$ group of rank~4, it is torsion-free
and $A\simeq\mathbf Z_\ell$. Choose a generator $u$ and a finite
$\ell$-group quotient $P\to Q$ in which its image $\bar u$ is
nonidentity. Put
\[
 P_0=\operatorname{preimage}(\langle\bar u\rangle)\subset P.
\]
The map $P_0\to\langle\bar u\rangle\to\mathbf Z/\ell$ takes
$u$ to a nonzero element. Its image in the finite-free
$\mathbf Z_\ell$-module $P_0^{\rm ab}$ is therefore nonzero.

The inverse image of $P_0$ in $\pi_1(U_{5,\overline K})$ is the
geometric covering subgroup of a connected finite \'etale curve
$Y_{\overline K}\to U_{5,\overline K}$. All equations and maps
descend to some finite extension $L/K$, and their finite-\'etale
property descends after enlarging $L$; the chosen geometric
connectedness is preserved. Its geometric maximal pro-$\ell$
group is exactly $P_0$. Indeed the Galois closure of this cover
has an $\ell$-group as deck group; the normal closure of a further
finite $\ell$-cover is obtained from the finitely many conjugate
$\ell$-covers over that Galois closure and again has an
$\ell$-group. This proves that every finite $\ell$-quotient of
$\pi_1(Y_{\overline K})$ factors through $P_0$, and proves the
maximal-pro-$\ell$ identification, not just a map onto it.

The arithmetic covering subgroup descends to an open subgroup
$\Pi_Y\subset\Pi$ with geometric kernel $P_0$ and projection
$G_L$. It does so because the kernel already divided out in
$\Pi$ is contained in its geometric covering subgroup. Set
$N_0=N\cap\Pi_Y$. Its index in $N$ is finite, hence
$p(N_0)$ is open and corresponds to a finite extension $L_0/L$.
The geometric kernel of $N_0\to p(N_0)$ lies in $P_0$ and acts
inner on $P_0$, so it acts trivially on $P_0^{\rm ab}$.
The nonzero homology line of $u$ is therefore an ordinary
$G_{L_0}$-stable line in the actual curve homology
\[
 P_0^{\rm ab}\otimes\mathbf Q_\ell
 =H_1(Y_{\overline K},\mathbf Q_\ell).
\]
No splitting of $N\to H$ or of $N_0\to G_{L_0}$ is used.

\paragraph{3. The same character has a first Magnus weight.}
The injective Magnus embedding of the rank-four free pro-$\ell$
group has a rational logarithm with a first nonzero homogeneous
term $u_d\in L_d(V_\ell B)$. This uses that $U_5$ has only one
geometric puncture, so its first homology is the proper homology
of $X_5$ and is pure of weight one. Every $L_d(V_\ell B)$ is
pure of weight $d$.

Conjugation by $n\in N$ acts as $u\mapsto u^{\theta(n)}$.
Geometric conjugation acts trivially on the first nonzero
homogeneous term, so comparison there shows that $\theta=1$
on $N\cap P$. Thus $\theta$ factors through $H$, and
\[
 \rho(p(n))u_d=\theta(n)u_d.
\]
This is the same character as on the detected homology line
of Step~2. The assertion is an equality of actions on a
nonzero vector, with a torsion-free coefficient field, so
there is no finite-character ambiguity at this stage.

\paragraph{4. Weight two forces the actual boundary character.}
Let $\overline Y$ be the smooth proper compactification and
$T=\overline Y\setminus Y$. The homology localization sequence is
\[
 0\longrightarrow W_T\longrightarrow
 H_1(Y_{\overline K},\mathbf Q_\ell)
 \longrightarrow H_1(\overline Y_{\overline K},\mathbf Q_\ell)
 \longrightarrow0,
\]
where
\[
 W_T\simeq\left(\bigoplus_{t\in T(\overline K)}
                    \mathbf Q_\ell(1)e_t\right)
          \big/\mathbf Q_\ell(1)\Bigl(\sum_t e_t\Bigr).
\]
Take a good arithmetic Frobenius for both curves outside the
finite set of excluded places. A nonzero image of the detected
line in proper homology would have eigenvalue absolute value
$q^{1/2}$. Comparison with its first Magnus term, whose
absolute value is $q^{d/2}$, would force $d=1$. But then
$u_1$ would be a one-dimensional Galois-stable subspace of
$V_\ell B$ on an open subgroup, excluded by Step~1.
Therefore the detected line has zero image in proper homology
and lies in $W_T$.

Choose a finite extension $L_1/L_0$ over which every point of
$T$ is rational. The displayed permutation representation
then becomes a direct sum and quotient of copies of the
same character $\chi_\ell$. Thus $G_{L_1}$ acts on all of
$W_T$ exactly by $\chi_\ell$, including on the nonzero
detected line. This proves the actual equality
$\theta=\chi_\ell$ on the inverse image of $G_{L_1}$ in $N_0$.
It also shows that the first Magnus degree was~2. No
semisimple splitting of the mixed homology extension is
needed: the line was first placed in its actual kernel.

\paragraph{5. Passage back to the full arithmetic normalizer.}
The inverse image $N_1\subset N_0$ of $G_{L_1}$ surjects
onto $G_{L_1}$. Every element of the full normalizer in the
arithmetic group over $L_1$, namely
$N\cap p^{-1}(G_{L_1})$, has the same arithmetic image as
some element of $N_1$. Their quotient lies in $N\cap P$,
where Step~3 proved that $\theta$ is trivial. Consequently
the actual full normalizer over $L_1$ acts on $A$ by
$\chi_\ell$ and has open (indeed surjective) arithmetic image.
Nakamura~1994, Theorem~2.1.1, p.~94, for the full class of finite\cite{AuditNakamura1994}
$\ell$-groups now gives an actual cusp inertia subgroup
containing $A$. The finite field extensions leave the
geometric fundamental group and its unique-cusp conjugacy
class unchanged. This proves the theorem.

\paragraph{Exact remaining scope.}
At a split prime $\ell\equiv1\pmod5$, the four rank-one
CM eigenspaces are actual Galois lines, so Step~1 supplies
no exclusion of the degree-one case. The theorem does not
assert that one of these lines lifts to a procyclic
arithmetic-normalizer subgroup. It specifies exactly the
prime range treated by the boundary argument alone. F.3.22
below excludes degree one at the remaining split primes
through simultaneous arithmetic conjugacy constraints.
\end{proof}

\begin{result}{F.3.22}{All-prime recognition on the genus-two CM target}{Proved; actual simultaneous normalizer constraints}
Let $X_5$ be the smooth proper curve $y^2=1-x^5$, let $O$ be its
unique point at infinity, and let $U_5=X_5\setminus\{O\}$.
Over every number field $K$, and for every rational prime $\ell$,
every nontrivial closed procyclic subgroup
\[
 A\subset P=\pi_1(U_{5,\overline K})^{(\ell)}
\]
whose arithmetic normalizer has open image in $G_K$ is contained
in actual cusp inertia. The result concerns this specified
genus-two curve; arbitrary-cover and full profinite A.22 remain
unresolved.
\end{result}
\begin{proof}
We impose the four conjugacy constraints on one actual arithmetic image and compute its entire relevant leading-term space.

\paragraph{1. Reduction to the four split CM lines.}

F.3.21 proves the assertion for $\ell=5$ or
$\ell\not\equiv1\pmod5$. We now take $\ell\equiv1\pmod5$.
In particular $\ell\ne2,5$. Write $E=\mathbf Q(\zeta_5)$,
$B=J_{X_5}$, and $V=V_\ell B$. F.3.19 proves absolute
simplicity, and the dimension argument of F.3.21 gives
\[
 \operatorname{End}^0_{\overline{\mathbf Q}}B=E.
\]
All these endomorphisms are defined over $E$. After extending
$K$ finitely, we may suppose $E\subset K$; the normalizer's
image stays open, and the geometric group $P$ is unchanged.
Fix an embedding $\iota:E\hookrightarrow\mathbf Q_\ell$.
The automorphism
\[
 t:(x,y)\longmapsto(\zeta_5x,y)
\]
splits $V$ into four lines $V_a$, $a=1,2,3,4$, on which $t$
has eigenvalue $\iota(\zeta_5)^a$. Their $G_E$-characters are
denoted $\chi_a$.

By Faltings\cite[\S5, S\"atze~3--4, p.~360]{GenFaltings1983}, the endomorphism algebra on every open arithmetic
subgroup is $E\otimes\mathbf Q_\ell=\mathbf Q_\ell^4$.
Semisimplicity then implies that its one-dimensional invariant
subspaces are precisely these four lines. Indeed a stable
line has an equivariant projector of rank one; the only
rank-one idempotents of the displayed rank-one module are
the four primitive coordinate projectors. This also excludes
coincidence of two characters on an open subgroup.

Let $u$ generate $A$. The actual finite-cover detection and
weight comparison of F.3.21 apply before its nonsplit-prime
exclusion. Thus either the first nonzero Magnus term has
degree one, or the detected cover-homology line lies in its
boundary kernel. In the latter case the exact boundary
argument and full-normalizer character transfer of F.3.21
give $\theta=\chi_\ell$ on an open subgroup and prove the
theorem. It remains to rule out degree one. Its nonzero
linear term lies on one of the four $V_a$.

\paragraph{2. All character relations, proved through local Hodge--Tate maps.}
The principal polarization is preserved by $t$. The Weil
pairing therefore pairs $V_a$ only with $V_{5-a}$, and gives
the exact character equalities
\[
 \chi_1\chi_4=\chi_2\chi_3=\chi_\ell.
 \tag{X5.1}
\]
There are no further finite-image monomial relations on any
open subgroup of $G_E$.

Here is a direct proof of the latter assertion, including
the local realization used. At every prime above $\ell$,
$X_5$ has good reduction. One smooth chart is
$y^2=1-x^5$; its chart at infinity is
$v^2=s^6-s$, where $s=1/x$ and $v=y/x^3$.
Their derivatives and $\ell\ne2,5$ give smoothness; they
glue to the proper double cover of $\mathbf P^1$ branched
at the five simple roots and infinity. Its Jacobian is an
abelian scheme. Tate's functorial Hodge--Tate decomposition
for its $\ell$-divisible group, as checked in B.1.4 from
Tate~1967, \S4.1, Theorem~3 and Corollary~2, identifies\cite{Tate1967} the
exponent-one summand of $V\otimes C_\ell$ with
$\operatorname{Lie}B\otimes C_\ell(1)$.
The forms $dx/y,x\,dx/y$ show that $t$ on the tangent
space has eigenvalues $\zeta_5,\zeta_5^2$.

The four embeddings of $E$ into $\mathbf Q_\ell$ give
exponent-one index sets
\[
 \{1,2\},\quad \{2,4\},\quad \{4,3\},\quad \{3,1\}.
 \tag{X5.2}
\]
Every finite extension of $E$ has a place over each of these
four places; restriction retains these exponents. A
finite-image character has exponent zero after its finite
trivializing extension, hence already has exponent zero.
For a finite-image product $\prod_a\chi_a^{n_a}$ the four
rows of (X5.2) give
\[
 n_1+n_2=n_2+n_4=n_4+n_3=n_3+n_1=0.
\]
Solving gives
$(n_1,n_2,n_3,n_4)=(b,-b,-b,b)$. By (X5.1) that
product is identically one. This proves the assertion on
every open subgroup, without a CM-image openness theorem.

Choose a symplectic basis
\[
 x_1\in V_1,\quad y_1\in V_4,\quad
 x_2\in V_2,\quad y_2\in V_3,
 \qquad \omega=[x_1,y_1]+[x_2,y_2].
\]
The characters of the diagonal action are
$(\alpha_1,\chi_\ell/\alpha_1,
  \alpha_2,\chi_\ell/\alpha_2)$, with
$\alpha_1=\chi_1,\alpha_2=\chi_2$.
It follows from the proved relation lattice that the
$\chi_\ell^3$-weight part of the degree-six free Lie algebra
is exactly
\[
 \bigoplus_{j=0}^3 L_6(j,j,3-j,3-j).
 \tag{X5.3}
\]
Here the entries are the multiplicities of
$x_1,y_1,x_2,y_2$. The assertion remains true on every open
arithmetic subgroup. Every other character appearing in the
finite-dimensional degree-six derivation space has quotient
by $\chi_\ell^3$ of infinite image.

\paragraph{3. Four actual conjugate subgroups and one common arithmetic group.}
Because $U_5$ is defined over $\mathbf Q$, work in the
actual arithmetic quotient
\[
 1\longrightarrow P\longrightarrow\Pi_{\mathbf Q}
   \xrightarrow{p}G_{\mathbf Q}\longrightarrow1.
\]
The group over $K$ is $p^{-1}(G_K)$, after compatible
geometric base points. Let $N$ be the normalizer of $A$
there, and $H=p(N)$, which is open in $G_{\mathbf Q}$.
The first Magnus line of $A$ is some $V_a$.

For $b\in(\mathbf Z/5)^\times$, choose
$\gamma_b\in G_{\mathbf Q}$ with
$\gamma_b(\zeta_5)=\zeta_5^b$, and choose an actual lift
$\widetilde\gamma_b\in\Pi_{\mathbf Q}$, taking
$\gamma_1=\widetilde\gamma_1=1$. Set
\[
 A_b=\widetilde\gamma_b A\widetilde\gamma_b^{-1},\qquad
 N_b=\widetilde\gamma_b N\widetilde\gamma_b^{-1},\qquad
 H_b=\gamma_bH\gamma_b^{-1}.
 \tag{X5.4}
\]
Each $N_b$ normalizes $A_b$ and maps onto the open subgroup
$H_b$. The formula
$\gamma_b t\gamma_b^{-1}=t^b$ gives
\[
 \gamma_b(V_a)=V_{ab^{-1}}.
\]
The four first Magnus lines of $A_b$ are therefore exactly
the four distinct CM lines. Changing a chosen arithmetic
lift by an element of $P$ changes only geometric conjugacy,
which does not change its first linear term.

Choose a finite Galois extension $L/\mathbf Q$ such that
\[
 G_L\subset G_E\cap\bigcap_b H_b.
 \tag{X5.5}
\]
Such an $L$ exists by taking the normal core of that open
intersection. In particular $G_L$ is contained in each
$G_{\gamma_b(K)}$, since $H_b$ is; all four arithmetic
covering groups can now be restricted to the same field.
For each $b$, the closed subgroup
$N_b\cap p^{-1}(G_L)$ maps \emph{onto} $G_L$.
No simultaneous section is chosen or needed.

\paragraph{4. A special leading representative of the same outer action.}
Let $G_{\mathbf Q}(6)$ be the inverse image of the sixth
weight-filtration subgroup under the outer action for
$U_5$. For any
$\sigma\in G_L\cap G_{\mathbf Q}(6)$, take an automorphism
representative trivial through degree six in the rational
Malcev Lie algebra of $P$. Its action on the boundary loop
$z$ is conjugation, since its action on $V$ and hence on
$\chi_\ell$ is trivial.

One may make this representative fix $z$ while retaining
that filtration. Here are the details needed in genus two.
Write $\log z=\omega+$ higher terms and let $W$ be the
logarithm of a conjugator between its two images. Comparison
in degrees three and four first gives
$[W_1,\omega]=[W_2,\omega]=0$. The centralizer
of $\omega$ in the homogeneous free Lie algebra is
$\mathbf Q_\ell\omega$: a nonzero homogeneous element has
a least Lyndon word; commutation with the least word
$x_1y_1$ of $\omega$ forces these two words to commute.
Two commuting words are powers of a common word, whereas
a Lyndon word is not a proper power. Thus the degree must
be two; subtracting its multiple of $\omega$ then rules
out any remaining leading word. Consequently $W_1=0$,
and the only possible lower term is $W_2=c\omega$.
Multiplying the conjugator on the right by
$\exp(-c\log z)$ leaves its conjugation of $z$ unchanged
and removes that term. For the adjusted conjugator,
comparison in degrees five, six and seven successively
gives $W_3=W_4=W_5=0$. The required adjustment starts in degree
six.

The resulting degree-six derivation $D_\sigma$ is special:
$D_\sigma(\omega)=0$. It is uniquely determined by the
outer leading class, since two such representatives differ
by $\operatorname{ad}w$ with $w\in L_6$ and
$[w,\omega]=0$, forcing $w=0$. The construction is
$\operatorname{GSp}(V)$-equivariant by this uniqueness.
It is only a rational leading representative; the Galois
element $\sigma$ itself remains an actual element of $G_L$.

For each of the four subgroups in (X5.4), choose a lift of
this \emph{same} $\sigma$ in its restricted $N_b$. Its
conjugation exponent on a generator $u_b$ is one, because
the first linear term of $u_b$ is nonzero and $\sigma$ acts
trivially on $V$. Thus this lift centralizes $u_b$.
The boundary-fixing representative of $\sigma$ differs
from that lift by an inner automorphism. If
$\log u_b=a_bv_b+$ higher terms, the conjugacy equation is
\[
 \widetilde\sigma(\log u_b)
  =\exp(\operatorname{ad}W_b)(\log u_b).
\]
Its degree-two part gives $(W_b)_1=c_bv_b$. Right
multiplication of the conjugator by
$\exp(-(c_b/a_b)\log u_b)$ removes this first term without
changing its conjugation action. The free-Lie centralizer
of the generator $v_b$ then successively gives
$(W_b)_2=\cdots=(W_b)_5=0$. In degree seven one obtains
\[
 D_\sigma(v_b)\in[v_b,L_6].
 \tag{X5.6}
\]
The four separate lifts of the one $\sigma$ have therefore
given four simultaneous constraints on its unique special
derivation $D_\sigma$. Nonprimitive linear coefficients
$a_b$ cause no problem, since this conjugator adjustment
is made over $\mathbf Q_\ell$.

\paragraph{5. Projection inside the actual arithmetic image.}
Nakamura~1996, Corollary~4.2 and table~4.3, p.~115,
Lemma~4.5, p.~116, and (4.7.1)--(4.7.2), p.~117, supply\cite{AuditNakamura1996}
the nonzero arithmetic degree-six quotient
\[
 q_6:\operatorname{gr}_6\Gamma^{\mathrm{univ}}_{2,1}
       \longrightarrow \mathcal G_{2,1}^{,6},
\]
on which the similitude acts by $\chi_\ell^3$.
The actual arithmetic image for $U_5/L$ maps onto this
quotient rationally. The source expressly retains this
graded image under finite base extension, so the field
chosen in (X5.5) does not remove the nonzero quotient.

Let $M$ be the image of $G_L\cap G_{\mathbf Q}(6)$ in the
degree-six graded space, embedded into special derivations
by Step~4. It is a compact additive group in a finite
dimensional $\mathbf Q_\ell$-space, hence a closed
$\mathbf Z_\ell$-module contained in a finite lattice.
It is invariant under conjugation by $G_L$ and every
member satisfies all four equations (X5.6).
Choose $m\in M$ with $q_6(m)\ne0$.

The action on this derivation space is diagonal in the
four CM characters. By Step~2 there exists
$\tau\in G_L$ for which $\chi_\ell(\tau)^3$ differs from
every other character value occurring in that finite
space. To justify existence, the kernel of any of the
finitely many nontrivial quotient characters has empty
interior: an open kernel would give finite image, which
Step~2 excludes. These kernels are closed, and the Baire
property of the compact group $G_L$ excludes their finite
union.

For $T=\operatorname{Ad}(\tau)$ use the polynomial
spectral projector $e(T)$ onto eigenvalue
$\chi_\ell(\tau)^3$. After multiplying by a power of
$\ell$, its coefficients are in $\mathbf Z_\ell$; hence
\[
 m_0=\ell^r e(T)m\in M,\qquad
 q_6(m_0)=\ell^r q_6(m)\ne0.
 \tag{X5.7}
\]
The defining image property of $M$ lifts $m_0$ to an actual
$\sigma_0\in G_L\cap G_{\mathbf Q}(6)$. Its leading
derivation is special, has character $\chi_\ell^3$, and
still satisfies all four tangency constraints. The same
$q_6$ is used throughout; it is the arithmetic projection
of the moduli problem for this curve, independent of the
four chosen normalizer lifts.

\paragraph{6. The complete finite Lie-algebra calculation.}
Put $a=x_1,b=y_1,c=x_2,d=y_2$ and order the alphabet
$a<b<c<d$. Let $b_w$ be the standard Lyndon bracket.
Tangency means
\[
 D(v)=[A_v,v]\quad(v=a,b,c,d),\qquad A_v\in L_6.
\]
The map $A_v\mapsto[A_v,v]$ is injective in degree six.
Equivariance and Step~2 therefore put each $A_v$ in the
balanced space (X5.3). Its four summands have respective
dimensions $3,30,30,3$, as given by their full Lyndon bases.

The equation $D([a,b]+[c,d])=0$ splits into five
multidegrees $(j,j,4-j,4-j)$, $j=0,\ldots,4$. For each
column labelled $(v,w)$ its polynomial is, respectively,
\[
 [[b_w,a],b],\quad [a,[b_w,b]],\quad
 [[b_w,c],d],\quad [c,[b_w,d]].
\]
The row indexed by a degree-eight Lyndon word is its
associative-word coefficient. Restriction to these rows
is injective on the homogeneous free Lie algebra: in a
nonzero combination of Lyndon brackets its least Lyndon
word has the unchanged nonzero coefficient. Thus these
matrices impose the full specialness equation.

The complete certificate
\texttt{check\_genus\_two\_tangent\_space.py} constructs
every column and every row from these definitions. It
gives
\[
\begin{array}{c|rrrrr}
 j&0&1&2&3&4\\\hline
 \text{rows}&8&140&312&140&8\\
 \text{columns}&6&66&120&66&6\\
 \text{rank over }\mathbf Q&6&66&119&66&6\\
 \text{nonzero minor modulo }101&100&1&99&100&100.
\end{array}
\tag{X5.8}
\]
The output includes every selected row and column index.
The nonzero minors prove the rational lower ranks, hence
the same ranks over every $\mathbf Q_\ell$. In the middle
block the following exact nonzero kernel supplies the
matching upper rank:
\[
 A_a=A_b=b_{abcdcd}+2b_{acdbcd}+b_{acdcdb},\qquad
 A_c=A_d=b_{ababcd}+b_{abacdb}.
 \tag{X5.9}
\]
The script checks this equality in the full associative
word expansion, not only in the selected rows.

Equivalently, putting
\[
 z_1=[a,b],\quad z_2=[c,d],\quad C=[z_1,z_2],
 \quad A_1=[z_2,C],\quad A_2=-[z_1,C],
\]
the entire space is the line spanned by
\[
 D_*(a)=[A_1,a],\quad D_*(b)=[A_1,b],\quad
 D_*(c)=[A_2,c],\quad D_*(d)=[A_2,d].
 \tag{X5.10}
\]
The coefficients in (X5.9) give $-D_*$. The identity
$C=[z_1,\omega]$ proves that every value of $D_*$ lies
in the Lie ideal $(\omega)$.

\paragraph{7. The actual universal quotient kills this entire line.}
For $r=0,1$, put
$\Pi_r^{\mathrm{mod}}=\pi_1(M_{2,r})$, write $p_r$ for
its arithmetic projection and $\Phi_r$ for universal outer
monodromy, and set
\[
 F_r(m)=\Phi_r^{-1}(\Gamma_{2,r}(m)),\qquad
 Q_r(m)=\overline{\mathbf Q}^{\,p_r(F_r(m))}.
\]
Takao, RIMS~1709 (November~2010), Theorem~3.5, p.~21,
proves $Q_0(m)=Q_1(m)$ for every $m$\cite{AuditTakao2010RIMS1709}. This theorem includes
the last-puncture step in genus two; the older NTU~1994
affine-only assertion is not substituted for it.

The actual morphism $M_{2,1}\to M_{2,0}$ and the quotient
of the surface group obtained by filling the puncture
give the commutative diagram
\[
\begin{array}{ccc}
 \operatorname{gr}_6\Gamma^{\mathrm{univ}}_{2,1}
 &\longrightarrow&
 \operatorname{gr}_6\Gamma^{\mathrm{univ}}_{2,0}\\
 q_{6,1}\downarrow&&\downarrow q_{6,0}\\
 \operatorname{Gal}(Q_1(7)/Q_1(6))\otimes\mathbf Q_\ell
 &\xrightarrow{\ =\ }&
 \operatorname{Gal}(Q_0(7)/Q_0(6))\otimes\mathbf Q_\ell.
\end{array}
\tag{X5.11}
\]
The lower arrow is induced by equal subfields of
$\overline{\mathbf Q}$. The two arithmetic projections
agree on each element of $F_1(6)$, which proves
commutativity. The maps descend to the graded monodromy
images because an element in $F_r(7)$ acts trivially on
$Q_r(7)$. Clearing denominators, and then any torsion in
an integral graded quotient, proves the same statement
for rational classes. On the affine side $q_{6,1}$ is
the exact map (4.7.1) of Nakamura~1996.

Here the precise Lie map needed for filling is as follows.
The inclusion $U_5\hookrightarrow X_5$ induces the quotient
of $P$ by the closed normal subgroup generated by its
boundary loop $z$. Its associated graded Lie algebra is
generated by degree one, and $z$ has initial term
$\omega$. The induced graded map therefore factors as
the canonical surjections
\[
 L(V)\longrightarrow L(V)/(\omega)
       \longrightarrow
 \operatorname{gr}\pi_1(X_{5,\overline{\mathbf Q}})^{(\ell)}
       \otimes\mathbf Q_\ell.
\]
This factorization, proved directly from the actual
boundary relator, is sufficient; no injectivity of the
second arrow is used. Since every $D_*(v)$ belongs to
$(\omega)$, all its values vanish in the proper graded
Lie algebra. Thus the top horizontal arrow kills
\[
 \mathbf Q_\ell[D_*]\cap
     \operatorname{gr}_6\Gamma^{\mathrm{univ}}_{2,1}.
\]
Diagram (X5.11) proves that $q_6$ kills this intersection.
Step~6 places the \emph{actual} $m_0$ of (X5.7) in this
intersection, contradicting $q_6(m_0)\ne0$.
This excludes the remaining degree-one case.

\paragraph{8. Return to the actual subgroup.}
With degree one excluded, the detecting-cover line lies
in the boundary kernel by the weight argument of F.3.21.
After a finite extension splitting those boundary points,
the full normalizer acts on $A$ by exactly $\chi_\ell$.
Nakamura~1994, Theorem~2.1.1, p.~94, then puts $A$ in\cite{AuditNakamura1994}
actual cusp inertia. Every operation involving formal
conjugators above was used only to derive a contradiction
for an actual arithmetic element; the final subgroup and
recognition take place in the original pro-$\ell$ group.

\paragraph{An exact consequence for the exhibited cover.}
The actual morphism $q:Y_{20}\to U_5$ constructed in
F.3.19 induces compatible arithmetic fundamental-group
maps. For a procyclic subgroup of the geometric group of
$Y_{20}$ whose image under $q_*$ is nontrivial, its
normalizer maps to a normalizing subgroup with the same
open arithmetic image. The theorem places that image
in cusp inertia of $U_5$. Its image in $H_1(U_5)$ is
therefore zero, since $U_5$ has one puncture. Consequently
no subgroup with open arithmetic normalizer can have
nonzero image under the actual map
$H_1(Y_{20})\to V_\ell J_{X_5}$ induced by $q$.
This closes that particular CM-surface detecting target.
It gives no elliptic or genus-two detecting morphism for
an arbitrary finite cover in the full A.22 assertion.
\end{proof}

\begin{result}{F.3.23}{An actual rational detecting cover and a rational cusp}{Proved; arithmetic normal core and punctured maps}
Put $U_0=E_0\setminus\{O\}$, where $E_0:y^2=x^3-x$ over
$\mathbf Q$, and choose its rational point $b=(0,0)$.
Let $F=\pi_1(U_{0,\overline{\mathbf Q}},\bar b)$ and
$\Pi=\pi_1(U_0,\bar b)$. The actual point $b$ gives a section
$s:G_{\mathbf Q}\to\Pi$ and hence the explicit identification
$\Pi=F\rtimes s(G_{\mathbf Q})$.

Let $A\subset F$ be a closed procyclic subgroup, with normalizer
$N\subset\Pi$ whose image in $G_{\mathbf Q}$ is open. Fix a prime
$\ell$ and an open geometric subgroup $H\subset F$ containing
$A$, such that the image of a generator of the pro-$\ell$ part
$A_\ell$ in $H^{\mathrm{ab}}\otimes\mathbf Q_\ell$ is nonzero.
There is an actual finite connected etale cover $Y\to U_0$ over
$\mathbf Q$, with a rational point above $b$, whose geometric
subgroup $S$ is normal in $\Pi$, such that
$S\subset H$ and $A_\ell\cap S$ still has nonzero rational
homology in $Y$. If the original class has nonzero image in
the proper compactification of the cover corresponding to $H$,
the new class has nonzero proper homology as well.
The tangential descent choice also supplies a rational
cusp $y$ and an actual map to the once-punctured curve
$\overline Y\setminus\{y\}$ over $\mathbf Q$, preserving
any nonzero detected proper class and the open arithmetic
normalizer image.

\end{result}
\begin{proof}
We construct the arithmetic covering subgroup and track the detected class through every map.
\paragraph{The finite normal core and actual descent.}
The profinite group $F$ is topologically generated by two
elements. For each positive integer $d$ it has only finitely
many open subgroups of index $d$: each such subgroup is the
stabilizer of a point in a continuous permutation action
$F\to\mathfrak S_d$, and a homomorphism is determined by the
images of the two fixed generators. Therefore
\[
 S=\bigcap_{\gamma\in\Pi}\gamma H\gamma^{-1}
\]
is an intersection of finitely many open subgroups. It is
open in $F$, is normal in $\Pi$, and lies in $H$.
The subgroup
\[
 \Pi_Y=S\rtimes s(G_{\mathbf Q})\subset\Pi
\]
is closed and open, with $\Pi_Y\cap F=S$ and arithmetic image
all of $G_{\mathbf Q}$. The finite-etale-cover correspondence
gives an actual connected cover $Y\to U_0$ over $\mathbf Q$;
the equality of its arithmetic image with $G_{\mathbf Q}$
makes $Y$ geometrically connected. Since the stabilizer
$\Pi_Y$ contains $s(G_{\mathbf Q})$, the corresponding point
of the finite fibre over $b$ is fixed by $G_{\mathbf Q}$.
It is a rational point of $Y$. This constructs the descent
using the displayed arithmetic subgroup, rather than a
field-of-moduli assertion.

\paragraph{Preservation of the detected class and normalizer.}
The subgroup $A'=A\cap S$ has finite index in $A$ and is
characteristic in $A$: an open subgroup of a procyclic group
is uniquely determined by its index. In particular $N$
normalizes $A'$. On the nonzero pro-$\ell$ part,
$A'_\ell=\ell^r A_\ell$ for some $r\geq0$. Write $u$ for
a generator of $A_\ell$ and $u'=u^{\ell^r}$ for the resulting
generator. The actual inclusion $S\hookrightarrow H$ induces
\[
 S^{\mathrm{ab}}\otimes\mathbf Q_\ell
 \longrightarrow H^{\mathrm{ab}}\otimes\mathbf Q_\ell,
 \qquad [u']\longmapsto\ell^r[u].
\]
Its image is nonzero, so $[u']$ is nonzero. The finite map of
geometric covers extends uniquely to their smooth proper
compactifications. Functoriality of their proper homology
gives the identical formula there, and proves the additional
proper-homology assertion.

The subgroup $N'=N\cap\Pi_Y$ has finite index in $N$,
normalizes $A'$, and has open arithmetic image. Its kernel
in $S$ acts by inner automorphisms on $S$ and trivially on
$S^{\mathrm{ab}}$. Thus the one-dimensional span of $[u']$
is an actual Galois-stable line for that open arithmetic
image. When its proper image is nonzero, the corresponding
line in $V_\ell J_Y$ has the same character. No splitting of
the boundary extension is required.

\paragraph{All finite arithmetic conjugates on the same curve.}
Because $Y$ is defined over $\mathbf Q$, its Jacobian and
proper Tate module $V=V_\ell J_Y$ are defined over
$\mathbf Q$. Let $W=\mathbf Q_\ell[u']\subset V$ be the
nonzero proper line. Its stabilizer in $G_{\mathbf Q}$ is
open, so its orbit is finite. Choose representatives
$g_1=1,g_2,\ldots,g_t$ and use the actual section of $Y$
coming from its rational point to lift them to
$\widetilde g_i\in\Pi_Y$. Define
\[
 A_i=\widetilde g_i A'_\ell\widetilde g_i^{-1},\qquad
 N_i=\widetilde g_i N'\widetilde g_i^{-1},\qquad
 W_i=g_iW.
\]
Choose a finite Galois extension $L/\mathbf Q$ such that
$G_L$ is contained in every $p(N_i)$ and $L$ contains a
field defining all endomorphisms of $J_Y$.
It follows directly that
\[
 N_i\cap p^{-1}(G_L)\longrightarrow G_L
\]
is surjective for every $i$. Every $\sigma\in G_L$ therefore
has an actual normalizing lift for every one of the lines
$W_i$. The lifts need not be equal and need not form a
section; their arithmetic element is exactly the same.
This is the simultaneous-lift input used in F.3.22, now
constructed on the general rational detecting cover.

\paragraph{A rational boundary point and an actual once-punctured target.}
There is a second descent choice which supplies a rational
boundary point instead of the rational point in the open
curve. Choose a nonzero rational tangent vector at $O$ on
$E_0$ and its actual tangential section
$s_O:G_{\mathbf Q}\to\Pi$. Changing the chosen geometric
basepoint transports all the groups above by a compatible
fundamental-group isomorphism over $G_{\mathbf Q}$.
The arithmetic-normal geometric core $S$ is unchanged
under this transported description. Form instead
\[
          \Pi_{Y_O}=S\rtimes s_O(G_{\mathbf Q}).
\]
Let $I_O\subset F$ be the chosen geometric cusp inertia.
The cusp decomposition group in $\Pi$ is
$\Delta_O=I_O\rtimes s_O(G_{\mathbf Q})$. For the
corresponding branch $y$ in the smooth compactification
$\overline Y_O$, the actual decomposition subgroup is
\[
 \Delta_y=\Delta_O\cap\Pi_{Y_O}
          =(I_O\cap S)\rtimes s_O(G_{\mathbf Q}).
\]
This is the decomposition-group formula for the extension
of the discrete valuation at $O$ through the finite-cover
function field. Its image on the residue Galois group is
all of $G_{\mathbf Q}$. Since the constants are $\mathbf Q$
and $y$ is a closed point on a curve over $\mathbf Q$,
this says exactly that $\mathbf Q(y)=\mathbf Q$.
Thus $y$ is a rational boundary point. Equivalently, in
the compatible Puiseux fibre, $s_O(G_{\mathbf Q})$ fixes
the chosen branch, so it fixes its residue point.

Now use the actual open immersion
\[
             Y_O\hookrightarrow
             X:=\overline Y_O\setminus\{y\}.
\]
It fills all the other boundary points and is defined
over $\mathbf Q$. The induced geometric fundamental-group
homomorphism is surjective; on proper homology both
curves map to the same $V_\ell J_{\overline Y_O}$.
Consequently a nonzero detected proper class remains
nonzero under this actual map. Its procyclic subgroup
image is nontrivial in the maximal pro-$\ell$ geometric
group of $X$, and the normalizer image has the same
open arithmetic projection. The smooth rational point
$y$ has a one-dimensional rational tangent space, so
$X$ also has a rational tangential basepoint. The
Galois-conjugate lines and common arithmetic subgroup
constructed above can therefore all be formed on this
single once-punctured curve over $\mathbf Q$.
No claim that these lines span its entire homology is
needed for this construction.

\paragraph{Exact remaining map.}
The construction gives the actual finite set
$\{W_1,\ldots,W_t\}$ and the Galois-stable subspace
$\sum_iW_i\subset V_\ell J_Y$. It does not identify that
subspace with all of $V_\ell J_Y$, or with all embedding
lines of every CM factor. The CM quotient of F.3.18 is an
actual abelian-variety quotient with an actual map from
the proper curve; the further punctured-curve maps used
in F.3.19 are explicit special features of that example.
For the arbitrary detecting cover, the remaining task is
to calculate the image of the simultaneous-tangency
subspace under the actual universal quotient $q_6$, or to
construct enough punctured target maps with already
proved recognition. Neither map is supplied merely by
the rational-cover construction above. Full LM6 A.22
therefore retains its existing status.

\end{proof}

\begin{result}{F.3.24}{The genus-three universal quotient on balanced tangencies}{Proved; geometric monodromy and exact trace equations}
Let $V$ be a symplectic vector space of dimension six over
$\mathbf Q_\ell$, with ordered symplectic basis
$a,b,c,d,e,f$. Put $z_1=[a,b]$, $z_2=[c,d]$,
$z_3=[e,f]$, and $\omega=z_1+z_2+z_3$ in the free Lie
algebra $\mathbb L(V)$. Let $\mathfrak h_6$ be its special
degree-six derivations. Give the basis the characters
\[
 a:\alpha_1,\ b:t/\alpha_1,\quad
 c:\alpha_2,\ d:t/\alpha_2,\quad
 e:\alpha_3,\ f:t/\alpha_3
\]
of the full diagonal symplectic-similitude torus. Define
$T\subset\mathfrak h_6$ to be its $t^3$-weight subspace
of derivations tangent to every one of the six basis lines:
\[
 D(v)\in[\mathbb L_6(V),v]\qquad(v=a,b,c,d,e,f).
\]
Let $U_6\subset\mathfrak h_6$ be the rational universal
arithmetic degree-six monodromy image, with the unique
special representative used in F.3.22, and let
$q_6:U_6\twoheadrightarrow Q_6$ be Nakamura's universal
arithmetic quotient, whose similitude character is $t^3$\cite{AuditNakamura1996}.
Then
\[
                         q_6(T\cap U_6)=0.                 \tag{G3.1}
\]
This is a statement about the specified universal maps.
It does not assert that an arbitrary detecting cover in
LM6 A.22 supplies the full diagonal torus or all six
simultaneous basis-line constraints.

\end{result}
\begin{proof}
We identify the completed geometric image before imposing the finite trace equations.
\paragraph{1. The exact comparison with geometric monodromy.}
Use the moduli stack of genus-three curves equipped with
one marked point and a nonzero tangent vector there.
Its geometric fundamental group is the profinite
completion of the mapping class group of a genus-three
surface with one boundary circle fixed pointwise.
The universal tangential basepoint gives a monodromy
homomorphism into automorphisms of the punctured
fundamental group fixing the boundary word geometrically.
The arithmetic action sends this word to its cyclotomic
power. These are actual homomorphisms from the exact
arithmetic fundamental-group sequence of this moduli stack.
In particular its geometric monodromy image is normal
in the full arithmetic image.

The tangential arithmetic action can also be described
on each number-field fibre without choosing a rational
point of the open curve. A local uniformizer $t$ whose
first derivative agrees with the chosen tangent defines
the coefficient action of $G_K$ on
$\overline K((t^{1/n}))$ for all compatible $n$th roots,
fixing $t^{1/n}$. The resulting action on tame inertia
is the cyclotomic action. Replacing the coordinate by
$t+O(t^2)$ gives the compatible root substitutions with
leading coefficient one, which identify these tangential
fibre functors. This is the marked-tangent version of
the universal action.

Write $R$ for the proalgebraic Zariski closure over
$\mathbf Q$ of the discrete geometric monodromy in
automorphisms of the rational Malcev Lie algebra of
the punctured surface. At every finite nilpotent level,
the actual profinite geometric image lies in
$R(\mathbf Q_\ell)$: its dense discrete subgroup does,
and all defining polynomials are continuous over
$\mathbf Q_\ell$. Zariski closure commutes with this
scalar extension. Indeed, in every fixed degree the
linear equations imposed on the coefficients of a
vanishing polynomial have rational entries; their
solution space after scalar extension is exactly the
extended rational solution space.

The precise identification of the graded Lie algebra
of this image comes from the following opened primary
locators in Hain's final 1997 article\cite{AuditHain1997Torelli}. Theorem~3.4,
printed p.~605, makes the prounipotent radical of
relative completion a quotient of the Malcev
completion of the Torelli group. Lemma~4.5 and
Corollary~4.8, pp.~610--611, give the compatible
action and its weight filtration. The punctured
surface and boundary-fixed action are treated
explicitly in Section~14.4, p.~643: all the action
and filling maps are morphisms of mixed Hodge
structures, hence strict. The image of the relative
completion is $R$ by its universal property and
Zariski density. More explicitly, Theorem~4.10,
p.~611, identifies the Torelli Malcev weight
filtration with its lower central series, and
Remark~4.11, p.~613, records the compatible action
on the Malcev Lie algebras. Strictness of the
punctured action therefore makes its graded image
generated by the degree-one Johnson image.
Enomoto--Satoh, Theorem~3.5 and Proposition~3.6,
p.~640, identify that generated algebra with
the actual rational Johnson image. Consequently
\[
 \operatorname{gr}_{-m}^W\operatorname{Lie}R_u
   =\operatorname{Im}\bigl(\tau_{g,1}(m)\otimes\mathbf Q\bigr)
   =:J_m,                                                  \tag{G3.2}
\]
as subspaces of the actual derivation space. Here the
weight filtration on the image is the filtration
induced from its action on the punctured Malcev Lie
algebra. Strictness is essential in this assertion.
This argument does not identify closures of individual
discrete filtration kernels by density alone.

The canonical special representative does not destroy
normalization of $R$. For an arithmetic element trivial
on homology and in outer filtration degree six, its
tangential action fixes the full boundary logarithm
$Z$, whose initial term is $\omega$. Any lower-degree
inner part has degree-two term $c\omega$ and no
degree-one term. Remove it by $\exp(-c\operatorname{ad}Z)$.
This is a point of the geometric algebraic image: the
boundary Dehn twist has logarithm
$\operatorname{ad}Z$, and its integer powers are
Zariski dense in the corresponding additive one-parameter
group. It therefore preserves normalization of $R$.
The free-Lie centralizer of $\omega$ then removes
degrees three through five exactly as in F.3.22.
The resulting degree-six derivation is the unique
special representative. It follows, by taking the
leading term of its action on $\operatorname{Lie}R$, that
\[
                       [D,J_1]\subset J_7
                       \qquad(D\in U_6).                  \tag{G3.3}
\]
The commutator used here is an actual normalizer action
on the algebraic image, followed by its strict graded
map (G3.2).

\paragraph{2. The trace equation and its precise genus range.}
The degree-one Johnson image is
$J_1=\Lambda^3V$, with the derivation corresponding to
$u\wedge v\wedge w$ given by
\[
 \delta_{uvw}(x)=\langle x,u\rangle[v,w]
                +\langle x,v\rangle[w,u]
                +\langle x,w\rangle[u,v].                 \tag{G3.4}
\]
Define $\operatorname{ES}_k$ on a degree-$k$ special
derivation by the exact composition
\[
 \mathfrak h_k\hookrightarrow V\otimes\mathbb L_{k+1}(V)
 \hookrightarrow V^{\otimes(k+2)}
 \xrightarrow{\,\text{contract slots }1,2\,}V^{\otimes k}
 \longrightarrow (V^{\otimes k})_{\mathbf Z/k}.             \tag{G3.5}
\]
The tensor representing $D$ is
$\sum_i(x_i\otimes D(y_i)-y_i\otimes D(x_i))$.
Changing the simultaneous duality convention changes
only the overall sign and does not change the kernel.

Enomoto--Satoh's final 2014 article\cite{AuditEnomotoSatoh2014}, Theorem~3.5 and
Proposition~3.6 on p.~640 and the actual diagram in
Section~7.1, p.~656, give
$\operatorname{ES}_7(J_7)=0$ in their stable genus
range. To apply it here, attach six additional handles,
giving the actual inclusion of boundary-fixed mapping
class groups $\Gamma_{3,1}\to\Gamma_{9,1}$ and the
symplectic inclusion $V\hookrightarrow V'$. Extend a
derivation by zero on the added generators. Both the
Johnson map and (G3.5) commute with this inclusion.
The map on cyclic tensor coinvariants is injective:
projection killing the added letters is its left
inverse. Since $9\geq7+2$, the stated stable theorem
therefore gives the required vanishing at genus three.
Combining it with (G3.3) gives
\[
       \operatorname{ES}_7([D,\delta])=0
       \quad(D\in U_6, \delta\in\Lambda^3V).              \tag{G3.6}
\]
No stable-range formula is directly substituted at
genus three without these explicit maps.

\paragraph{3. The complete tangent space.}
Every element of $T$ has a unique expression
$D(v)=[A_v,v]$ with $A_v$ of degree six and character
$t^3$. Uniqueness follows from the centralizer of a
free generator, and the character assertion follows
by the injectivity of $A_v\mapsto[A_v,v]$ in this degree.
Its balanced multidegrees have three pairs and are
indexed by $(m_1,m_2,m_3)$ with sum three.
The specialness equation $D\omega=0$ splits into
blocks indexed by $(n_1,n_2,n_3)$ with sum four.
Up to permutation they have the following sizes and
rational ranks:
\[
\begin{array}{c|c|c|c}
 (n_1,n_2,n_3)&\text{rows}&\text{columns}&\text{rank}\\\hline
 (4,0,0)&8&6&6\\
 (3,1,0)&140&66&66\\
 (2,2,0)&312&120&119\\
 (2,1,1)&1260&360&359.
\end{array}                                               \tag{G3.7}
\]
The first three blocks are exactly the two-handle
blocks certified in F.3.22, with absent letters omitted.
For the new block a nonzero $359$-minor modulo $101$
has determinant $100$, and the following exact integer
kernel provides the matching upper bound.

For each pair $i<j$, put $C_{ij}=[z_i,z_j]$ and define
$P_{ij}$ by common conjugators on each handle:
\[
 A_i=[z_j,C_{ij}],\qquad A_j=-[z_i,C_{ij}],\qquad A_k=0.
\]
For each $i$, with $\{i,j,k\}=\{1,2,3\}$, define $T_i$ by
\[
 \begin{split}
 A_i&=[z_j,[z_i,z_k]]+[z_k,[z_i,z_j]],\\
 A_j&=-[z_i,[z_i,z_k]],\qquad
 A_k=-[z_i,[z_i,z_j]].
 \end{split}                                               \tag{G3.8}
\]
Each handle's two generators have the same displayed
conjugator. Jacobi gives $\sum[A_i,z_i]=0$ exactly.
The six elements
\[
                  P_{12},P_{13},P_{23},T_1,T_2,T_3
\]
are therefore a basis of $T$ by (G3.7).

\paragraph{4. The exact normalizer trace kernel.}
Apply (G3.5) to brackets with the twenty derivations
(G3.4) indexed by triples of distinct basis letters.
Using cyclic words of length seven as coordinates
gives an integer matrix with $2208$ nonzero rows and
six columns in the order above. Four of its rows,
labelled by $(\delta,\text{cyclic word})$, are
\[
 (abc,abccdcd),\quad(abc,aabcfeb),\quad
 (abc,abccdef),\quad(abe,aabcdeb).
\]
On the first four columns their matrix is
\[
 \begin{pmatrix}
 -1&0&0&0\\
 0&-1&0&1\\
 0&0&1&-1\\
 -1&0&0&1
 \end{pmatrix},\qquad\det=1.                              \tag{G3.9}
\]
Exact expansion of all the bracket traces gives identical
columns for $T_1,T_2,T_3$. Thus the full kernel, with
both bounds certified over the integers, is
\[
 \{D\in T:\operatorname{ES}_7([D,\Lambda^3V])=0\}
       =\mathbf Q_\ell(T_2-T_1)+\mathbf Q_\ell(T_3-T_1).
                                                               \tag{G3.10}
\]
The complete code keeps the exact associative-word
expansions; the minor alone supplies only the lower
bound and is used together with the full-column equality.

\paragraph{5. The arithmetic quotient.}
Permuting the three symplectic handles conjugates the
three derivations $T_i$ transitively. Every
$\operatorname{Sp}(V)$-invariant linear functional on
$\mathfrak h_6$ therefore vanishes on both differences
in (G3.10). The universal image $U_6$ is an algebraic
$\operatorname{Sp}(V)$-submodule: conjugation by the
geometric mapping-class image preserves it, and that
image is Zariski dense in $\operatorname{Sp}(V)$ on
homology. Nakamura's equivariance makes $q_6$ invariant
under this symplectic group, since its target character
is $t^3$. Complete reducibility in characteristic zero
extends this invariant linear map from $U_6$ to
$\mathfrak h_6$ (or, equivalently, gives an equivariant
projection $\mathfrak h_6\to U_6$). Equations (G3.6)
and (G3.10) now imply (G3.1).

\paragraph{Scope and the failure of the simpler filling calculation.}
The two-handle derivation $P_{12}$ has nonzero image on
the proper genus-three quotient. The actual free-Lie
substitution
\[
 a,c,e\mapsto x,\qquad b\mapsto y,\qquad d\mapsto z,
 \qquad f\mapsto-y-z
\]
kills $\omega$, whereas the image of $P_{12}(a)$ has
coefficient one on the word $xxyxzxz$. This proves the
nonvanishing through an explicit quotient morphism.
The additional geometric-normalizer equation (G3.6)
is what removes the relevant arithmetic quotient.
No element of a common-kernel subgroup with an open
arithmetic normalizer is constructed or ruled out on
an arbitrary LM6 A.22 detecting cover by this lemma
alone. Its exact further requirements are the actual
CM-character projection and enough simultaneous
orbit-line constraints on that cover.

\end{proof}

\paragraph{Unresolved obligations retained.}
The inventory does not promote the six unspecified global essential images, general Nori Hodge fullness, general numerical equals homological equivalence, the historical integral-model identification of F.3.12--F.3.13, the residual LM6 subgroup problem of F.3.10, the unlocated historical standard-conjecture reduction through singular hypersurfaces in cote14, the two undefined entries of the historical A--D dictionary, or the unexplained pre-minimal normalization to established results. The arithmetic period and splitting obligations are exactly those retained in A.2--A.5. The requested named-conjecture implications prove their stated consequences only; they are not assertions of their premises.


\begin{result}{F.3.25}{Rational CM endomorphism orbits and genus-three normalizer recognition}{Proved; actual endomorphism and arithmetic maps}
Let $A/\Q$ be absolutely simple of dimension $g$, with
$E=\operatorname{End}^0_{\Qb}(A)$ a CM field of degree $2g$.
Then $E/\Q$ is Galois, the actual endomorphism action
$G_{\Q}\to\operatorname{Aut}_{\Q}(E)$ is surjective, and
its kernel is $G_E$. At a completely split prime its
actual action is transitive on the $2g$ CM eigenlines,
which form conjugate symplectic pairs.

For every smooth proper genus-three curve $C/\Q$ with
absolutely simple Jacobian and point $y\in C(\Q)$,
every prime $\ell$, and every finite extension $K/\Q$,
a nontrivial closed procyclic subgroup of
$\pi_1((C\setminus\{y\})_{\overline K})^{(\ell)}$
whose arithmetic normalizer has open image in $G_K$
is contained in an actual cusp inertia subgroup.
\end{result}
\begin{proof}
We first determine the exact endomorphism action from the rational
Hodge filtration, then apply the universal quotient of F.3.24.

\paragraph{The type and its right stabilizer.}
Let $L$ be the normal closure of $E$, put
$G=\operatorname{Gal}(L/\mathbf Q)$ and
$H=\operatorname{Gal}(L/E)$, and fix the inclusion
$E\hookrightarrow L\subset\overline{\mathbf Q}$.
The embeddings of $E$ are identified with the right
cosets $G/H$ by $gH\mapsto g|_E$. Let $S\subset G$
be the inverse image of the CM type occurring in
$\operatorname{Fil}^1H^1_{\mathrm{dR}}(A/\mathbf Q)
\otimes\overline{\mathbf Q}$. Thus $S$ is right
$H$-invariant and
\[
                           S\sqcup cS=G,                  \tag{CMQ.1}
\]
where $c$ is the central complex conjugation of the
Galois CM field $L$. Here this assertion about the normal
closure has an explicit field proof. The field $L$ is
the compositum of the conjugate CM fields $g(E)$.
Complex conjugation on each $g(E)$ is its intrinsic
CM involution, so every embedding carries it to the
CM involution on the image. Hence complex conjugation
commutes with $G$ on each generating subfield and
therefore on $L$. Its fixed field is totally real:
for $x$ fixed by $c$ and $g\in G$, one has
$c(gx)=g(cx)=gx$. Since $L$ contains the totally
imaginary field $E$, the involution is nontrivial.
Thus $L/L^{\langle c\rangle}$ is a totally
imaginary quadratic extension of a totally real field.

Its right stabilizer is exactly $H$. Indeed let
\[
                    K=\{k\in G:Sk=S\}.
\]
It contains $H$, and $c\notin K$ by (CMQ.1).
The fixed field $F=L^K$ is a CM subfield of $E$:
centrality of $c$ makes its restriction to $F$ a
nontrivial involution, its fixed field lies in the
totally real field $L^{\langle c\rangle}$, and no
embedding of $F$ is real because $c$ is central and
acts nontrivially. Right $K$-invariance says precisely
that the type of $E$ is the inverse image, under
$\operatorname{Hom}(E,\mathbf C)\to
 \operatorname{Hom}(F,\mathbf C)$, of a CM type of $F$.

If $K$ were strictly larger than $H$, this would give
an actual proper abelian subvariety of $A$. To see the
map explicitly, identify $H_1(A(\mathbf C),\mathbf Q)$
with the rank-one $E$-module $E$. The induced-type
condition says that its complex structure, expressed
as multiplication in $E\otimes\mathbf R$, belongs to
$F\otimes\mathbf R$. Choose an $F$-basis of $E$ and
project onto one coordinate. The resulting rational
idempotent $p:E\to E$ commutes with the complex
structure and has rational rank $[F:\mathbf Q]$.
Choose an integer $n>0$ clearing its denominators on
$H_1(A(\mathbf C),\mathbf Z)$. The map $np$ gives a
complex-linear homomorphism of the actual analytic
torus $A(\mathbf C)$. Its graph is an algebraic
subvariety of the projective product by Chow's theorem,
so it is an abelian-variety endomorphism. The base-change
map
$\operatorname{End}_{\overline{\mathbf Q}}(A)
 \to\operatorname{End}_{\mathbf C}(A_{\mathbf C})$
is an isomorphism: spread any complex endomorphism over
a finitely generated extension of $\overline{\mathbf Q}$,
which is regular, and apply the graph-descent assertion
of Milne's original Corollary~20.4(b), pp.~146--147~\cite{AuditMilne1986AV}.
This is also the rigidity argument obtained by spreading
over a connected base and using the discrete Hom functor.
The actual endomorphism $np$ therefore descends to
$\overline{\mathbf Q}$ and has image of dimension
$[F:\mathbf Q]/2$ strictly between zero and $g$.
Absolute simplicity excludes this. Hence
\[
                              K=H.                        \tag{CMQ.2}
\]

\paragraph{The semilinear de Rham equation.}
The actual action $a_\sigma$ on $E$ corresponds to
some $n_\sigma H\in N_G(H)/H$:
\[
                       a_\sigma=n_\sigma|_E.
\]
For an $E$-eigenvector $w$ of de Rham character
$\varphi$ in the rational Hodge filtration, applying
$\sigma$ coefficientwise to $e^*w=\varphi(e)w$ gives
\[
             a_\sigma(e)^*\sigma(w)
                    =\sigma(\varphi(e))\sigma(w).
\]
The new character is
$\sigma\circ\varphi\circ a_\sigma^{-1}$.
The Hodge filtration is defined over $\mathbf Q$,
so the corresponding type identity is
\[
                         \sigma S=S n_\sigma.             \tag{CMQ.3}
\]
The right action has the stated orientation because
composing $g|_E$ with $n|_E$ gives $(gn)|_E$.
This also verifies compatibility with the
homomorphism law for actual endomorphism conjugation.

If $\sigma$ fixes $L$ pointwise, the left side of
(CMQ.3) is $S$. Equations (CMQ.2)--(CMQ.3) force
$n_\sigma\in H$ and $a_\sigma=1$. Thus $a$ factors
through the finite group $G$. For this factored
homomorphism, its kernel is exactly
\[
                 H^*=\{g\in G:gS=S\},                    \tag{CMQ.4}
\]
the left stabilizer of $S$. In particular $H^*$ is
normal in $G$.

Normality turns left stabilization into right
stabilization. For $s\in S$ and $h\in H^*$, write
$sh=(shs^{-1})s$. Since $shs^{-1}\in H^*$, the
right side belongs to $S$. Hence $SH^*=S$ and
\[
                                H^*\subset H.              \tag{CMQ.5}
\]
Now the exact index comparison is
\[
 [G:H]\leq [G:H^*]=|\operatorname{im}a|
       \leq|N_G(H)/H|\leq[G:H].                            \tag{CMQ.6}
\]
Every inequality is an equality. Thus $H^*=H$,
$N_G(H)=G$, and the endomorphism action is onto
$N_G(H)/H$. The subgroup $H$ is normal, so $E/\mathbf Q$
is Galois. Since $L$ was its normal closure, $L=E$
and $H=1$. Formula (CMQ.4) therefore proves that the
kernel of the original action is exactly $G_E$.

\paragraph{The actual Tate-module permutation.}
At a completely split prime let
$V_\varphi\subset V_\ell A$ be the eigenline for
$\varphi:E\to\mathbf Q_\ell$. The arithmetic action
is $\mathbf Q_\ell$-linear, and its endomorphism
conjugation gives the exact formula
\[
 \rho_\ell(\sigma)V_\varphi
                    =V_{\varphi\circ a_\sigma^{-1}}.       \tag{CMQ.7}
\]
Surjectivity of $a$ makes the orbit all the
$2g$ embeddings. A polarization over $\mathbf Q$
has Rosati involution equal to complex conjugation
on $E$. Its alternating pairing is therefore zero
between all eigenlines except the pairs
$V_\varphi,V_{\bar\varphi}$, and is nondegenerate
on each such pair. Scaling their generators yields
an actual symplectic basis made of the orbit lines.
This supplies both the arithmetic permutation and
the pairing needed by the simultaneous normalizer
calculation.


\paragraph{The projector after the actual simultaneous constraints.}
Keep the genus-three universal submodule $U_6\subset\mathfrak h_6$
and quotient $q_6$ of F.3.24. Fix any symplectic basis of $V$.
Let $T_{\mathrm{all}}$ be the full subspace of special degree-six
derivations tangent to all six basis lines, without imposing a
character condition. Then
\[
                         q_6(U_6\cap T_{\mathrm{all}})=0.       \tag{QCM.1}
\]

Indeed, let $T_{\mathrm{diag}}$ be the full split diagonal
similitude torus for this basis, with multiplier $t$. The
subspace $T_{\mathrm{all}}$ is stable under it: for a diagonal
element $h$, conjugation sends $D(v)=[A_v,v]$ to another
bracket with the same line $\mathbf Q_\ell v$. The universal
submodule $U_6$ is stable under $\operatorname{GSp}(V)$.
Its symplectic stability was proved in F.3.24; scalar matrices
act on degree-six derivations by their sixth power, so this
stability extends across the central isogeny
$\operatorname{Sp}(V)\times\mathbf G_m\to\operatorname{GSp}(V)$.

There is an explicit $\mathbf Q_\ell$-linear weight projector
\[
 e_{t^3}:U_6\cap T_{\mathrm{all}}
       \longrightarrow U_6\cap T_{\mathrm{all}}[t^3].          \tag{QCM.2}
\]
To construct it, choose $h\in T_{\mathrm{diag}}(\mathbf Q_\ell)$
at which $t^3$ differs from every other weight in the finite
dimensional derivation space. Such an $h$ exists because a
finite union of the proper zero sets of the Laurent
polynomials $\lambda-t^3$ does not cover the split torus
over the infinite field $\mathbf Q_\ell$. The Lagrange
interpolation polynomial in the operator $h$ gives (QCM.2).
Every factor preserves both displayed subspaces. The
equivariance of $q_6$ gives
\[
                              q_6 e_{t^3}=q_6.
\]
F.3.24 kills its target and proves (QCM.1).
In an arithmetic application the six tangency constraints
are first proved for an actual Galois element. The operator
(QCM.2) is then applied to that already established linear
constraint space inside the universal image. No lift of
this later algebraic projection to an arithmetic element
is asserted or needed. The same argument applies to the
four-line genus-two quotient calculation in F.3.22.



\paragraph{The local algebra and the nonsplit case.}
Put $B=J_C$ and $V=V_\ell B$. A degree-one arithmetic
normalizer line is stable on an open subgroup $G_L$.
Enlarge $L$ to define all geometric endomorphisms.
The semisimplicity and endomorphism statements of
Faltings, original Section~5, S\"atze~3--4, p.~360
\cite{GenFaltings1983}, give a rank-one idempotent
in $\operatorname{End}^0_{\Qb}(B)\otimes\Q_\ell$.
The exact rank formula and actual quotient construction
of F.3.18 apply. Since $B$ is absolutely simple, the
simple quotient there is isogenous to $B$ itself;
the formula forces the full geometric endomorphism
algebra to be a CM field $E$ of degree~6. The preceding
CM argument therefore makes $E/\Q$ Galois.

The comparison respecting actual endomorphisms makes
$V$ a free rank-one module over
\[
                E\otimes\Q_\ell=\prod_{\lambda\mid\ell}E_\lambda.
                                                               \tag{QCM.3}
\]
Every idempotent in this commutative product selects
some factors, and its $\Q_\ell$-rank is the sum of
the selected local degrees $[E_\lambda:\Q_\ell]$.
Galois transitivity on primes of $E$ makes all those
degrees equal, including at ramified primes.
Thus a rank-one idempotent exists only when every
local degree is one, precisely when $\ell$ splits
completely in $E$. At every other prime the
hypothesized degree-one normalizer line is excluded
by this actual endomorphism-algebra calculation.
This argument applies after every finite base
extension because the endomorphisms have already
all been defined.

\paragraph{The actual split-prime orbit of one normalizer.}
At a split prime the rank-one idempotent selects
exactly one embedding eigenline. Its actual
$G_{\Q}$ orbit is all six lines by (CMQ.7).
They are distinct as characters on every sufficiently
small open arithmetic subgroup: two equal characters
would allow a nonzero off-diagonal intertwiner,
contradicting Faltings' equality of the full commuting
algebra with the diagonal algebra (QCM.3).
The rational polarization supplies the three
nondegenerate conjugate pairs described above.

Choose a nonzero rational tangent vector at $y$.
It gives the actual tangential section of the
arithmetic group over $\Q$, as in F.3.24. For six
$\tau_i\in G_{\Q}$ carrying the initial eigenline
to the six specified lines, conjugate the original
procyclic subgroup and its normalizer by these
actual section lifts. Write the resulting subgroups
as $A_i$ and $N_i$. Their arithmetic images are
$\tau_i p(N)\tau_i^{-1}$ and are open. Choose a
finite Galois extension $L/\Q$ containing a field
defining all endomorphisms, with $G_L$ contained
in every one of these six images and in $G_K$.
Then each $N_i\cap p^{-1}(G_L)$ surjects onto
this same $G_L$. All groups act on the same
geometric punctured curve; no target curve or
homology space is replaced in this construction.

The arithmetic quotient and finite-base invariance
in Nakamura's original Lemma~4.5, Corollary~4.6,
and (4.7.1)--(4.7.2), pp.~115--117
\cite{AuditNakamura1996}, supply an actual
$\sigma\in G_L(6)$ whose degree-six class
$D\in U_6$ has $q_6(D)\ne0$. Here the quotient
is nonzero in degree~6 and has similitude character
$t^3$, with exactly the maps fixed in F.3.24.
For each $i$, lift this same $\sigma$ to $N_i$.
Because it acts trivially on $V$ and the linear
term of a generator of $A_i$ is nonzero, its
normalizer exponent is exactly one.

We spell out the rational correction so that no
primitivity of the original integral generator is
used. Scale that nonzero linear term to a basis
vector $v_i$, and write its full logarithm
$U_i=v_i+$ higher degrees. The canonical special
representative of $\sigma$ gives an equality
\[
             \widetilde\sigma(U_i)=\exp(\operatorname{ad}W_i)U_i
                                                               \tag{QCM.4}
\]
for a rational completed free-Lie element $W_i$.
Its degree-one term centralizes $v_i$, and hence
is a scalar multiple of $v_i$. Multiplying the
conjugating element by the appropriate rational
power of $\exp(U_i)$ removes that term while
leaving its action on $U_i$ unchanged. In each
successive degree below~6, the centralizer of
$v_i$ in the free Lie algebra forces the next
term to vanish. The degree-seven part of
(QCM.4) is therefore
\[
                         D(v_i)=[(W_i)_6,v_i].              \tag{QCM.5}
\]
This calculation is performed separately for each
$i$ with the same $D$. It proves that $D$ is
in $U_6\cap T_{\mathrm{all}}$. Equation (QCM.1)
then contradicts $q_6(D)\ne0$ and excludes the
split-prime degree-one case as well.

\paragraph{The exact passage from the first Magnus term to inertia.}
For completeness, the reduction used here is the
finite-cover argument of F.3.21 with the present
curve substituted through its actual maps. Let
$u$ generate the nontrivial procyclic subgroup.
A finite $\ell$-group quotient in which $u$ survives,
followed by the preimage of its cyclic image,
gives a finite etale cover $Y$ where $u$ has
nonzero rational abelian image. Its arithmetic
model exists over a finite extension. The
intersection of the normalizer with the covering
group still has open arithmetic image; its
geometric kernel acts inner on that covering
group and hence trivially on its homology.
Thus the detected homology line has exactly the
same character $\theta$ as the first nonzero
Magnus term $u_d\in\mathbb L_d(V)$.

The latter is pure of weight $d$ because the
original curve has exactly one puncture.
A nonzero image of the detected line in proper
homology of the compactification of $Y$ would,
at a common good Frobenius, force $d=1$.
That case was just excluded at every prime.
The line therefore lies in the actual kernel
of the homology map to the compactification,
namely the quotient of the boundary permutation
module tensored with $\Q_\ell(1)$ by its diagonal
copy. Over a finite extension defining all
boundary points, this entire kernel has character
$\chi_\ell$, so $\theta=\chi_\ell$ there.
The first Magnus term also shows that $\theta$
is trivial on the original geometric normalizer.
A lift from the intersected normalizer and any
lift from the full normalizer of the same
arithmetic element differ by that geometric
kernel. Hence the full normalizer over this
extension acts by $\chi_\ell$ and still has
open arithmetic image. Nakamura's original
Theorem~2.1.1, p.~94 \cite{AuditNakamura1994},
now places the original subgroup in actual
cusp inertia. Finite arithmetic base changes
leave the geometric group and this inertia
conjugacy class unchanged.
\end{proof}

\paragraph{Scope and source comparison.}
The general CM assertion has been derived from the exact
rational filtration and stabilizer maps, with no novelty
claim. For comparison, the opened Lemma~2.5 of
\cite{AuditGoodman2026} states the usual reflex-field
endomorphism formula under its explicit Galois-CM-field
hypothesis and points to Shimura's 1998 Proposition~30,
p.~65~\cite{AuditShimura1998Gate}; that original book passage has not been obtained
and is not an input to the proof above.
The curve theorem concerns a whole absolutely simple
Jacobian of dimension~3. A general detecting cover in
F.3.23 comes with its actual CM quotient, whose dimension
may be smaller than that of its Jacobian. The maps to
that quotient and its conjugates are retained; the
proof does not identify their orbit span with the whole
curve homology. In particular an original higher-genus
cover of $E_0\setminus O$ still has its norm quotient
to $E_0$. The full LM6 A.22 subgroup assertion remains
open in its stated scope.


\begin{result}{F.3.26}{The universal quotient on full symplectic basis tangencies in every genus}{Proved; exact support and normalizer maps}
Let $g\geq3$, let $V$ have symplectic basis
$x_1,y_1,\ldots,x_g,y_g$ over $\Q_\ell$, and put
$z_i=[x_i,y_i]$, $\omega=\sum_i z_i$.
Write $\mathfrak h_{g,6}$ for the special degree-six
derivations of $\mathbb L(V)$, and $U_{g,6}$ for the
universal arithmetic image with its canonical special
representative and the actual quotient
\[
                 q_{g,6}:U_{g,6}\longrightarrow Q_{g,6}
\]
of Nakamura, as in F.3.24~\cite[pp.~115--117]{AuditNakamura1996}. Its character is the cube
of the similitude multiplier. Define $T_{g,\mathrm{all}}$
by tangency to every one of the specified basis lines:
\[
 D(x_i)=[A_i,x_i],\quad D(y_i)=[B_i,y_i],
 \qquad A_i,B_i\in\mathbb L_6(V).
\]
Then
\[
                 q_{g,6}(U_{g,6}\cap T_{g,\mathrm{all}})=0.
                                                               \tag{AG.1}
\]
This statement retains the full homology space of the
specified curve. A proper CM summand supplies only its
own line constraints and is not substituted for $V$.

\end{result}
\begin{proof}
We reduce every balanced tangency to its exact support on at most four handles and retain the actual universal submodule throughout.
\paragraph{1. The actual geometric-normalizer equation in every genus.}
The construction of F.3.24 uses the moduli of a genus-$g$
curve with one point and one nonzero tangent vector.
Its boundary-fixed geometric monodromy has proalgebraic
image $R_g$ in automorphisms of the rational punctured
Malcev group. The actual profinite geometric image lies
in $R_g(\Q_\ell)$ at each finite nilpotent level by
polynomial continuity and the dense discrete image.
Its algebraic graded image is the Johnson image $J_m$:
Hain's original Theorem~3.4, p.~605, explicitly applies
for all $g\geq3$; Theorem~4.10, p.~611, identifies
the Torelli Malcev weight filtration with its lower
central series for $g\ne2$; Remark~4.11, p.~613,
and the boundary action in Section~14.4, p.~643,
give the strict action maps. Together with the exact
degree-one generation in Enomoto--Satoh, Theorem~3.5
and Proposition~3.6, p.~640, these are the same
typed maps used in F.3.24, now in genus~$g$~\cite{AuditHain1997Torelli,AuditEnomotoSatoh2014}.

Tangential arithmetic monodromy normalizes this actual
algebraic image. Removing the lower inner boundary
term by $\exp(-c\operatorname{ad}Z)$ remains inside
$R_g$, because the boundary Dehn twist is a geometric
element and its powers are Zariski dense in that
one-parameter group. Here $Z$ is the full boundary
logarithm with initial term $\omega$. The free-Lie
centralizer removes the remaining lower inner terms.
Thus the canonical special representative satisfies
\[
       [D,J_1]\subset J_7\qquad(D\in U_{g,6}),
       \qquad J_1=\Lambda^3V.                              \tag{AG.2}
\]
These are leading terms of the actual normalizer
action at a finite nilpotent level. No closure of
an individual discrete filtration kernel is inferred.

Use the exact contraction and cyclic-coinvariant
map $\operatorname{ES}_7$ of F.3.24. It vanishes on
$J_7$ for every $g\geq3$. For $g\geq9$ this is
the stated stable scope of Enomoto--Satoh. For
$3\leq g<9$, the actual boundary-fixed inclusion
of surfaces from genus~$g$ to genus~9 extends
derivations by zero. Both the Johnson map and
the contraction commute with this inclusion;
the cyclic-tensor projection killing the added
letters is a left inverse. Consequently
\[
     \operatorname{ES}_7([D,\delta])=0
      \quad(D\in U_{g,6},\ \delta\in\Lambda^3V).           \tag{AG.3}
\]

\paragraph{2. The complete balanced tangent space has support at most four.}
Give the full diagonal torus its characters
$x_i:a_i$, $y_i:t/a_i$. The weight-$t^3$ subspace
$T_g$ of $T_{g,\mathrm{all}}$ has unique conjugators
$A_i,B_i$ of degree~6 and weight $t^3$. They are
balanced: every handle occurs equally often through
its two generators, and the sum of the handle
multiplicities is three. The equation $D\omega=0$
therefore splits into blocks indexed by
\[
               n=(n_1,\ldots,n_g),\qquad \sum_i n_i=4.
\]
Every nonzero $n_i$ is positive, so each block involves
at most four handles. Up to permutation its possible
types, sizes, and rational ranks are
\[
\begin{array}{c|r|r|r}
\text{positive entries of }n&\text{rows}&\text{columns}&\text{rank}\\\hline
(4)&8&6&6\\
(3,1)&140&66&66\\
(2,2)&312&120&119\\
(2,1,1)&1260&360&359\\
(1,1,1,1)&5040&960&958.
\end{array}                                               \tag{AG.4}
\]
The first four blocks are the exact matrices already
certified in F.3.22 and F.3.24. The last is the only
new block, certified by
\texttt{code/check\_genus\_four\_tangent\_block.py}.
Its ordered $958$-minor has determinant~1 modulo~101.
Two exact independent integer kernels give the
matching upper bound; their independence minor is
$-1$. All matrix entries are coefficients of standard
Lyndon commutators in the free associative algebra.
Thus this finite calculation proves the full rational
rank, rather than only a numerical approximation.

Here is the resulting complete basis description.
There is one $P_{ij}$ on each pair of handles, with
common conjugators on each paired pair of generators
\[
 A_i=B_i=[z_j,[z_i,z_j]],\quad
 A_j=B_j=-[z_i,[z_i,z_j]],
\]
and zero conjugators elsewhere. On each triple
$\{i,j,k\}$ there are three $T_{i;jk}$, indexed
by the repeated handle $i$, with
\[
 \begin{split}
 A_i=B_i&=[z_j,[z_i,z_k]]+[z_k,[z_i,z_j]],\\
 A_j=B_j&=-[z_i,[z_i,z_k]],\\
 A_k=B_k&=-[z_i,[z_i,z_j]].
 \end{split}                                               \tag{AG.5}
\]
On each four-element set there are two elements
$Q^{(1)},Q^{(2)}$. Number those four handles
$1,2,3,4$, and let $L_{ijk}$ be the standard
Lyndon bracket for the word $ijk$ in formal
ordered letters $z_1,z_2,z_3,z_4$. Its values
for the orders used below are fixed by
$L_{ijk}=[z_i,[z_j,z_k]]$ for $i<j<k$ and
$L_{ikj}=[[z_i,z_k],z_j]$ for $i<j<k$.
The common conjugators of the two generators
on each handle are the following columns:
\[
\begin{array}{c|c|c}
 &Q^{(1)}&Q^{(2)}\\\hline
1&L_{234}+L_{243}&-L_{243}\\
2&L_{143}&-L_{134}-L_{143}\\
3&-L_{142}&-L_{124}\\
4&L_{123}&L_{132}.
\end{array}                                               \tag{AG.6}
\]
The program verifies $\sum_i[A_i,z_i]=0$ as an
exact associative-word identity for each column,
then substitutes $z_i=[x_i,y_i]$ and verifies the
two corresponding full $960$-coordinate kernel
vectors. Since the displayed elements occupy
different blocks except for the independent pair
on each quadruple, (AG.4) proves that all these
elements form a basis of $T_g$. Its dimension is
\[
                  \binom{g}{2}+3\binom{g}{3}+2\binom{g}{4}.
                                                               \tag{AG.7}
\]

\paragraph{3. Exact restriction of the trace to three handles.}
Fix a three-handle set $S$, let $W$ be its symplectic
span, and let $p:\mathbb L(V)\to\mathbb L(W)$ kill
the other generators. Put $D_S=pD|_W$.
It is special, since $p(D\omega)=D_S(\omega_W)$;
it remains tangent to the six basis lines of $W$.
For $\delta\in\Lambda^3W$, extend the corresponding
degree-one derivation by zero outside $W$.
On words in $W$ the projected bracket is precisely
$[D_S,\delta]$.

We check the extra contraction terms from every
outside handle $(x_i,y_i)$ explicitly. Because
$\delta(x_i)=\delta(y_i)=0$, one has
\[
 [D,\delta](x_i)=-[\delta A_i,x_i],\qquad
 [D,\delta](y_i)=-[\delta B_i,y_i].
\]
In the dual tensor
$x_i\otimes[D,\delta](y_i)
 -y_i\otimes[D,\delta](x_i)$, contraction of
the first two slots followed by projection to
$W$ gives
\[
                           p\delta A_i+p\delta B_i.
\]
Indeed the surviving terms are $y_i\delta B_i$
and $x_i\delta A_i$, with the signs of the two
symplectic pairings cancelling the displayed
dual-tensor signs. Every other term retains an
outside letter and projects to zero. Each
displayed term is a Lie polynomial of degree~7
in $W$. It is zero in cyclic coinvariants,
because every Lie bracket is an associative
commutator and its two products are cyclic
rotations. Therefore the exact identity is
\[
 p_{\mathrm{cyc}}\operatorname{ES}_{7,V}([D,\delta])
       =\operatorname{ES}_{7,W}([D_S,\delta]).              \tag{AG.8}
\]
This calculation needs only tangency; equality
$A_i=B_i$ was not assumed in proving it.

For $D\in U_{g,6}\cap T_g$, (AG.3) and (AG.8)
put $D_S$ in the exact genus-three trace kernel
computed in F.3.24. Every basis element in
(AG.5)--(AG.6) whose support is not contained
in $S$ restricts to zero: for an inside active
generator its conjugator contains a letter
from every other handle in its support, while
an outside active generator is simply omitted.
Hence, writing the complete expansion of $D$
in the basis above, F.3.24 gives
\[
    \text{every coefficient of }P_{ij}=0,
    \qquad
    c_{i;jk}+c_{j;ik}+c_{k;ij}=0
    \text{ on each triple}.                               \tag{AG.9}
\]
Each pair belongs to a three-handle set because
$g\geq3$. Thus the remaining triple contributions
are differences of the three $T$ elements.

\paragraph{4. The actual quotient kills the remaining components.}
Every handle permutation is an actual element of
$\operatorname{Sp}(V)$. It permutes the three
$T$ elements on a fixed triple transitively.
Therefore every symplectic-invariant linear
functional has the same value on those three
elements and vanishes on their differences.

On a fixed quadruple the program verifies the
stronger exact identities
\[
       \sum_{\pi\in\mathfrak S_4}\pi Q^{(1)}=0,
       \qquad
       \sum_{\pi\in\mathfrak S_4}\pi Q^{(2)}=0.            \tag{AG.10}
\]
The action permutes both the handle indices of
the conjugators and every formal $z_i$ inside
them. These are integer associative-word
identities, checked before the substitution
$z_i=[x_i,y_i]$. There is also a direct independent
proof of these averages. Fix an output handle. For
any input conjugator $A_i$, summing the permutations
that carry handle $i$ to that output sums all six
permutations of its multilinear degree-three Lie
word on the remaining three $z$ variables. Write
that word as a sum of nested brackets. Each nested
bracket has zero such sum, by pairing permutations
that exchange the two arguments of its inner
bracket. Sum this identity over the four possible
input handles. It gives (AG.10) without a rank or
representation-theoretic inference. If a functional
is invariant, its value on the left side is $24$
times its value on the original vector. The coefficient
field has characteristic zero, including when
$\ell=2$ or~3; hence the functional vanishes
on both quadruple basis elements.

The actual submodule $U_{g,6}$ is
$\operatorname{Sp}(V)$-stable, and $q_{g,6}$
is symplectic-invariant. Complete reducibility
of finite-dimensional symplectic representations
extends it to an invariant functional on
$\mathfrak h_{g,6}$. This is only a linear
extension; it assigns no arithmetic meaning
to a vector outside $U_{g,6}$. Equations
(AG.9)--(AG.10) show that the extension
vanishes on $D$, proving
\[
                         q_{g,6}(U_{g,6}\cap T_g)=0.
                                                               \tag{AG.11}
\]

Finally, both $U_{g,6}$ and $T_{g,\mathrm{all}}$
are stable under the full diagonal similitude
torus. For $U_{g,6}$, symplectic stability and
the scalar action $s^6$ on degree-six derivations
give similitude stability through the central
isogeny. The explicit Lagrange weight projector
of F.3.25 onto the $t^3$ eigenspace preserves
their intersection and satisfies $q_{g,6}e=q_{g,6}$.
Applying (AG.11) proves (AG.1). This projection
is applied after an actual arithmetic class
has satisfied all line tangencies; no lift
of this subsequent linear operator is needed.

\paragraph{Remaining proper-summand map.}
The proof has established a linear restriction
map $D\mapsto D_S$ and the precise trace identity
(AG.8). It has not identified $q_{g,6}(D)$ with
the degree-six arithmetic quotient of a proper
symplectic summand. Such an identification, or
sufficient additional mixed Johnson equations,
is still needed when the actual conjugate CM
normalizer lines fill only part of the homology
of the detecting curve in F.3.23.

\end{proof}


\begin{result}{F.3.27}{All-genus normalizer recognition for a whole absolutely simple Jacobian over Q}{Proved; every prime and finite arithmetic base change}
Let $C/\Q$ be a smooth proper geometrically connected curve
of genus $g\geq3$, with absolutely simple Jacobian and
$y\in C(\Q)$. Put $U=C\setminus\{y\}$.
For every prime $\ell$ and finite extension $K/\Q$,
a nontrivial closed procyclic subgroup of
$P=\pi_1(U_{\overline K})^{(\ell)}$ whose arithmetic
normalizer has open image in $G_K$ is contained in
an actual cusp inertia subgroup.
\end{result}
\begin{proof}
We retain the original curve and its full homology while
constructing every conjugate normalizer and universal map.

\paragraph{The first Magnus term and an actual finite detector.}
Let $A$ be the subgroup and $N$ its arithmetic normalizer.
The group $P$ is free pro-$\ell$ of rank $2g$, hence
torsion-free, so $A\simeq\Z_\ell$. For a generator $u$
its rational Magnus logarithm has a first nonzero term
$u_d\in\mathbb L_d(V)$, where $V=V_\ell J_C$.
Its normalizer character $\theta$ is trivial on
$N\cap P$, because geometric conjugation fixes that
initial term. It therefore gives an arithmetic
character on the open image of $N$.

The exact finite-cover construction in F.3.21 works
in this rank: take a finite $\ell$-quotient where
$u$ is nontrivial, and the preimage of its cyclic
image. The resulting actual finite etale cover $Y$
has nonzero rational homology class $[u]$ after
finite arithmetic descent. The normalizer intersection
with its arithmetic covering group still has open
image. Its geometric kernel acts inner and trivially
on the covering homology. Thus $[u]$ and $u_d$
have the very same arithmetic character $\theta$.
The original curve has one puncture, so $V$ is
pure of weight one and $\mathbb L_d(V)$ of weight
$d$. A nonzero proper-homology image of $[u]$
would force $d=1$ by a common good Frobenius.
It remains to exclude this degree-one case on
the original $V$.

\paragraph{The full CM field and all local primes.}
For $d=1$, the nonzero line $\Q_\ell u_1\subset V$
is stable on an open arithmetic subgroup. After
defining every geometric endomorphism, Faltings'
semisimplicity and endomorphism theorem identify
its projector with an idempotent of
$\operatorname{End}^0_{\Qb}(J_C)\otimes\Q_\ell$.
The rank formula in F.3.18 applies to this very
Jacobian. Absolute simplicity forces its full
geometric endomorphism algebra to be a CM field
$E$ of degree $2g$; the resulting simple quotient
is isogenous to $J_C$ itself.

The general CM/Q part of F.3.25 makes $E/\Q$
Galois and the actual endomorphism action
$G_{\Q}\to\operatorname{Aut}(E)$ surjective.
The module $V$ is rank one over
$E\otimes\Q_\ell=\prod_{\lambda\mid\ell}E_\lambda$.
Every idempotent rank is a sum of the positive
local degrees. Those degrees are all equal by
Galois transitivity, even at ramified primes.
A rank-one projector therefore requires complete
splitting of $\ell$ in $E$. This already excludes
$d=1$ at every nonsplit prime, with no prime
excluded from the argument.

\paragraph{All actual normalizer lines at a split prime.}
At a split prime the initial line is one of
the $2g$ embedding lines. Choose a nonzero
rational tangent vector at $y$ and its actual
arithmetic section. The exact formula
\[
 \rho_\ell(\tau)V_\varphi
          =V_{\varphi\circ a_\tau^{-1}}
\]
from F.3.25 supplies $2g$ elements $\tau_i\in G_{\Q}$
whose section lifts carry that line onto every
embedding line. Conjugate $A$ and $N$ by those
lifts. All resulting groups $A_i,N_i$ lie in
the arithmetic group of the same $U/\Q$.
Choose a finite Galois extension $L/\Q$ containing
$K$ and a field defining all endomorphisms, with
$G_L$ contained in every conjugate open image
$p(N_i)$. Every restricted $N_i$ now surjects
onto one common $G_L$.

Rosati duality pairs exactly the conjugate
embedding lines and is nondegenerate on each
pair, so generators of these $2g$ actual orbit
lines give a symplectic basis after scaling.
The arithmetic quotient of Nakamura's original
\S4, pp.~115--117, applies to every hyperbolic
$(g,n)$ with $n\geq1$. In particular Lemma~4.5,
Corollary~4.6, the degree-six rank table, and
the actual maps (4.7.1)--(4.7.2) give, for this
$(g,1)$ and this finite base $L$, an actual
$\sigma\in G_L(6)$ whose canonical special
derivation $D\in U_{g,6}$ satisfies
\[
                            q_{g,6}(D)\ne0.
\]
The quotient has dimension one and character
$\chi_\ell^3$ in this genus range
\cite[pp.~115--117]{AuditNakamura1996}.

Lift this same $\sigma$ to each $N_i$.
Its action on $V$ is the identity, so its
normalizer exponent on each $A_i$ is one.
The actual canonical representative and the
normalizing representative differ by an inner
automorphism of the punctured group. The rational
completed free-Lie calculation (QCM.4)--(QCM.5)
in F.3.25, now repeated on these $2g$ logarithms,
removes the degree-one inner term by a rational
power of the full subgroup logarithm and then
removes degrees two through five by the
centralizer of its nonzero linear term. It gives
\[
                         D(v_i)\in[\mathbb L_6(V),v_i]
                         \qquad(1\leq i\leq2g).
\]
Thus the same actual $D$ lies in
$U_{g,6}\cap T_{g,\mathrm{all}}$. F.3.26 proves
that its quotient is zero, a contradiction.
The degree-one case is excluded also at every
split prime.

\paragraph{The boundary character and the full normalizer.}
Return to the actual detector $Y$. Its nonzero
line $[u]$ has zero image in proper homology
because that would force the excluded case
$d=1$. It belongs to the actual boundary kernel
\[
 \left(\bigoplus_{t\in(\overline Y\setminus Y)(\Qb)}
          \Q_\ell(1)e_t\right)
 \big/\Q_\ell(1)\Bigl(\sum_t e_t\Bigr).
\]
Over a finite extension defining those points,
this entire kernel has character $\chi_\ell$.
Hence $\theta=\chi_\ell$ on the intersected
normalizer. Any element of the full normalizer
over that extension has the same arithmetic
image as an element of the intersected one;
their quotient lies in $N\cap P$, where
$\theta$ is trivial. Thus the full normalizer
also acts cyclotomically and has open arithmetic
image. Nakamura's Theorem~2.1.1, p.~94, for
the full class of finite $\ell$-groups gives
actual cusp inertia containing $A$
\cite{AuditNakamura1994}. The finite arithmetic
extensions change neither the geometric group
nor this unique-puncture inertia conjugacy class.
\end{proof}

\paragraph{Exact relation to the general detecting-cover problem.}
The assertion concerns the whole absolutely simple
Jacobian of a specified once-punctured curve over
$\Q$. The rational normal-core construction of
F.3.23 gives an actual once-punctured $\Q$-curve,
but its Jacobian need not be absolutely simple.
For a proper CM quotient the conjugate detected
lines still lie in its actual images inside the
larger curve homology. The mixed normalizer-trace
equations or an actual punctured target map
needed there are retained as unresolved steps;
the full LM6 A.22 subgroup assertion is unchanged.


\begin{result}{F.3.28}{Two paired tangencies on a proper summand in genus three}{Proved; exact geometric contraction image}
Let $V$ be symplectic of dimension six with basis
$a,b,c,d,e,f$, paired as $(a,b),(c,d),(e,f)$.
Let $U_{3,6}\subset\mathfrak h_{3,6}$ and
$q_6:U_{3,6}\to Q_6$ be the actual universal
arithmetic image and quotient of F.3.24.
Write $P_{\mathrm{all}}$ for the subspace of
special degree-six derivations tangent to the
first four basis lines:
\[
 D(x)\in[\mathbb L_6(V),x]\qquad(x=a,b,c,d).
\]
The values on $e,f$ are unrestricted except by
the specialness equation. Then
\[
                         q_6(U_{3,6}\cap P_{\mathrm{all}})=0.
                                                               \tag{PT.1}
\]
This gives a partial-summand statement on the original
six-dimensional space. No arithmetic curve quotient
or compatibility of $q_6$ with a restriction map is used.

\end{result}
\begin{proof}
We keep the complementary handle and all mixed Johnson equations, and compare invariant contractions with two actual geometric brackets.
\paragraph{1. The actual geometric quotient and all invariant functionals.}
The original exact sequence (4.7.1) of Nakamura,
printed p.~117~\cite{AuditNakamura1996}, identifies the universal quotient
as the quotient of the universal image by its
geometric image. The last paragraph of that page
also explicitly places the topological Johnson
image in the geometric image. Under the strict
punctured-action map of F.3.24, this gives
\[
                    J_6\subset U_{3,6},\qquad q_6(J_6)=0.
                                                               \tag{PT.2}
\]
Here $J_6$ denotes the actual degree-six geometric
Johnson image in the canonical special derivations.
In particular every sixfold Lie bracket of
degree-one Johnson derivations belongs to this
specified kernel. These are actual geometric
image vectors, not Lie polynomials merely
declared to be arithmetic.

Choose a nonzero basis of the one-dimensional $Q_6$,
whose dimension was verified in F.3.27, to regard $q_6$
as a scalar-valued map. By complete reducibility, the
symplectic-invariant map $q_6$ extends to a symplectic-invariant linear
functional $F$ on $\mathfrak h_{3,6}$, retaining
its zero value on $J_6$. The natural tensor map
\[
 \mathfrak h_{3,6}\hookrightarrow V\otimes\mathbb L_7(V)
                      \hookrightarrow V^{\otimes8}
\]
sends $D$ to
\[
 a\otimes D(b)-b\otimes D(a)
 +c\otimes D(d)-d\otimes D(c)
 +e\otimes D(f)-f\otimes D(e).
\]
Extend $F$ once more, equivariantly, to $V^{\otimes8}$.
Every invariant functional there is a linear
combination of the 105 complete symplectic
contractions indexed by perfect matchings of
the eight tensor slots. This is precisely the
all-genus surjectivity of Morita--Sakasai--Suzuki,
original v2, Proposition~2.1, p.~8, with the
dual contraction formula on p.~7~\cite{AuditMoritaSakasaiSuzuki2015v2}. Its bijectivity
requires $g\geq4$ here, but its stated surjectivity
holds for every $g$, including $g=3$.
The program uses the unsigned products of the
four pairings; these differ from that paper's
convention by an invertible sign on each of
the 105 coordinates, which changes neither
their span nor the image inclusion below.

Let $K:\mathfrak h_{3,6}\to\Q_\ell^{105}$ be
the vector of those contractions. To show
$F(D)=0$ it therefore suffices to show
\[
                    K(D)\in K(J_6).                       \tag{PT.3}
\]
This implication uses a functional on the original
derivation space. It makes no identification with
an arithmetic quotient belonging to a smaller curve.

\paragraph{2. The complete balanced partial-tangent space.}
Give the full diagonal similitude torus the weights
$a:\alpha_1,b:t/\alpha_1,c:\alpha_2,d:t/\alpha_2,
e:\alpha_3,f:t/\alpha_3$, and let $P$ be the
$t^3$ part of $P_{\mathrm{all}}$.
The specialness equation has balanced output
multidegrees $(n_1,n_1,n_2,n_2,n_3,n_3)$ with
$n_1+n_2+n_3=4$.

For a constrained generator $x\in\{a,b,c,d\}$,
the column space consists exactly of
$[A_x,x]$, with $A_x$ in the corresponding
degree-six Lyndon basis. For an unconstrained
generator $x\in\{e,f\}$, the column space
is the complete degree-seven Lyndon basis
of the required multidegree. Applying
$D\omega$ gives the following full block table:
\[
\begin{array}{c|r|r|r|r}
n&\text{rows}&\text{columns}&\text{rank}&\text{nullity}\\\hline
(4,0,0),(0,4,0)&8&6&6&0\\
(0,0,4)&8&10&8&2\\
(3,1,0),(1,3,0)&140&66&66&0\\
(3,0,1),(0,3,1)&140&100&99&1\\
(1,0,3),(0,1,3)&140&126&123&3\\
(2,2,0)&312&120&119&1\\
(2,0,2),(0,2,2)&312&240&236&4\\
(2,1,1),(1,2,1)&1260&660&657&3\\
(1,1,2)&1260&840&834&6.
\end{array}                                               \tag{PT.4}
\]
For rows listing two multidegrees, the displayed
sizes and nullity apply to each block. They sum
to $\dim P=31$.

These are exact rational ranks. The complete
program \texttt{check\_genus\_three\_two\_pair\_}\allowbreak\texttt{partial.py}
performs elimination modulo~101 for the lower
bounds and lifts each kernel vector to an
integer vector, checking every original
specialness equation over the integers.
The nonzero multiplier of each lift modulo~101
preserves the independence of the modular
free-coordinate kernel basis. The upper and
lower bounds therefore agree over $\Q$.
This argument does not identify a modular
kernel with a rational one without checking
the lifted equations.

For all twenty degree-one derivations
$\delta_{uvw}$ of F.3.24, form
$\operatorname{ES}_7([D,\delta_{uvw}])$.
The exact matrix on the 31 basis vectors has
3920 nonzero coordinate rows and rank~15.
The program supplies sixteen independent
integer kernel vectors and verifies all their
entries against the original integer matrix.
Thus the complete subspace
\[
 P^0=\{D\in P:\operatorname{ES}_7([D,\delta])=0
                      \text{ for every }\delta\in\Lambda^3V\}
\]
has dimension~16 over $\Q$ and over $\Q_\ell$.
The actual geometric-normalizer equation of
F.3.24 places $U_{3,6}\cap P$ inside $P^0$.

\paragraph{3. Exact contraction image in two specified geometric brackets.}
Write brackets right-nested, so that
$[D_1,D_2,\ldots,D_6]$ means
$[D_1,[D_2,[D_3,[D_4,[D_5,D_6]]]]]$.
Define the actual geometric degree-six elements
\[
\begin{split}
 Z_1&=[\delta_{cdf},\delta_{cef},\delta_{aef},
          \delta_{def},\delta_{bef},\delta_{cde}],\\
 Z_2&=[\delta_{bce},\delta_{adf},\delta_{adf},
          \delta_{abd},\delta_{bce},\delta_{abc}].
\end{split}                                               \tag{PT.5}
\]
They lie in $J_6$ by its graded Lie bracket.
The contraction vectors $K(Z_1),K(Z_2)$ are
independent: on the matchings
$(12)(34)(56)(78)$ and $(12)(36)(47)(58)$
their two-by-two determinant is $-10752$.

For each of the sixteen exact basis vectors
$D_j$ of $P^0$, the program computes its
105 integer contractions and solves for the
two coefficients using that nonzero minor.
It then verifies the resulting identity in
every one of the 105 coordinates over the
integers after multiplying by $-10752$.
The complete list of integer coefficient
triples is retained in the output. Hence
\[
       K(P^0)\subset\operatorname{span}_{\Q}
                         \{K(Z_1),K(Z_2)\}\subset K(J_6).
                                                               \tag{PT.6}
\]
The nonzero rational determinant remains invertible
in every $\Q_\ell$; no integral division is needed.
Both the source kernel and the image inclusion
are exact; a finite-field rank coincidence
alone is not used for (PT.6).

For $D\in U_{3,6}\cap P$, choose the invariant
extension $F$ above. Express it as a linear
combination of the 105 contractions. Equation
(PT.6) then gives
$F(D)=\alpha F(Z_1)+\beta F(Z_2)=0$
by (PT.2). Since $F(D)=q_6(D)$ on the actual
universal submodule, the balanced assertion follows.

Finally $P_{\mathrm{all}}$ is stable under the
full diagonal torus despite imposing no
constraint on $e,f$. The explicit weight
projector of F.3.25 preserves its intersection
with $U_{3,6}$, projects to $P$, and preserves
$q_6$. This proves (PT.1) for all degrees of
the unconstrained torus weights as well.

\paragraph{Exact scope and the one-pair test.}
The theorem imposes four actual line tangencies
spanning a nondegenerate four-dimensional
symplectic summand of this genus-three homology.
It retains both complementary generators and
all mixed Johnson equations. It makes no
assertion that arbitrary CM quotient lines
necessarily supply that nondegenerate summand.
A separate discovery calculation with only
one paired handle has an additional contraction
direction modulo~101; it has not supplied an
arithmetic class or a proof of a converse.
That calculation is preserved as work provenance
and is not used in the theorem.

\end{proof}

\begin{result}{F.3.29}{A single normalizer line in genus three}{Proved; every prime and unrestricted Jacobian}
Let $V$ be a six-dimensional symplectic vector space over $\Q_\ell$.
Let $U_{3,6}\subset\mathfrak h_{3,6}$ and
$q_6:U_{3,6}\to Q_6$ be the actual universal arithmetic image and
the one-dimensional quotient in F.3.24 and Nakamura's original
(4.7.1)--(4.7.2), printed p.~117.
For every nonzero $a\in V$, put
\[
 P_a^{\rm all}=\{D\in\mathfrak h_{3,6}:
                         D(a)\in[\mathbb L_6(V),a]\}.
\]
Then
\[
                    q_6(U_{3,6}\cap P_a^{\rm all})=0.       \tag{SL.1}
\]
Only this one line is constrained. In particular, no conjugate line,
CM field, multiplicity space, or nondegenerate span of Galois conjugates
is an assumption of (SL.1).

For every smooth proper genus-three curve $C/\Q$ with a rational
point $y$, every prime $\ell$, and every finite extension $K/\Q$,
a nontrivial closed procyclic subgroup of
$\pi_1((C\setminus\{y\})_{\overline K})^{(\ell)}$
whose arithmetic normalizer has open image is contained in actual
cusp inertia. The Jacobian is unrestricted.
\end{result}
\begin{proof}
We use the actual mixed geometric-normalizer equations on the original
six-dimensional space and an exact contraction certificate.

\paragraph{The actual image maps.}
Choose a symplectic basis $a,b,c,d,e,f$.
Its pairs are $(a,b)$, $(c,d)$ and $(e,f)$.
Write $J_m$ for the actual geometric Johnson image in special
degree-$m$ derivations. The strict geometric-action construction of
F.3.24 gives
\[
 J_6\subset U_{3,6},\qquad q_6(J_6)=0,\qquad
                  [D,J_1]\subset J_7\quad(D\in U_{3,6}).   \tag{SL.2}
\]
These maps use Hain's strict weight filtration and the actual arithmetic
normalization of the geometric monodromy image, as checked in F.3.24.
The first inclusion and quotient are also explicit in Nakamura's
(4.7.1) and the last paragraph of printed p.~117
\cite{AuditNakamura1996}. Thus the brackets in (SL.2) are vectors in
the geometric image. The argument does not identify a formal
closed-surface ideal with that image.

For $u,v,w\in V$, set
\[
 \delta_{uvw}(x)=\langle x,u\rangle[v,w]
              +\langle x,v\rangle[w,u]
              +\langle x,w\rangle[u,v].
\]
These derivations give $J_1=\Lambda^3V$. In particular, the twenty
triples of distinct basis letters span $J_1$.

\paragraph{One additional primary-source detector.}
The Enomoto--Satoh map vanishes on $J_7$ by the source maps already
verified in F.3.24. The newly opened primary preprint
\cite[Propositions~4.1--4.2, p.~11]{AuditMoritaSakasaiSuzuki2026v1}
gives an additional symplectic-equivariant map
$R:V^{\otimes9}\to V$ whose restriction to $J_7$ is zero.
Explicitly,
\[
\begin{split}
 R(u_1\otimes\cdots\otimes u_9)
  =\bigl(&-2\langle u_1,u_2\rangle\langle u_3,u_5\rangle
                \langle u_4,u_7\rangle\langle u_6,u_8\rangle\\
         &+2\langle u_1,u_2\rangle\langle u_3,u_6\rangle
                \langle u_4,u_7\rangle\langle u_5,u_8\rangle\\
         &+2\langle u_1,u_2\rangle\langle u_3,u_6\rangle
                \langle u_4,u_8\rangle\langle u_5,u_7\rangle\\
         &+2\langle u_1,u_2\rangle\langle u_3,u_7\rangle
                \langle u_4,u_6\rangle\langle u_5,u_8\rangle\\
         &+\langle u_1,u_3\rangle\langle u_2,u_6\rangle
                \langle u_4,u_7\rangle\langle u_5,u_8\rangle\bigr)u_9.
\end{split}                                                   \tag{SL.3}
\]
The tensor map for a degree-seven derivation is
$D\mapsto\sum_i(x_i\otimes D(y_i)-y_i\otimes D(x_i))$.
For the genus-three use, choose an embedding of the boundary surface
in a surface of genus $G$ in the source's stable range. Extending a
mapping class by the identity gives the actual geometric map on
Johnson images, and its derivation extends by zero on the added
handles. For the symplectic inclusion $i:V\hookrightarrow V_G$,
every contraction in (SL.3) gives
$R_G(i^{\otimes9}T)=iR_3(T)$. The source's vanishing on the
extended actual Johnson image and the injectivity of $i$ therefore
give $R_3(J_7)=0$. No stable multiplicity count is imposed on
the original genus-three space.
Consequently every actual $D\in U_{3,6}$ satisfies
\[
 \operatorname{ES}_7([D,\delta])=0\quad(\delta\in J_1),
 \qquad [c]R([D,\delta_{abc}])=0,                         \tag{SL.4}
\]
where $[c]$ denotes the coefficient of $c$ in $V$.
Only this one scalar component of $R$ is used. In particular, the
final proof does not impose a raw exterior-power equation in place
of the quotient target of the separate map $S$ in that preprint.

The source gives a sketch of proof for Proposition~4.1. Our certificate
copies each of the five displayed contractions, and independently
checks its vanishing on the forty actual degree-seven brackets
$[Z_i,\delta_{uvw}]$ used below. The universal vanishing on $J_7$
is the stated primary-source result; the forty checks are a separate
test of the implemented formula and conventions.

\paragraph{The complete balanced one-line space.}
Use the diagonal similitude torus with weights
$a:\alpha_1,b:t/\alpha_1,c:\alpha_2,d:t/\alpha_2,
e:\alpha_3,f:t/\alpha_3$. Let $P_a$ be the $t^3$ part of
$P_a^{\rm all}$. Its specialness output has multidegree
$(n_1,n_1,n_2,n_2,n_3,n_3)$, where $n_1+n_2+n_3=4$.
For $D(a)$ use the entire image $[\mathbb L_6,a]$ in the
required multidegree. For every other generator, including $b$,
use the entire degree-seven Lyndon basis in its required multidegree.
The specialness equations have the following exact ranks:
\[
\begin{array}{c|r|r|r|r}
n&\text{rows}&\text{columns}&\text{rank}&\text{nullity}\\\hline
(0,0,4),(0,4,0)&8&10&8&2\\
(0,1,3),(0,3,1)&140&160&140&20\\
(0,2,2)&312&360&312&48\\
(1,0,3),(1,3,0)&140&143&140&3\\
(1,1,2),(1,2,1)&1260&1290&1260&30\\
(2,0,2),(2,2,0)&312&300&296&4\\
(2,1,1)&1260&1200&1188&12\\
(3,0,1),(3,1,0)&140&130&129&1\\
(4,0,0)&8&8&8&0.
\end{array}                                                   \tag{SL.5}
\]
Each of the two listed multidegrees has the displayed sizes.
The nullities sum to $180$.

The complete certificate is
\texttt{check\_genus\_three\_}\allowbreak\texttt{single\_line.py}.
For every block it computes a rank lower bound modulo~101 and
checks independent integer kernel lifts against the original
integer specialness matrix. The lower and upper bounds agree,
so these are rational ranks. Conjugation by $a$ is injective on
each degree-six multidegree used here: its free-Lie centralizer
consists of multiples of $a$. Indeed, the associative centralizer
of the free generator $a$ is $\Q[a]$: the coefficient equation
$aF=Fa$ along each finite chain of words with a non-$a$ core
forces those coefficients to vanish. The only primitive elements
of $\Q[a]$ in characteristic zero are the multiples of $a$,
as is seen from the binomial formula for its coproduct.
Hence $\Q[a]\cap\mathbb L_6=0$. Thus the selected column coordinates
describe the full derivation space without an unaccounted kernel.

On these 180 basis vectors the equations (SL.4) form an integer
matrix with $7977$ nonzero coordinate rows and rank $48$.
To certify the rank over $\Q$, the program selects 48 independent
rows modulo~101, performs exact rational row reduction on those
rows, clears denominators in all 132 free-coordinate kernel
vectors, and checks them against every original integer row.
Their independent free coordinates and the modular lower bound
prove that the full rational kernel has dimension $132$.
All these equations and ranks remain valid after extension to
every $\Q_\ell$.

\paragraph{The contraction image lies in two actual geometric brackets.}
Use right-nested brackets and define
\[
\begin{split}
 Z_1&=[\delta_{cdf},\delta_{cef},\delta_{aef},
                  \delta_{def},\delta_{bef},\delta_{cde}],\\
 Z_2&=[\delta_{bce},\delta_{adf},\delta_{adf},
                  \delta_{abd},\delta_{bce},\delta_{abc}].
\end{split}                                                   \tag{SL.6}
\]
Both belong to $J_6$ by the actual geometric Lie bracket.
Let $K$ be the vector of all 105 complete contractions of
the tensor associated with a special degree-six derivation.
The two rows indexed by
$(12)(34)(56)(78)$ and $(12)(36)(47)(58)$ give
\[
 \begin{pmatrix}-104&-288\\32&192\end{pmatrix},
             \qquad\det=-10752.                           \tag{SL.7}
\]
For every one of the 132 integer kernel vectors $D_j$, solve
for the two coefficients using (SL.7). The program checks
\emph{all} 105 coordinates of
\[
                 K(D_j)\in\operatorname{span}\{K(Z_1),K(Z_2)\}
                                                               \tag{SL.8}
\]
over the integers after multiplication by $-10752$.
The largest coefficient in the verified integer kernel is
$247830$, and all fixed-width multiplication bounds are checked.
The determinant is nonzero over every $\Q_\ell$; this step
makes no integral divisibility assertion.

Choose a basis of the one-dimensional $Q_6$ and regard $q_6$
as scalar valued. Symplectic complete reducibility extends it
equivariantly first to $\mathfrak h_{3,6}$ and then through
$\mathfrak h_{3,6}\hookrightarrow V^{\otimes8}$ to the full
tensor power. Every invariant functional on that tensor power
is a linear combination of the complete contractions, by
\cite[Proposition~2.1, p.~8]{AuditMoritaSakasaiSuzuki2015v2}.
Since it vanishes on $Z_1,Z_2$ by (SL.2), (SL.8) proves
$q_6(U_{3,6}\cap P_a)=0$.

\paragraph{The actual post-constraint projection.}
The diagonal torus preserves both $U_{3,6}$ and $P_a^{\rm all}$.
The quotient $Q_6$ has weight $t^3$. Take the algebraic weight
projector $e$ onto that weight in the finite-dimensional
derivation representation; it is a polynomial in a generic
diagonal torus element, with the denominators cleared in the
usual Lagrange formula. It satisfies
\[
 e(U_{3,6}\cap P_a^{\rm all})\subset U_{3,6}\cap P_a,
                    \qquad q_6(eD)=q_6(D).
\]
Applying the proved balanced assertion gives (SL.1).
The projection is applied to an already constrained vector in
the actual universal image. It is not asserted to be a new
individual Galois element.
\end{proof}

\paragraph{Arithmetic consequence on the original genus-three curve.}
Let $C/\Q$ be a smooth proper geometrically connected curve of
genus three, let $y\in C(\Q)$, and put $X=C\setminus\{y\}$.
For every prime $\ell$ and finite extension $K/\Q$, every nontrivial
closed procyclic subgroup of $\pi_1(X_{\overline K})^{(\ell)}$
whose arithmetic normalizer has open image in $G_K$ is contained
in an actual cusp inertia subgroup. No restriction on the
Jacobian is needed.

\begin{proof}
We apply (SL.1) to the original homology, so no CM decomposition
or permutation of character lines enters the proof.
Let $u$ generate the subgroup, $N$ be its arithmetic normalizer,
and $u_d$ its first nonzero rational Magnus term.
The geometric kernel of $N$ acts trivially on $u_d$, so the
normalizer exponent gives a character $\theta$ on the open
arithmetic image. The exact finite-cover detector of F.3.21
gives an actual covering $Y$ in whose rational homology $u$
has nonzero class, after finite arithmetic descent.
The intersection of the normalizer with that covering group
still has open arithmetic image; its geometric kernel acts
inner and trivially on the covering homology. Consequently
$u_d$ and this class have the same character $\theta$.

The original curve has one puncture, so its degree-one
homology is pure of weight one and its degree-$d$ free-Lie
part has weight $d$. If the detected class has nonzero image
in proper homology, a common good Frobenius gives $d=1$.
We exclude this case directly. Put $a=u_1\ne0$ and choose
a symplectic basis extending $a$. Over a finite extension
$L$ with $G_L\subset p(N)$, the actual finite-base-invariant
arithmetic quotient of Nakamura, pp.~115--117, supplies
\[
 \sigma\in G_L(6),\qquad D_\sigma\in U_{3,6},
                         \qquad q_6(D_\sigma)\ne0.
\]
Lift this same $\sigma$ to $N$. It acts identically on
homology, so its exponent on $u$ is one. If $X_u=\log u$
and $Z$ is the logarithm of the inner correction between
the normalizing lift and the canonical special representative,
the resulting equality is
\[
                   \phi_\sigma(X_u)=\exp(\operatorname{ad}Z)X_u.
\]
The left side first differs from $X_u$ in degree seven.
The degree-one term of $Z$ centralizes $a$ and is a rational
multiple of $a$. Multiplying the correction by a rational
power of the full $u$ removes that term without changing
the displayed conjugation. Successive degrees two through
five of $Z$ then centralize $a$; the associative-centralizer and
primitive-element argument above shows that the free Lie algebra
has no such homogeneous centralizer part. Comparing
degree seven gives
\[
                     D_\sigma(a)=[Z_6,a].
\]
This is exactly the nonprimitive-logarithm correction already
proved in F.3.22 and F.3.25, applied to one line. Now (SL.1)
contradicts the chosen nonzero quotient.

The detector's nonzero class therefore belongs to its boundary
kernel. After defining all boundary points, that kernel has
character $\chi_\ell$. Thus $\theta=\chi_\ell$ on the
intersected normalizer. Every element of the full normalizer
over that finite extension differs from such a lift by an
element of its geometric kernel, where $\theta$ is trivial.
The full normalizer is consequently cyclotomic with open
arithmetic image. Nakamura's original Theorem~2.1.1,
printed p.~94, applies to the class of finite $\ell$-groups
and gives actual cusp inertia containing the subgroup
\cite{AuditNakamura1994}. The argument concerns every prime.
\end{proof}

\paragraph{Separate all-genus balanced tangency lemma.}
For every genus $g\geq3$, a special degree-six derivation in
the $t^3$ weight that is tangent to a nonzero basis vector
$a$ is also tangent to its chosen symplectic partner $b$.
This assertion concerns the balanced derivation space; an
all-genus version of (SL.1) still needs its own proof.

\begin{proof}
Write $T=\sum_i(x_i\otimes D(y_i)-y_i\otimes D(x_i))$.
The specialness equation says $(1-\rho)T=0$, where $\rho$
moves the first tensor letter to the end. The tensor is thus
cyclic. In every balanced block its handle multiplicities
$n_i$ sum to four. If the constrained handle is absent,
both $D(a)$ and $D(b)$ vanish: their required multidegrees
are obtained by subtracting the partner letter from the
tensor multidegree and would otherwise have a negative entry.

If this handle occurs once, write $D(a)=[A,a]$; the
multidegree of $A$ contains neither $a$ nor $b$. The
$b$-leading part of $T$ is
\[
                      -bAa+baA.
\]
Each tensor word has a unique occurrence of $b$, so this
part determines its full cyclic orbit and hence determines
the whole block of $T$. Rotating to the unique $a$-leading
positions gives $-abA+aAb=a[A,b]$. Therefore
$D(b)=[A,b]$, as asserted, even for a block on four handles.

If the constrained handle has multiplicity at least two,
there are at most two other handles in its support. The
complete single-line block bases in (SL.5) are all
certified to be paired by the same final program.
For each polynomial $D(b)$ it groups words as
$b^i w b^j$, with $w$ beginning and ending in a letter
different from $b$, and solves the coefficient chains of
$[B,b]=D(b)$. It then verifies that exact commutator and
the Dynkin identity $\operatorname{Dyn}_6(B)=6B$, which
places $B$ in $\mathbb L_6$. These checks cover all 180
basis vectors, so they prove the actual inclusion of the
single-line space in the paired-line space over $\Q$ and
every $\Q_\ell$; the reverse inclusion is immediate.
Rename the at most three support
handles to obtain the required equality in every genus.
This argument gives the actual inclusion and equality of
the two specified balanced tangent spaces.
\end{proof}

\paragraph{Retained scope and source status.}
The full LM6 A.22 subgroup assertion still permits a rational
once-punctured detecting curve of arbitrary genus. Its genus-three
case is now handled without a simplicity or pairing hypothesis.
The all-genus mixed-normalizer calculation is not established
by the support bound alone. The MSS source is the opened
18 August 2026 v1 preprint. Its two separately recorded arithmetic
slips in the proofs on pp.~22 and 25 do not occur in (SL.3),
and the final argument above imports neither of the affected
scalar expansions.


